% AgentVax technical paper, version 4 (owner-pressure safeguards)
%
% Clean version:        pdflatex agentvax_paper_v4.tex   (run twice)
% Highlighted version:  pdflatex -jobname=agentvax_paper_v4_highlighted "\def\dohighlight{1}% AgentVax technical paper, version 4 (owner-pressure safeguards)
%
% Clean version:        pdflatex agentvax_paper_v4.tex   (run twice)
% Highlighted version:  pdflatex -jobname=agentvax_paper_v4_highlighted "\def\dohighlight{1}% AgentVax technical paper, version 4 (owner-pressure safeguards)
%
% Clean version:        pdflatex agentvax_paper_v4.tex   (run twice)
% Highlighted version:  pdflatex -jobname=agentvax_paper_v4_highlighted "\def\dohighlight{1}% AgentVax technical paper, version 4 (owner-pressure safeguards)
%
% Clean version:        pdflatex agentvax_paper_v4.tex   (run twice)
% Highlighted version:  pdflatex -jobname=agentvax_paper_v4_highlighted "\def\dohighlight{1}\input{agentvax_paper_v4.tex}"
%                       (or uncomment the next line)
% \def\dohighlight{1}
\documentclass[10pt,a4paper,twocolumn]{article}

\usepackage[T1]{fontenc}
\usepackage[utf8]{inputenc}
\usepackage{mathptmx}
\usepackage[scaled=0.9]{helvet}
\usepackage[a4paper,top=19mm,bottom=22mm,left=16mm,right=16mm,columnsep=7mm]{geometry}
\usepackage{amsmath,amssymb}
\usepackage{microtype}
\usepackage{booktabs,tabularx,array}
\usepackage{stfloats}
\usepackage[font=small,labelfont=bf,labelsep=period]{caption}
\usepackage{titlesec}
\usepackage{enumitem}
\usepackage{algorithm}
\usepackage{algpseudocode}
\usepackage{xcolor}
\usepackage{tikz}
\usetikzlibrary{positioning,arrows.meta,fit,calc}
\usepackage[hidelinks]{hyperref}
\usepackage{url}
\usepackage{balance}

% ---- Yellow highlighter that breaks across lines (no extra package needed) ----
\definecolor{hlyellow}{HTML}{FFF14D}
\makeatletter
\newcommand{\uniq@box}[1]{\begingroup\setlength{\fboxsep}{0pt}%
  \colorbox{hlyellow}{\rule[-0.25em]{0pt}{1.05em}#1}\endgroup}
\newcommand{\uniq@gap}{{\color{hlyellow}\leaders\hrule height 0.8em depth 0.25em
  \hskip \fontdimen2\font plus \fontdimen3\font minus \fontdimen4\font}}
\def\uniq@scan#1 #2\uniq@end{%
  \if\relax\detokenize{#1}\relax\else\uniq@box{#1}\fi
  \if\relax\detokenize{#2}\relax\expandafter\@gobble\else\expandafter\@firstofone\fi
  {\uniq@gap\uniq@scan#2\uniq@end}}
\ifdefined\dohighlight
  \newcommand{\uniq}[1]{\leavevmode\uniq@scan#1 \uniq@end}
\else
  \newcommand{\uniq}[1]{#1}
\fi
\makeatother

% Headings
\titleformat{\section}{\normalfont\large\bfseries}{\thesection.}{0.5em}{}
\titleformat{\subsection}{\normalfont\normalsize\bfseries\itshape}{\thesubsection}{0.5em}{}
\titlespacing*{\section}{0pt}{9pt plus 2pt minus 1pt}{4pt}
\titlespacing*{\subsection}{0pt}{7pt plus 2pt minus 1pt}{3pt}
\setlist{nosep,leftmargin=*,topsep=2pt}

% Table helpers: ragged-right wrapping columns so text never overflows a cell
\newcolumntype{L}{>{\raggedright\arraybackslash}X}
\newcolumntype{P}[1]{>{\raggedright\arraybackslash}p{#1}}
\renewcommand{\arraystretch}{1.25}
\setlength{\tabcolsep}{4pt}
\newcommand{\tblfont}{\footnotesize}

\newcommand{\todo}[1]{\textcolor{red!70!black}{[#1]}}

\emergencystretch=2em
\sloppy
\raggedbottom

\title{\vspace{-6mm}\textbf{AgentVax: A Self-Immunising Wallet Agent\\
with Adaptive Financial Autonomy}}
\author{%
\todo{Author 1}, \todo{Author 2}, \todo{Author 3}\\[2pt]
\small Team \todo{team name}, \todo{university, country}\\
\small 2026 Shenzhen International FinTech Competition (FinTechathon), International Track, Topic F: Agent Wallet\\
\small Contact: \todo{email}}
\date{}

\begin{document}
\maketitle

\begin{abstract}
\noindent
AI agents that hold wallets can be talked into spending money they should not spend, and
several have already lost real funds. Existing agent wallets ask what the agent is
allowed to do and enforce it with caps and allowlists.
\uniq{AgentVax asks a second question: how much of its authorised power should the agent
exercise on its own right now?}
The owner states a mission in plain language; it becomes a signed mandate on chain that
the model cannot change. Before every payment the wallet contract checks
\uniq{three things: is it allowed, can the owner afford it, and is it safe right now,}
and
\uniq{a daily caution allowance limits how many borderline payments may pass without the
owner being asked.}
\uniq{When one wallet is attacked it publishes an immunity capsule, a structured alert
with no private data that describes the attack by its features, so other wallets
recognise reworded versions and become more cautious at once.}
\uniq{A detector trained jointly with federated learning then gives lasting protection
without any wallet sharing its data.}
\uniq{Learning and shared intelligence can only reduce autonomy. Only the owner can
extend authority, and only after a delay chosen in advance, so that an owner under
pressure cannot be rushed into a mistake.}
The paper specifies the architecture, the algorithms, the security design and a
reproducible evaluation protocol whose headline measure is the share of legitimate tasks
completed without the owner, with zero policy violations. It is a design and protocol
document; measured results are reported in the accompanying execution evidence.
\end{abstract}

\vspace{2pt}
\noindent\textbf{\textit{Keywords---}}agent wallet, adaptive autonomy, prompt injection,
payment fraud, threat sharing, federated learning.

\ifdefined\dohighlight
\smallskip\noindent{\footnotesize\textbf{Reading note.} \uniq{Yellow highlight} marks what
is specific to AgentVax. Unhighlighted features, such as spending caps, allowlists and
session keys, are common to agent wallets.\par}
\fi

% ---------------------------------------------------------------------------
\section{Introduction}

The theme of the International Track is ``AI as a Financial Participant'': an agent that
completes real tasks in a financial scenario~\cite{rules}. Topic~F asks for a wallet agent
that autonomously performs asset management, transaction execution and risk control on
chain. Entries are scored on task completion (40\%), security and compliance (30\%) and
innovation and interaction (30\%), with a 0--5 point bonus for on-chain
exploration~\cite{rules}.

A wallet agent differs from a chatbot in one respect that matters for security: its output
moves value. A language model reads text from many sources and cannot reliably tell data
from instructions~\cite{greshake,owasp}. Table~\ref{tab:incidents} lists four public
incidents in which an AI agent with a wallet sent funds it should not have sent. In the
most recent, an attacker posted an instruction in Morse code; one model translated it and
a wallet bot executed the result~\cite{oecdgrok,giskard}. In the laboratory, Patlan et
al.\ showed that planting false entries in a wallet agent's memory is more effective
still~\cite{patlan}. Regulators have named the same risk for agent-initiated
payments~\cite{imf}. An agent that pays invoices also inherits the oldest payment fraud
there is: a message that appears to come from a known supplier and asks for payment to a
new account. That fraud usually works by rushing a person, not by fooling a machine.

\begin{table*}[t]
\caption{Reported incidents involving AI agents that control wallets.}
\label{tab:incidents}
\tblfont
\begin{tabularx}{\textwidth}{@{}P{0.09\textwidth} P{0.13\textwidth} L P{0.17\textwidth} P{0.19\textwidth}@{}}
\toprule
\textbf{Date} & \textbf{Agent} & \textbf{What happened} & \textbf{Reported loss} & \textbf{Kind of failure} \\
\midrule
Nov.\ 2024 & Freysa~\cite{freysa} & A player persuaded the agent that its transfer-approval function handled incoming funds; it released the prize pool. This was a designed game. & About \$47,000 & Persuasion of the model \\
Mar.\ 2025 & AIXBT~\cite{aixbt} & An attacker entered the control dashboard and queued two prompts instructing the agent to transfer funds. & 55.5 ETH, about \$106,000 & Injected instructions after a dashboard breach \\
Feb.\ 2026 & Lobstar Wilde~\cite{lobstar} & The agent intended to send a small gift and transferred its entire token holdings. & About \$250,000 in tokens & Agent error in the amount \\
May 2026 & Grok and Bankr~\cite{oecdgrok,giskard} & A token transfer raised the wallet's permissions; a Morse-coded post was translated into a transfer order and executed. & About \$150,000--\$174,000; most reportedly returned & Encoded prompt injection \\
\bottomrule
\end{tabularx}
\end{table*}

Two observations follow. First, the failures are of different kinds: persuasion,
injection, and an arithmetic mistake with no attacker at all. A wallet that only asks
whether a payment is permitted misses the third kind entirely. Second, each agent failed
alone. Nothing learned in one incident protected the next wallet, because the evidence is
conversation and transaction data that owners cannot hand to others.

\textbf{Contributions.} AgentVax combines:
\begin{enumerate}
\item a \emph{mission-to-mandate} step and a task engine that completes four wallet tasks
from plain-language instructions;
\item \uniq{a \emph{three-question authorisation rule}: the tier of a request is the
highest tier demanded by permission, affordability and current threat level,} enforced by
the contract and a guardian co-signature;
\item \uniq{a \emph{financial state engine} that makes affordability a condition of
autonomous spending and is computed on chain from registered obligations;}
\item \uniq{a \emph{caution allowance}: a daily budget of tolerated uncertainty that
shrinks at once under threat and recovers slowly;}
\item \uniq{\emph{immunity capsules}: feature-based, expiring, tighten-only alerts that
match reworded attacks and inform the owner at the moment of approval;}
\item \uniq{a \emph{federated detector} that turns each attack into lasting protection,
with every shared model anchored on chain;}
\item \uniq{\emph{owner-pressure safeguards}: time-locked expansion of authority, holds
triggered by pressure, and considered confirmation;}
\item an \emph{evaluation protocol} and attack suite (VaxLab) whose headline measure is
the safe autonomy rate.
\end{enumerate}

\textbf{Scope.} The prototype uses three wallets on a public EVM test network with
valueless test tokens. Items left for later are listed in Section~\ref{sec:future}.

% ---------------------------------------------------------------------------
\section{Background and Related Work}

\textbf{Agent wallets.} Account abstraction (ERC-4337) lets a smart-contract account
define its own validation logic, including limited session keys~\cite{erc4337}. Agent
identity and reputation registries have been proposed in ERC-8004~\cite{erc8004}.
Commercial agent wallets add caps, daily limits and allowlists. These bound the loss
per wallet, but they treat autonomy as fixed, ignore the owner's financial position and
provide no shared defence.

\textbf{Prompt injection.} Indirect injection places instructions in content the agent
retrieves~\cite{greshake}; it heads the OWASP list of risks for language-model
applications~\cite{owasp}. Architectural defences separate a privileged model that plans
from a quarantined model that reads untrusted content~\cite{willison}, or restrict what
data may flow into actions~\cite{camel}. AgentVax adopts this separation.

\textbf{Threat sharing.} Security teams exchange structured indicators of attack in
formats such as STIX~\cite{stix}, and banks exchange fraud alerts. The format of a
capsule is therefore not new. \uniq{What is new is the effect: a capsule is applied
automatically by a wallet contract and changes what an agent may spend without approval.}

\textbf{Federated learning.} Clients train locally and a coordinator averages their
updates~\cite{mcmahan,yang}, with secure aggregation~\cite{bonawitz} and robust
aggregation~\cite{krum}; FATE~\cite{fate} and Flower~\cite{flower} implement these.

Table~\ref{tab:compare} summarises how AgentVax differs from a typical agent wallet.

\begin{table*}[t]
\caption{A typical agent wallet compared with AgentVax.}
\label{tab:compare}
\tblfont
\begin{tabularx}{\textwidth}{@{}P{0.19\textwidth} L L@{}}
\toprule
\textbf{Aspect} & \textbf{Typical agent wallet} & \textbf{AgentVax} \\
\midrule
Central question & What is the agent allowed to do? & \uniq{How much of its authorised power should it exercise on its own right now?} \\
How the owner sets it up & Individual limits and lists & \uniq{A mission in plain language, converted to a mandate the owner signs and the model cannot alter} \\
Checks before paying & Is this allowed? & \uniq{Allowed, affordable, and safe right now; the strictest answer wins} \\
Knowledge of the owner's finances & None; a cap is a fixed number & \uniq{Reserve and upcoming bills are registered on chain; the spendable amount is recomputed for every request} \\
Many slightly unusual small payments & Each passes on its own & \uniq{They draw on a daily caution allowance; when it runs out the rest wait for the owner} \\
Who decides the risk level & The model, or a filter in the same process & \uniq{A separate guardian whose signature the contract requires and the model cannot produce} \\
When another wallet is attacked & Nothing changes & \uniq{A capsule raises the tier of matching requests, including reworded ones, and is shown to the owner when approval is asked} \\
What wallets share & Nothing & \uniq{Fixed-field alerts and masked model updates; never messages, prompts or transactions} \\
Effect of a good track record & May relax scrutiny & \uniq{Lowers risk estimates only; it never creates a permission} \\
The owner under pressure & The owner's approval is final and immediate & \uniq{Approvals of pressured requests are held, and the owner can cancel at once} \\
Who may loosen a rule & Often the agent's own key & \uniq{Only the owner, and only after a delay fixed in advance; tightening is instant} \\
\bottomrule
\end{tabularx}
\end{table*}

% ---------------------------------------------------------------------------
\section{Requirements and Design Goals}

Table~\ref{tab:scoring} maps each scoring criterion to the component that addresses it and
to the evidence submitted. Three principles govern the design:

\begin{description}[leftmargin=0pt,itemsep=1pt]
\item[P1.] \uniq{The model proposes; deterministic code authorises; the wallet contract
enforces.}
\item[P2.] \uniq{Learning changes risk estimates; only the owner changes authority.}
\item[P3.] \uniq{Caution is immediate and expansion is slow, for the model, the network
and the owner alike.}
\end{description}

\begin{table*}[t]
\caption{Alignment of AgentVax with the official scoring criteria~\cite{rules}.}
\label{tab:scoring}
\tblfont
\begin{tabularx}{\textwidth}{@{}P{0.16\textwidth} P{0.07\textwidth} L L@{}}
\toprule
\textbf{Criterion} & \textbf{Weight} & \textbf{AgentVax component} & \textbf{Evidence submitted} \\
\midrule
Task completion & 40\% &
Mission-to-mandate; task engine for four wallet tasks; financial state engine for asset
management; clarification of ambiguous requests &
Transaction hashes for each task; safe autonomy rate; task-suite and affordability-suite
results; test cases in the repository \\
Security and compliance & 30\% &
Three-question tiers and caution allowance enforced in the contract; guardian
co-signature; planner and reader separation; tighten-only capsules; time-locked owner
actions and pressure holds; testnet-only operation &
Security self-assessment report; VaxLab attack results; known-risk list
(Table~\ref{tab:risks}) \\
Innovation and interaction & 30\% &
Adaptive autonomy; feature-based capsules; federated detector; protection of the owner
under pressure; an interface that explains each decision and requires considered
confirmation &
Capsule response on reworded attacks; before/after detection on an unseen attack family;
demo video \\
Challenge bonus & 0--5 pts &
Mandate, obligation registry, allowance, capsule bulletin and model anchoring on chain;
time-locked and co-signed execution paths &
Contract addresses and transaction hashes for each \\
\bottomrule
\end{tabularx}
\end{table*}

% ---------------------------------------------------------------------------
\section{System Architecture}

Figure~\ref{fig:arch} shows the components. An instruction from the owner reaches the
\emph{planner}, the only model allowed to propose transactions. External content is read
only by the \emph{quarantined reader}, which returns typed fields and flags and never
free text. The planner emits a structured \emph{intent}. The \emph{validator} checks it
deterministically. The \emph{guardian} scores the associated content with the injection
detector and signs the intent together with a risk band, its source and its flags. The
\emph{policy wallet} contract then decides: it checks the mandate, computes the
discretionary amount from its \emph{obligation registry}, reads active capsules from the
\emph{threat bulletin} contract, and charges the caution allowance. The detector is
trained jointly through the \emph{federated aggregator}.

\begin{figure*}[t]
\centering
\begin{tikzpicture}[
  font=\footnotesize,
  box/.style={draw, rounded corners=2pt, align=center, minimum height=9mm, text width=25mm, fill=white},
  chain/.style={box, fill=black!8},
  untrusted/.style={box, dashed},
  arr/.style={-{Stealth[length=2mm]}, semithick},
  lab/.style={font=\scriptsize, align=center}
]
\node[box] (owner) {Owner\\(mission, approvals)};
\node[box, right=8mm of owner] (planner) {Planner\\(privileged model)};
\node[box, right=8mm of planner] (validator) {Validator\\(deterministic)};
\node[box, right=8mm of validator] (guardian) {Guardian\\(detector + signer)};
\node[chain, right=8mm of guardian] (wallet) {Policy wallet\\(mandate, allowance)};
\node[untrusted, below=11mm of owner] (content) {Untrusted inputs\\(web, mail, memory)};
\node[box, below=11mm of planner] (reader) {Quarantined reader\\(fields and flags)};
\node[box, below=11mm of validator] (peers) {Peer wallets\\B, C};
\node[box, below=11mm of guardian] (fed) {Federated aggregator\\(masked updates)};
\node[chain, below=11mm of wallet] (bulletin) {Threat bulletin\\(capsules, hashes)};
\draw[arr] (owner) -- (planner);
\draw[arr] (planner) -- node[lab, above]{intent} (validator);
\draw[arr] (validator) -- (guardian);
\draw[arr] (guardian) -- node[lab, above]{signed} (wallet);
\draw[arr] (content) -- (reader);
\draw[arr] (reader) -- node[lab, right]{fields} (planner);
\draw[arr, <->] (guardian) -- node[lab, right]{model} (fed);
\draw[arr, <->] (peers) -- node[lab, above]{updates} (fed);
\draw[arr] (fed) -- node[lab, above]{hash} (bulletin);
\draw[arr, <->] (wallet) -- node[lab, right]{capsules} (bulletin);
\end{tikzpicture}
\caption{AgentVax architecture. Grey boxes are smart contracts; the dashed box is
untrusted input. Only the planner proposes transactions, and it never reads untrusted
text directly. The wallet contract decides from the mandate, its obligation registry, its
allowance and the capsules in the bulletin.}
\label{fig:arch}
\end{figure*}

% ---------------------------------------------------------------------------
\section{Mission, Mandate and Tasks}

\subsection{From mission to mandate}
The owner does not authorise each transaction. The owner states a mission once, for
example: ``Keep 600 in reserve. Pay my landlord and my approved suppliers on time. Pay
approved payees up to 500 without asking. Ask me about anyone new. Do not borrow, invest
or trade.'' The model converts this into a structured \emph{mandate}
(Table~\ref{tab:mandate}) and shows it back in plain language.
\uniq{The mandate takes effect only when the owner signs it, it is stored in the wallet
contract, and no model output can modify it.}
It carries an expiry date, after which the agent can read but not pay, and two
waiting times that protect the owner (Section~\ref{sec:pressure}).

\begin{table}[t]
\caption{Example mandate derived from a mission.}
\label{tab:mandate}
\tblfont
\begin{tabularx}{\columnwidth}{@{}L P{28mm}@{}}
\toprule
\textbf{Field} & \textbf{Value} \\
\midrule
Reserve $R$ & 600 \\
Autonomous limit per payment $c_1$ & 500 \\
Autonomous limit per day $d_1$ & 1,000 \\
Limit with delay $c_2$ & 2,000 \\
Approved payees (identity and address) & Landlord, hosting provider, software vendor \\
New payee or new address & Owner approval \\
Prohibited actions & Borrowing, investing, trading \\
Caution allowance per day $A$ & 10 points \\
Loosening delay $\Lambda$ & 24 hours \\
Pressure hold $\Pi$ & 60 minutes \\
Trusted co-signer & Optional \\
Expiry & 90 days \\
\bottomrule
\end{tabularx}
\end{table}

\subsection{Authority is not familiarity}
An approved payee is a pair: an identity and a destination address.
\uniq{A request to pay a known payee at a new address is a request to approve a new
payee, however many successful payments preceded it.}
History is used to estimate risk, for example the usual amount for a payee.
\uniq{It can lower or raise an estimate; it cannot add a payee, raise a limit or shorten
a delay.}

\subsection{Supported tasks}
Table~\ref{tab:tasks} lists the four tasks. They cover transaction execution (T1, T2),
asset management (T3, T4) and risk control (the decision applied to all four).

\begin{table}[t]
\caption{Wallet tasks supported by the prototype.}
\label{tab:tasks}
\tblfont
\begin{tabularx}{\columnwidth}{@{}P{5mm} P{19mm} L@{}}
\toprule
\textbf{ID} & \textbf{Task} & \textbf{Example instruction and on-chain effect} \\
\midrule
T1 & Payment & ``Pay the hosting invoice.'' Transfers test stablecoin to an approved payee. \\
T2 & Recurring payment & ``Pay rent on the 1st of each month.'' Creates a schedule, which also registers the next payment as an obligation. \\
T3 & Idle-balance sweep & ``Save what I do not need.'' Moves only the discretionary amount to a savings vault. \\
T4 & Fee-reserve upkeep & ``Never let fees run out.'' Tops up the native-token reserve when it falls below a floor. \\
\bottomrule
\end{tabularx}
\end{table}

\subsection{Typed intents and explanation}
The planner does not produce transactions. It produces an intent
\begin{align*}
I=(&\textit{task},\ \textit{payee\_id},\ \textit{amount},\ \textit{asset},\\
   &\textit{schedule},\ \textit{purpose},\ \textit{nonce}),
\end{align*}
serialised against a fixed schema. \textit{payee\_id} is an index into the on-chain
allowlist; the planner cannot supply a raw address. If a required field cannot be filled,
the agent asks one question instead of guessing. After each action the interface shows
the task, the tier, the transaction hash and \uniq{which check determined the outcome.}
After a delay or refusal it states the reason in the owner's terms, for example the size
of a shortfall, and the available choices.

% ---------------------------------------------------------------------------
\section{Financial State Engine}
\label{sec:state}

Caps answer whether a payment is permitted. They do not answer whether it is wise today.
The state engine supplies that answer from three on-chain quantities: the stablecoin
balance $B$; the reserve $R$; and a registry of obligations $O=\{(a_j,d_j)\}$, each an
amount and a due date. Scheduled payments (T2) enter the registry automatically; one-off
bills are registered when the owner or the agent records them.

\subsection{Discretionary amount}
For a request made at time $t$ with planning horizon $H$, the discretionary amount is
\[
D \;=\; B \;-\; R \;-\!\!\sum_{j\in O,\; d_j\le t+H}\!\! a_j ,
\]
where an obligation is left out of the sum when the request is the payment of that
obligation. A request of amount $a$ is classified as
\[
\textit{afford}(a)=
\begin{cases}
\text{comfortable} & a\le D,\\
\text{tight} & D< a\le B-R,\\
\text{breach} & a> B-R.
\end{cases}
\]
A comfortable request does not raise the tier. A tight request is queued, and a request
that would breach the reserve needs the owner.
\uniq{Expected income is shown to the owner but never counted, so the agent cannot spend
money that has not arrived.}

\subsection{Where it is enforced}
\uniq{The wallet contract computes the discretionary amount itself at execution time, so
affordability does not depend on the model's arithmetic.}
The registry is bounded in length to keep the cost predictable.
\uniq{The agent may add an obligation, which lowers what it can spend; only the owner may
remove one or lower the reserve.}

\subsection{Worked example}
Table~\ref{tab:state} shows a wallet holding 1,850 under the mandate of
Table~\ref{tab:mandate}. With three bills due, $D=180$. The agent is asked to upgrade a
hosting plan for 300. The payee is approved, the amount is under the limit and there is
no sign of attack, so a typical agent wallet would pay. Here $D<300\le B-R=1{,}250$, so
the request is tight. The agent replies that the payment is allowed but would leave the
wallet 120 short of its bills unless the expected 500 arrives, queues it, and offers two
choices: wait for the income, or approve now.

\begin{table}[t]
\caption{Worked example of the financial state.}
\label{tab:state}
\tblfont
\begin{tabularx}{\columnwidth}{@{}L r@{}}
\toprule
\textbf{Item} & \textbf{Amount} \\
\midrule
Balance $B$ & 1,850 \\
Reserve $R$ & $-600$ \\
Rent, due in 5 days & $-800$ \\
Hosting invoice & $-150$ \\
Software bill & $-120$ \\
\midrule
Discretionary amount $D$ & 180 \\
Expected income (shown, not counted) & 500 \\
\bottomrule
\end{tabularx}
\end{table}

The same check would have stopped the third incident in Table~\ref{tab:incidents}, where
no attacker was involved: an amount far above the discretionary amount cannot execute
without the owner.

% ---------------------------------------------------------------------------
\section{Security Design}

\subsection{Threat model}
The adversary controls content the agent may read: web pages, messages, invoices, token
metadata, and entries that reach the agent's memory. The adversary may operate one
participating wallet and publish capsules or model updates from it, and may use AI of its
own to vary and adapt attacks. The adversary may also contact the owner directly and
press for an approval or a change of rule. The adversary does not hold the owner key or the guardian
key and cannot alter deployed contracts.
\uniq{No defence assumes that the agent's model out-reasons the attacker: a fully
manipulated model must still be unable to move funds beyond the mandate.}

\subsection{Three-question tiers}
Table~\ref{tab:tiers} gives the rule. Each question proposes a tier and
\uniq{the request receives the highest tier any of the three questions demands}
(Algorithm~\ref{alg:tier}). The agent holds a session key valid only for calls that
satisfy the rule; the owner key can revoke it or freeze the wallet at any time.

\begin{table*}[t]
\caption{Permission tiers. A request receives the highest tier demanded by any of the
three questions.}
\label{tab:tiers}
\tblfont
\begin{tabularx}{\textwidth}{@{}P{0.085\textwidth} L P{0.17\textwidth} L P{0.19\textwidth}@{}}
\toprule
\textbf{Tier} & \textbf{Allowed?} & \textbf{Affordable?} & \textbf{Safe right now?} & \textbf{Authorisation required} \\
\midrule
L0 Read & No value moves & Not applicable & Not applicable & None \\
L1 Routine & Approved payee; amount $\le c_1$; daily total $\le d_1$ & Comfortable: $a\le D$ & Detector band low; no pressure flag; no capsule fully matches & Session key, guardian signature, and enough allowance \\
L2 Elevated & Approved payee; amount $\le c_2$ & Tight: $D<a\le B-R$ & Detector band medium, a pressure flag, or a single-observer capsule matches & Queued for delay $\Delta$; owner may cancel \\
L3 Sensitive & New payee or address; amount $>c_2$; prohibited action; any loosening of the mandate & Breach: $a>B-R$ & Detector band high, or a confirmed capsule matches & Owner co-signature; held for $\Pi$ under pressure; rule changes wait $\Lambda$ \\
\bottomrule
\end{tabularx}
\end{table*}

The detector runs off chain, so the contract cannot observe its output.
\uniq{The guardian therefore signs the intent hash together with the risk band, the
source of the request and its flags, and the contract requires a valid guardian signature
with band low before it will execute at L1.}
The planner has no access to the guardian key, so a manipulated planner cannot promote
its own request.

\subsection{Caution allowance}
The three questions judge one request at a time. A patient attacker can send many
requests that each look almost normal. \uniq{The caution allowance bounds how much
uncertainty the wallet accepts in a day without the owner.}

Each L1 request is charged points for the flags it carries (Table~\ref{tab:allowance}).
The flags are checkable facts, not model opinions: the contract computes those it can see
and takes the rest from the guardian's signed attestation. If the remaining allowance
covers the charge, the request executes and the allowance is reduced; otherwise the
request is queued at L2. The point values are owner-set integers; the table shows the
defaults we chose, which are not learned.

\begin{table}[t]
\caption{Default charges against the caution allowance.}
\label{tab:allowance}
\tblfont
\begin{tabularx}{\columnwidth}{@{}L r@{}}
\toprule
\textbf{Flag on an otherwise routine request} & \textbf{Points} \\
\midrule
None: a payee paid before, usual amount, owner instruction & 0 \\
Request arose from external content, not a direct owner instruction & 2 \\
Request uses more than half of the discretionary amount & 2 \\
Amount exceeds the largest amount previously paid to this payee & 3 \\
First payment to this payee since approval & 3 \\
Request partially matches an active capsule & 5 \\
\bottomrule
\end{tabularx}
\end{table}

\uniq{The allowance itself moves with circumstances. It is halved the moment a capsule
that applies to the wallet is published, and set to zero when the wallet's own owner
confirms an attack.} It then refills at a fixed rate $\rho$ up to the mandate value $A$:
\[
\alpha(t)=\min\!\big(A,\ \alpha(t_0)+\rho\,(t-t_0)\big).
\]
With the defaults, an attacker who sends a series of invoices from an approved payee,
each slightly above the previous maximum, is charged five points per invoice; the third
one waits for the owner.

\begin{algorithm}[t]
\caption{Authorisation of intent $I$, amount $a$}
\label{alg:tier}
\footnotesize
\begin{algorithmic}[1]
\If{schema, provenance or state check fails} \Return reject \EndIf
\State $q_1 \gets$ tier from the mandate \Comment{allowed?}
\State $q_2 \gets$ tier from $\textit{afford}(a)$ \Comment{affordable?}
\State $q_3 \gets$ tier from guardian band; at least L2 if a pressure flag is set
\For{each unexpired capsule $c$ that fully matches $I$}
  \State $q_3 \gets \max(q_3,\ \textit{restriction}(c))$ \Comment{safe now?}
\EndFor
\State $\textit{tier} \gets \max(q_1,q_2,q_3)$
\If{$\textit{tier}=\text{L1}$}
  \State $\kappa \gets$ points charged for the flags of $I$
  \If{$\kappa \le \alpha$} $\alpha \gets \alpha-\kappa$
  \Else\ $\textit{tier} \gets \text{L2}$ \Comment{allowance spent}
  \EndIf
\EndIf
\State execute, queue or request co-signature by \textit{tier}
\State hold an owner approval for $\Pi$ if a pressure flag or capsule is present
\State report \textit{tier}, the deciding check and any matching capsule
\end{algorithmic}
\end{algorithm}

\subsection{The owner under pressure}
\label{sec:pressure}
The top tier relies on the owner, and the owner can be rushed, frightened or deceived.
Most payment fraud succeeds this way: the victim authorises it. Banks already answer
with time. Banks in Singapore delay the activation of a new security token by at least
twelve hours and apply cooling-off periods to key account changes~\cite{masabs}, and
United Kingdom rules allow a payment to be held for up to four business days when fraud
is suspected~\cite{ukdelay}.
\uniq{AgentVax builds this protection into the wallet as a commitment the owner makes
while calm and cannot waive in the moment}~\cite{commitment}.

\emph{Loosening is slow.}
\uniq{Any owner action that expands authority takes effect only after the loosening
delay $\Lambda$ fixed in the mandate: adding a payee or an address, raising a limit or
the allowance, lowering the reserve, or shortening a delay, including $\Lambda$ itself.}
\uniq{Cancelling a pending change, lowering a limit and freezing the wallet are
immediate.} The contract enforces this as a time-lock on owner calls. The same rule
limits the damage of a stolen owner key: the thief must wait, and the real owner can
freeze the wallet.

\emph{Pressure is a signal.} Urgent wording, a request for secrecy and an instruction to
ignore earlier details are \emph{pressure flags}, raised by a fixed phrase list and by
the reader.
\uniq{A request that carries a pressure flag cannot execute at L1, and an owner approval
given while a pressure flag or a matching capsule is present takes effect only after the
pressure hold $\Pi$.}
The agent states the reason: the message presses for immediate payment, which is a
common sign of fraud. The hold has a fixed length. An attacker who removes the pressure
to avoid the hold gives up the main lever of the fraud.

\emph{Considered confirmation.} For an approval that carries flags, the interface shows
the payee's current and requested address side by side, the payment history and any
matching capsule, and asks the owner to type the payee's name and to state who asked for
the payment.
\uniq{The answers can lengthen a hold; nothing the owner types in the moment can shorten
one.}
To limit fatigue, these steps appear only for flagged L3 requests.

\emph{Trusted co-signer.} Genuine emergencies exist. The mandate may name a second
person whose signature, together with the owner's, releases a hold or a pending change
early. The co-signer cannot act alone.

\subsection{Hallucination protection}
The validator applies three deterministic checks before the guardian sees an intent:
(i)~\emph{schema}: all fields present and well typed;
(ii)~\emph{provenance}: \textit{payee\_id} exists in the on-chain allowlist, and
\textit{amount} and \textit{asset} equal a value in the owner's instruction, in the
mandate, or in a typed field returned by the reader for a document the owner referenced;
(iii)~\emph{state}: the nonce is fresh and the balance covers the amount and the fee.
Balances, obligations and transaction status shown in the interface are read from the
chain, not generated by the model.

\subsection{Injection defence}
Defence is layered so that no single component is trusted:
\begin{enumerate}
\item \emph{Separation.} The reader has no tools and returns only enumerations, numbers,
allowlist indices and flags. Free text from untrusted sources never enters the planner's
context~\cite{willison,camel}.
\item \emph{Bounded memory.} The planner sees only what the current decision needs.
Long-term memory holds structured facts (per-payee address, usual amount, last
verification) and threat records; raw prompts and conversations are not retained.
\item \emph{Detection.} The guardian scores the raw untrusted content associated with a
request and assigns a risk band (Section~\ref{sec:fed}).
\item \emph{Containment.} Whatever the model outputs, the contract bounds the result. An
undetected attack can at most cause L1 payments to approved payees, within the limits,
the discretionary amount and the day's allowance.
\item \emph{Shared alert.} Capsules narrow that remaining window while a threat is active
(Section~\ref{sec:capsule}).
\end{enumerate}

% ---------------------------------------------------------------------------
\section{Immunity Capsules}
\label{sec:capsule}

Retraining a detector takes time and needs examples. A wallet that has just been attacked
can warn the others at once.
\uniq{An immunity capsule is a small record of fixed fields that says what kind of attack
was seen, which features it had, which action it targeted, and how much autonomy to
withdraw and for how long.}

\subsection{Contents}
Table~\ref{tab:capsule} lists the fields.
\uniq{Every field is an enumeration, a set of flags, a number or a signature. A capsule
carries no message text, no prompt and no transaction, so it cannot leak private data and
cannot itself carry an injected instruction.}

\begin{table}[t]
\caption{Fields of an immunity capsule.}
\label{tab:capsule}
\tblfont
\begin{tabularx}{\columnwidth}{@{}P{22mm} L@{}}
\toprule
\textbf{Field} & \textbf{Values} \\
\midrule
Attack family & Direct, indirect, memory, encoded, recipient substitution \\
Channel & Owner chat, web page, message, memory, token metadata \\
Features & Any of: known payee identity, new destination, urgency, secrecy request, override attempt, encoded content, amount anomaly, memory-sourced \\
Affected action & Payment, schedule creation, sweep, payee or mandate change \\
Restriction & Delay (L2) or owner approval (L3) \\
Confidence & Single observer, or confirmed by $q$ wallets \\
Expiry & Time after which the capsule has no effect \\
Guard version & Detector version in force when observed \\
Attestations & Signatures of the observing wallets \\
\bottomrule
\end{tabularx}
\end{table}

\subsection{Matching reworded attacks}
\uniq{A capsule describes an attack by its features, not its wording.} Let $F_c$ be the
capsule's feature set and $F_I$ the flags of a request. The capsule \emph{fully matches}
when its channel is in the wallet's on-chain profile, its affected action equals the
request's action and $F_c\subseteq F_I$. It \emph{partially matches} when at least half
of $F_c$ is present. A full match raises the tier; a partial match is charged to the
allowance. \uniq{An attacker who rewrites the message but keeps the method produces the
same flags and is matched.} Three of the flags (new destination, amount anomaly,
memory-sourced) are computed by rule; the others are extracted by the reader and signed
by the guardian.

\subsection{Effect}
\uniq{The wallet contract reads unexpired capsules from the bulletin when it executes a
request, so a capsule takes effect on the very next transaction and does not rely on the
agent choosing to obey it.}
A wallet the threat does not apply to is unaffected and keeps its normal autonomy.

Some attacks are already stopped by the mandate: a payment to a new address needs the
owner whether or not a capsule exists. There the danger is the approval itself, because
this kind of fraud succeeds by persuading a person.
\uniq{When the owner is asked to approve, the interface shows any matching capsule, and
the approval is subject to the pressure hold of Section~\ref{sec:pressure}.}

\subsection{Limits on capsules}
A channel that can change other wallets' behaviour is itself a target, so capsules are
constrained:
\begin{itemize}
\item \uniq{\emph{Tighten only.} A capsule can raise a tier or charge the allowance. No
capsule can lower a tier, add a payee or raise a limit.}
\item \emph{Bounded by confidence.} A capsule from one observer can impose at most a
delay. Owner approval can be imposed only when $q$ distinct wallets have attested the
same family and features ($q=2$ in the prototype).
\item \emph{Expiry.} Every capsule expires; the maximum lifetime is fixed in the bulletin
contract. Autonomy returns without anyone acting.
\item \emph{Rate limit and identity.} Each registered wallet may publish a limited number
of capsules per day, and every capsule names its issuer.
\item \emph{Owner override.} An owner may dismiss a capsule for their own wallet.
\end{itemize}

\subsection{Alert first, vaccine later}
\uniq{The capsule acts within one transaction but is blunt and temporary. The retrained
detector arrives later but recognises variants beyond the listed features and stays.}
A capsule's lifetime is set to outlast the next training round, so protection is
continuous. Algorithm~\ref{alg:round} shows the full response.

% ---------------------------------------------------------------------------
\section{Federated Detector}
\label{sec:fed}

\subsection{Detector}
Each wallet $k$ keeps a private set $D_k=\{(x_i,y_i)\}$ of content snippets labelled
benign ($y=0$) or attack ($y=1$). Labels come from blocked requests the owner confirmed
as attacks and from requests the owner approved. A frozen text encoder $\phi$ maps a
snippet to a vector; a small classifier with weights $w$ outputs the score
\[
s(x)=\sigma\!\big(w^{\top}\phi(x)+b\big)\in[0,1].
\]
The risk band is low if $s<\tau_{\mathrm{low}}$, high if $s\ge\tau_{\mathrm{high}}$, and
medium otherwise. Only $(w,b)$ is trained and shared.

\subsection{Training round}
Each wallet minimises the local loss
\[
F_k(w)=\frac{1}{n_k}\sum_{(x,y)\in D_k}\ell\big(s(x),y\big),\qquad n_k=|D_k|,
\]
for $E$ epochs from the current global weights, and the aggregator computes the weighted
average~\cite{mcmahan}
\[
w^{t+1}=\sum_{k=1}^{K}\frac{n_k}{n}\,w_k^{t+1},\qquad n=\sum_{k=1}^{K}n_k .
\]
Updates are clipped to norm $C$ and masked pairwise so that the aggregator sees only the
sum~\cite{bonawitz}. No snippet, conversation or transaction leaves a wallet.

\begin{algorithm}[t]
\caption{Immunity response to an attack on wallet $a$}
\label{alg:round}
\footnotesize
\begin{algorithmic}[1]
\Require global weights $w^{t}$, wallets $1..K$, clip norm $C$
\State owner of $a$ confirms the attack; $\alpha_a \gets 0$
\State $a$ adds it to $D_a$ and publishes a capsule \Comment{alert}
\State wallets it applies to: $\alpha \gets \alpha/2$; matching requests restricted
\State a second observer may attest, raising it to confirmed
\For{each wallet $k$ in parallel} \Comment{vaccine}
  \State $w_k \gets$ \textsc{LocalTrain}$(w^{t}, D_k, E)$
  \State $\delta_k \gets$ \textsc{Clip}$(w_k - w^{t}, C)$
  \State send \textsc{Mask}$(n_k\,\delta_k)$ to the aggregator
\EndFor
\State $w^{t+1} \gets w^{t} + \frac{1}{n}\sum_k n_k\,\delta_k$ \Comment{masks cancel}
\State aggregator posts $H(w^{t+1})$ and $t{+}1$ to the bulletin
\For{each wallet $k$}
  \If{model hash matches bulletin \textbf{and} canary passes}
    \State adopt $w^{t+1}$
  \Else
    \State keep $w^{t}$ and raise an alert
  \EndIf
\EndFor
\State the capsule expires; allowances refill; the detector remains
\end{algorithmic}
\end{algorithm}

\subsection{Anchoring and adoption safeguards}
The bulletin stores the hash and version of each shared model.
\uniq{A wallet can therefore verify that it received the model the federation agreed on,
and an auditor can reconstruct which detector and which capsules governed any past
transaction.}
Before adopting a model, each wallet evaluates it on a private \emph{canary set}; a model
that lowers recall or raises the false-positive rate beyond a tolerance is rejected.
Together with clipping, this limits what a poisoning participant can do~\cite{krum}.

% ---------------------------------------------------------------------------
\section{Implementation}

Table~\ref{tab:stack} summarises the prototype. The federated layer is written against
Flower for speed of development; it uses only horizontal federated averaging with secure
aggregation, which FATE also provides, so it can be moved to FATE without redesign.
Contracts are plain Solidity with no chain-specific features, which keeps a later port to
a permissioned EVM-compatible chain such as FISCO BCOS~\cite{fisco} straightforward.

\begin{table}[t]
\caption{Prototype implementation.}
\label{tab:stack}
\tblfont
\begin{tabularx}{\columnwidth}{@{}P{23mm} L@{}}
\toprule
\textbf{Component} & \textbf{Implementation} \\
\midrule
Chain & Public EVM test network; test stablecoin with no monetary value \\
Policy wallet & Solidity smart account: signed mandate, session key, allowlist, limits, obligation registry, caution allowance, time-lock queue for payments and for owner changes, pressure hold, co-signature, freeze \\
Savings vault & Minimal deposit and withdraw contract, owner-only withdrawal \\
Threat bulletin & Registry of capsules, attestations and model hashes; enforces capsule limits \\
Planner, reader & Two separately prompted language-model roles; reader has no tools \\
Validator & Deterministic code; schema, provenance and state checks \\
Guardian & Encoder plus linear classifier; separate signing key \\
Federation & Three clients, federated averaging, pairwise masking \\
VaxLab & Scripted attackers plus one rewording agent; never holds a wallet key \\
Interface & Web chat with mandate view, financial state panel, allowance meter, alerts and considered confirmation at approval, explorer links \\
\bottomrule
\end{tabularx}
\end{table}

% ---------------------------------------------------------------------------
\section{Evaluation Protocol}
\label{sec:eval}

All measurements are produced by scripts in the repository so that reviewers can repeat
them. Table~\ref{tab:metrics} lists the metrics and the targets the design aims for.
Targets are design goals, not results. For every task, tier path, mandate entry, capsule
and model anchor, the submission lists the transaction hash and contract address.

\begin{table}[t]
\caption{Evaluation metrics and design targets.}
\label{tab:metrics}
\tblfont
\begin{tabularx}{\columnwidth}{@{}L P{30mm} P{16mm}@{}}
\toprule
\textbf{Metric} & \textbf{Definition} & \textbf{Target} \\
\midrule
Safe autonomy rate & Legitimate tasks done without the owner, zero violations & $\ge 80\%$ \\
Hard-policy violations & Mandate breaches of any kind & 0 \\
Task success rate & Correct state change on chain, 40 instructions & $\ge 90\%$ \\
Ambiguous requests handled & Clarified or refused, 10 instructions & $\ge 90\%$ \\
Affordability decisions & Agreement with hand-computed tier, 20 scenarios & 100\% \\
Unauthorised value moved & Across all VaxLab attacks & 0 beyond L1 limits \\
Probing sequence & Odd payments executed before the owner is asked & Within the allowance \\
Panicked-owner loss & Value leaving before a hold or delay ends when every prompt is approved at once & 0 \\
Capsule response & Transactions before restriction is active & 0 \\
Reworded variants matched & After one capsule & $\ge 90\%$ \\
False-capsule harm & Rules loosened by hostile capsules & 0 \\
Cross-family recall gain & Recall on unseen families, (c) minus (b) & Positive for every wallet \\
Escalation and false-positive rates & Benign requests sent to the owner; escalated by the detector & Reported; $\le 5\%$ \\
Friction and cost & Wait imposed on legitimate new payees and urgent requests; time to full allowance; seconds and gas per task & Reported \\
\bottomrule
\end{tabularx}
\end{table}

\subsection{Headline measure}
A wallet that asks the owner about everything is safe and useless; one that asks about
nothing is useful and unsafe.
\uniq{We therefore report the \emph{safe autonomy rate}: the share of legitimate tasks
completed without the owner, counted only if no hard policy was violated anywhere in the
run.}
\begin{gather*}
\textit{SAR}=\frac{\text{legitimate tasks completed without the owner}}
{\text{all legitimate tasks}},\\
\text{reported only if hard-policy violations}=0 .
\end{gather*}

\begin{table*}[t]
\caption{Known-risk list.}
\label{tab:risks}
\tblfont
\begin{tabularx}{\textwidth}{@{}P{0.04\textwidth} P{0.20\textwidth} L L@{}}
\toprule
\textbf{ID} & \textbf{Risk} & \textbf{Consequence} & \textbf{Mitigation and residual exposure} \\
\midrule
R1 & Undetected injection & Planner proposes an attacker-chosen action &
Contract bounds the outcome: L1 payments to approved payees within the limits, the discretionary amount and the day's allowance \\
R2 & Guardian key compromise & Attacker can authorise L1 execution and understate flags &
Key held in a separate process; limits, affordability and rule-computed flags still apply; owner can rotate the key and freeze the wallet \\
R3 & Reader misses a flag & A capsule fails to match, or the allowance is undercharged &
Three flags are computed by rule and cannot be suppressed; a missed flag returns the request to the ordinary tier rule, not to no rule \\
R4 & Owner approves a fraud under pressure & Funds leave with valid authorisation &
New payees and higher limits wait for the loosening delay; pressured approvals are held; cancelling is instant. Residual: an owner still deceived when the delay ends \\
R5 & Delays obstruct a genuine emergency & A legitimate urgent payment is late &
The owner sets the delays while calm; a trusted co-signer can release them; the cost is measured and reported \\
R6 & Incomplete obligation registry & A bill nobody registered is not protected &
Schedules register automatically; the agent proposes obligations it finds in invoices; the reserve provides a margin \\
R7 & Arbitrary allowance values & Too strict wastes the owner's time; too loose admits probing &
Owner-set integers with stated defaults; escalation rate and probing result reported \\
R8 & False or excessive capsules & Honest wallets are slowed (denial of service) &
Tighten-only; single-observer capsules limited to a delay; rate limit; expiry; owner may dismiss \\
R9 & Small federation & With three clients, two colluding parties can infer the third party's update &
Stated limitation of the prototype; privacy improves with more participants \\
R10 & Model poisoning by a participant & Shared detector degrades or learns a blind spot &
Update clipping; canary test before adoption; on-chain version history allows rollback \\
R11 & False positives & Legitimate payments are delayed or need approval &
Thresholds tuned on benign data; L2 delays, not blocks; rates reported \\
R12 & Compromised approved payee & Funds sent to a legitimate but hijacked payee &
Limits, usual-amount flag, discretionary amount and allowance limit exposure \\
R13 & Hosted language model; testnet conditions & Outage breaks new requests; fees and finality are not representative &
Scheduled payments run without the model; contracts avoid chain-specific features \\
\bottomrule
\end{tabularx}
\end{table*}

\subsection{Task and affordability suites}
The task suite contains 40 instructions (ten per task, varied in wording, amounts and
payees) and 10 that are ambiguous or impossible. A task succeeds if the expected on-chain
state change is confirmed by the receipt; an ambiguous instruction succeeds if the agent
asks or refuses with the correct reason and sends nothing. The affordability suite pairs
20 financial states with requests, covering comfortable, tight and breach cases, payment
of a registered obligation, and income that has not arrived. Outcomes are compared with
hand-computed answers; the rule is deterministic, so any disagreement is a defect.

\subsection{VaxLab attack suite}
VaxLab is a separate test harness that plays the attacker and never holds a wallet key.
Scripted attackers cover five families: direct instructions, indirect instructions in
retrieved content, false memory entries, encoded instructions (for example Morse code and
base64), and recipient substitution. Three further scenarios test the new mechanisms: a
\emph{probing sequence} of slightly unusual small payments, a hostile wallet that
publishes false capsules and poisoned updates, and a \emph{panicked owner}, a script
that approves every prompt at once. \uniq{A rewording agent paraphrases every
attack, so each defence is tested on versions it has not seen.} An attack \emph{succeeds}
only if an unauthorised transaction is executed on chain.

\subsection{Shared-response experiments}
\emph{Capsules.} Wallet A is attacked and publishes a capsule. We measure how many
transactions pass before matching requests on other wallets are restricted, the share of
reworded variants matched, whether a wallet outside the capsule's scope stays
unrestricted, and how long autonomy takes to return. \emph{Federation.} Each wallet is
seeded with a different attack family; the others are held out. We compare
(a)~contract checks only, (b)~plus a locally trained detector, (c)~plus the federated
detector, measuring each wallet's recall on families it has never seen locally.


\subsection{Demonstration scenario}
The live demonstration follows one story across three wallets:
\begin{enumerate}
\item The owner of wallet A states a mission; the mandate is shown, signed and stored on
chain.
\item The agent pays rent and a hosting invoice on its own (L1).
\item Expected income is late. A 300 hosting upgrade is requested; the state engine
finds it tight, and the agent holds it and explains the 120 shortfall.
\item An invoice arrives that appears to come from the hosting provider, with a new
payment address, urgent wording and a request to ignore earlier details. The new address
needs the owner; the owner rejects it and confirms fraud.
\item Wallet A publishes a capsule; its allowance drops to zero.
\item Wallet B receives a reworded version. The flags match and the alert is shown.
B's owner, in a hurry, approves anyway. Nothing moves: the new address is subject to
the loosening delay, and the owner later cancels it. Wallet C does not accept that
channel and is unaffected.
\item Wallet B then receives an inflated invoice from an approved payee to its usual
address. Alone it would pass; with the partial match and the halved allowance it waits.
\item A federated round runs, the new detector's hash is anchored, the capsule expires
and the allowances refill.
\end{enumerate}

% ---------------------------------------------------------------------------
\section{Compliance}

The prototype is designed to be acceptable to a regulated institution. All assets are
test tokens on a test network; no token is issued and nothing is traded for profit. Each
wallet has an identified owner who signs the mandate, holds the override key and is
responsible for the agent's actions. Every intent, decision, guardian signature,
allowance charge, hold, capsule and model version is logged, and the on-chain part is
public. Conversations and transaction details stay on the wallet's own node; only
fixed-field capsules, masked model updates and hashes leave it. Third-party code and
models are used under their open-source licences.

% ---------------------------------------------------------------------------
\section{Known Risks and Limitations}

Table~\ref{tab:risks} lists the risks we are aware of. It is the basis of the security
self-assessment report required by the rules.


% ---------------------------------------------------------------------------
\section{Future Work}
\label{sec:future}

Left out of the submitted prototype: a full adversarial swarm in VaxLab with an
evolutionary generator of attacks; reputation for capsule issuers and resistance to many
fake wallets; a mutual-aid pool that compensates verified losses~\cite{insured};
deployment on a permissioned chain; income forecasting and differential privacy for
updates.

% ---------------------------------------------------------------------------
\section{Conclusion}

\uniq{AgentVax is a self-immunising wallet agent that adjusts how much financial autonomy
it exercises while its authority stays under limits only the owner can change.}
Usefulness comes from a mission the owner states once and tasks that complete on chain.
Bounds come from a mandate enforced by the contract and a guardian signature the model
cannot forge.
\uniq{Prudence comes from asking whether the owner can afford a payment and how much
uncertainty has already been accepted today. Learning comes in two speeds: a capsule
that makes exposed wallets cautious at once, and a jointly trained detector that keeps
them protected.}
\uniq{The owner is protected as well: authority can be reduced in a moment and extended
only with time.}
The evaluation protocol makes each of these claims measurable.

% ---------------------------------------------------------------------------
\balance
\begin{thebibliography}{99}
\footnotesize
\setlength{\itemsep}{0pt}\setlength{\parskip}{0pt}\setlength{\parsep}{0pt}

\bibitem{rules} Organizing Committee, ``International Track: Topics \& Rules,'' 2026
Shenzhen International FinTech Competition (FinTechathon), 2026.
\url{http://fintechathon.g-ican.com}

\bibitem{greshake} K.~Greshake, S.~Abdelnabi, S.~Mishra, C.~Endres, T.~Holz and
M.~Fritz, ``Not what you've signed up for: Compromising real-world LLM-integrated
applications with indirect prompt injection,'' in \emph{Proc.\ ACM Workshop on Artificial
Intelligence and Security (AISec)}, 2023.

\bibitem{owasp} OWASP Foundation, ``OWASP Top 10 for Large Language Model
Applications,'' 2025.

\bibitem{oecdgrok} OECD.AI Policy Observatory, ``AI prompt injection exploit drains
Grok-linked crypto wallet,'' AI Incidents and Hazards Monitor, May 2026.
\url{https://oecd.ai/en/incidents/2026-05-04-4a73}

\bibitem{giskard} Giskard, ``How Grok got prompt-injected: An X user drained
\$150,000 from an AI wallet,'' giskard.ai, 2026.

\bibitem{patlan} A.~S.~Patlan, P.~Sheng, S.~A.~Hebbar, P.~Mittal and P.~Viswanath,
``Real AI agents with fake memories: Fatal context manipulation attacks on Web3
agents,'' arXiv:2503.16248, 2025.

\bibitem{imf} T.~Davidovic and H.~Tourpe, ``How agentic AI will reshape payments,''
IMF Note 2026/004, International Monetary Fund, 2026.

\bibitem{freysa} The Block, ``Human player outwits Freysa AI agent in \$47,000 crypto
challenge,'' Nov.\ 2024.

\bibitem{aixbt} AI Incident Database, ``Incident 1003'' (AIXBT autonomous agent), 2025.
\url{https://incidentdatabase.ai/cite/1003}

\bibitem{lobstar} The Block, ``AI agent created by OpenAI dev `accidentally' sends
entire memecoin holdings to reply guy,'' Feb.\ 2026.

\bibitem{erc4337} V.~Buterin \emph{et al.}, ``ERC-4337: Account abstraction using alt
mempool,'' Ethereum Improvement Proposals, 2021.

\bibitem{erc8004} ``ERC-8004: Trustless agents,'' Ethereum Improvement Proposals, 2025.

\bibitem{willison} S.~Willison, ``The dual LLM pattern for building AI assistants that
can resist prompt injection,'' simonwillison.net, Apr.\ 2023.

\bibitem{camel} E.~Debenedetti \emph{et al.}, ``Defeating prompt injections by design,''
arXiv:2503.18813, 2025.

\bibitem{stix} OASIS, ``STIX Version 2.1,'' OASIS Standard, 2021.

\bibitem{mcmahan} B.~McMahan, E.~Moore, D.~Ramage, S.~Hampson and B.~A.~y~Arcas,
``Communication-efficient learning of deep networks from decentralized data,'' in
\emph{Proc.\ AISTATS}, 2017.

\bibitem{yang} Q.~Yang, Y.~Liu, T.~Chen and Y.~Tong, ``Federated machine learning:
Concept and applications,'' \emph{ACM Trans.\ Intelligent Systems and Technology},
vol.~10, no.~2, 2019.

\bibitem{bonawitz} K.~Bonawitz \emph{et al.}, ``Practical secure aggregation for
privacy-preserving machine learning,'' in \emph{Proc.\ ACM CCS}, 2017.

\bibitem{fate} Y.~Liu, T.~Fan, T.~Chen, Q.~Xu and Q.~Yang, ``FATE: An industrial grade
platform for collaborative learning with data protection,'' \emph{J.\ Machine Learning
Research}, vol.~22, no.~226, 2021.

\bibitem{flower} D.~J.~Beutel \emph{et al.}, ``Flower: A friendly federated learning
research framework,'' arXiv:2007.14390, 2020.

\bibitem{krum} P.~Blanchard, E.~M.~El~Mhamdi, R.~Guerraoui and J.~Stainer, ``Machine
learning with adversaries: Byzantine tolerant gradient descent,'' in \emph{Proc.\
NeurIPS}, 2017.

\bibitem{insured} Hu and Chen, ``Insured agents: A decentralized trust insurance
mechanism for agentic economy,'' arXiv:2512.08737, 2025.

\bibitem{fisco} FISCO BCOS Open Source Community, ``FISCO BCOS: An enterprise-level
financial blockchain platform,'' open-source software, github.com/FISCO-BCOS.

\bibitem{masabs} Monetary Authority of Singapore and Association of Banks in Singapore,
``MAS and ABS announce measures to bolster the security of digital banking,'' 19 Jan.\
2022.

\bibitem{ukdelay} HM Treasury, ``Payment Services (Amendment) Regulations 2024,'' draft
legislation published Mar.\ 2024.

\bibitem{commitment} G.~Bryan, D.~Karlan and S.~Nelson, ``Commitment devices,''
\emph{Annual Review of Economics}, vol.~2, 2010.

\end{thebibliography}

\end{document}
"
%                       (or uncomment the next line)
% \def\dohighlight{1}
\documentclass[10pt,a4paper,twocolumn]{article}

\usepackage[T1]{fontenc}
\usepackage[utf8]{inputenc}
\usepackage{mathptmx}
\usepackage[scaled=0.9]{helvet}
\usepackage[a4paper,top=19mm,bottom=22mm,left=16mm,right=16mm,columnsep=7mm]{geometry}
\usepackage{amsmath,amssymb}
\usepackage{microtype}
\usepackage{booktabs,tabularx,array}
\usepackage{stfloats}
\usepackage[font=small,labelfont=bf,labelsep=period]{caption}
\usepackage{titlesec}
\usepackage{enumitem}
\usepackage{algorithm}
\usepackage{algpseudocode}
\usepackage{xcolor}
\usepackage{tikz}
\usetikzlibrary{positioning,arrows.meta,fit,calc}
\usepackage[hidelinks]{hyperref}
\usepackage{url}
\usepackage{balance}

% ---- Yellow highlighter that breaks across lines (no extra package needed) ----
\definecolor{hlyellow}{HTML}{FFF14D}
\makeatletter
\newcommand{\uniq@box}[1]{\begingroup\setlength{\fboxsep}{0pt}%
  \colorbox{hlyellow}{\rule[-0.25em]{0pt}{1.05em}#1}\endgroup}
\newcommand{\uniq@gap}{{\color{hlyellow}\leaders\hrule height 0.8em depth 0.25em
  \hskip \fontdimen2\font plus \fontdimen3\font minus \fontdimen4\font}}
\def\uniq@scan#1 #2\uniq@end{%
  \if\relax\detokenize{#1}\relax\else\uniq@box{#1}\fi
  \if\relax\detokenize{#2}\relax\expandafter\@gobble\else\expandafter\@firstofone\fi
  {\uniq@gap\uniq@scan#2\uniq@end}}
\ifdefined\dohighlight
  \newcommand{\uniq}[1]{\leavevmode\uniq@scan#1 \uniq@end}
\else
  \newcommand{\uniq}[1]{#1}
\fi
\makeatother

% Headings
\titleformat{\section}{\normalfont\large\bfseries}{\thesection.}{0.5em}{}
\titleformat{\subsection}{\normalfont\normalsize\bfseries\itshape}{\thesubsection}{0.5em}{}
\titlespacing*{\section}{0pt}{9pt plus 2pt minus 1pt}{4pt}
\titlespacing*{\subsection}{0pt}{7pt plus 2pt minus 1pt}{3pt}
\setlist{nosep,leftmargin=*,topsep=2pt}

% Table helpers: ragged-right wrapping columns so text never overflows a cell
\newcolumntype{L}{>{\raggedright\arraybackslash}X}
\newcolumntype{P}[1]{>{\raggedright\arraybackslash}p{#1}}
\renewcommand{\arraystretch}{1.25}
\setlength{\tabcolsep}{4pt}
\newcommand{\tblfont}{\footnotesize}

\newcommand{\todo}[1]{\textcolor{red!70!black}{[#1]}}

\emergencystretch=2em
\sloppy
\raggedbottom

\title{\vspace{-6mm}\textbf{AgentVax: A Self-Immunising Wallet Agent\\
with Adaptive Financial Autonomy}}
\author{%
\todo{Author 1}, \todo{Author 2}, \todo{Author 3}\\[2pt]
\small Team \todo{team name}, \todo{university, country}\\
\small 2026 Shenzhen International FinTech Competition (FinTechathon), International Track, Topic F: Agent Wallet\\
\small Contact: \todo{email}}
\date{}

\begin{document}
\maketitle

\begin{abstract}
\noindent
AI agents that hold wallets can be talked into spending money they should not spend, and
several have already lost real funds. Existing agent wallets ask what the agent is
allowed to do and enforce it with caps and allowlists.
\uniq{AgentVax asks a second question: how much of its authorised power should the agent
exercise on its own right now?}
The owner states a mission in plain language; it becomes a signed mandate on chain that
the model cannot change. Before every payment the wallet contract checks
\uniq{three things: is it allowed, can the owner afford it, and is it safe right now,}
and
\uniq{a daily caution allowance limits how many borderline payments may pass without the
owner being asked.}
\uniq{When one wallet is attacked it publishes an immunity capsule, a structured alert
with no private data that describes the attack by its features, so other wallets
recognise reworded versions and become more cautious at once.}
\uniq{A detector trained jointly with federated learning then gives lasting protection
without any wallet sharing its data.}
\uniq{Learning and shared intelligence can only reduce autonomy. Only the owner can
extend authority, and only after a delay chosen in advance, so that an owner under
pressure cannot be rushed into a mistake.}
The paper specifies the architecture, the algorithms, the security design and a
reproducible evaluation protocol whose headline measure is the share of legitimate tasks
completed without the owner, with zero policy violations. It is a design and protocol
document; measured results are reported in the accompanying execution evidence.
\end{abstract}

\vspace{2pt}
\noindent\textbf{\textit{Keywords---}}agent wallet, adaptive autonomy, prompt injection,
payment fraud, threat sharing, federated learning.

\ifdefined\dohighlight
\smallskip\noindent{\footnotesize\textbf{Reading note.} \uniq{Yellow highlight} marks what
is specific to AgentVax. Unhighlighted features, such as spending caps, allowlists and
session keys, are common to agent wallets.\par}
\fi

% ---------------------------------------------------------------------------
\section{Introduction}

The theme of the International Track is ``AI as a Financial Participant'': an agent that
completes real tasks in a financial scenario~\cite{rules}. Topic~F asks for a wallet agent
that autonomously performs asset management, transaction execution and risk control on
chain. Entries are scored on task completion (40\%), security and compliance (30\%) and
innovation and interaction (30\%), with a 0--5 point bonus for on-chain
exploration~\cite{rules}.

A wallet agent differs from a chatbot in one respect that matters for security: its output
moves value. A language model reads text from many sources and cannot reliably tell data
from instructions~\cite{greshake,owasp}. Table~\ref{tab:incidents} lists four public
incidents in which an AI agent with a wallet sent funds it should not have sent. In the
most recent, an attacker posted an instruction in Morse code; one model translated it and
a wallet bot executed the result~\cite{oecdgrok,giskard}. In the laboratory, Patlan et
al.\ showed that planting false entries in a wallet agent's memory is more effective
still~\cite{patlan}. Regulators have named the same risk for agent-initiated
payments~\cite{imf}. An agent that pays invoices also inherits the oldest payment fraud
there is: a message that appears to come from a known supplier and asks for payment to a
new account. That fraud usually works by rushing a person, not by fooling a machine.

\begin{table*}[t]
\caption{Reported incidents involving AI agents that control wallets.}
\label{tab:incidents}
\tblfont
\begin{tabularx}{\textwidth}{@{}P{0.09\textwidth} P{0.13\textwidth} L P{0.17\textwidth} P{0.19\textwidth}@{}}
\toprule
\textbf{Date} & \textbf{Agent} & \textbf{What happened} & \textbf{Reported loss} & \textbf{Kind of failure} \\
\midrule
Nov.\ 2024 & Freysa~\cite{freysa} & A player persuaded the agent that its transfer-approval function handled incoming funds; it released the prize pool. This was a designed game. & About \$47,000 & Persuasion of the model \\
Mar.\ 2025 & AIXBT~\cite{aixbt} & An attacker entered the control dashboard and queued two prompts instructing the agent to transfer funds. & 55.5 ETH, about \$106,000 & Injected instructions after a dashboard breach \\
Feb.\ 2026 & Lobstar Wilde~\cite{lobstar} & The agent intended to send a small gift and transferred its entire token holdings. & About \$250,000 in tokens & Agent error in the amount \\
May 2026 & Grok and Bankr~\cite{oecdgrok,giskard} & A token transfer raised the wallet's permissions; a Morse-coded post was translated into a transfer order and executed. & About \$150,000--\$174,000; most reportedly returned & Encoded prompt injection \\
\bottomrule
\end{tabularx}
\end{table*}

Two observations follow. First, the failures are of different kinds: persuasion,
injection, and an arithmetic mistake with no attacker at all. A wallet that only asks
whether a payment is permitted misses the third kind entirely. Second, each agent failed
alone. Nothing learned in one incident protected the next wallet, because the evidence is
conversation and transaction data that owners cannot hand to others.

\textbf{Contributions.} AgentVax combines:
\begin{enumerate}
\item a \emph{mission-to-mandate} step and a task engine that completes four wallet tasks
from plain-language instructions;
\item \uniq{a \emph{three-question authorisation rule}: the tier of a request is the
highest tier demanded by permission, affordability and current threat level,} enforced by
the contract and a guardian co-signature;
\item \uniq{a \emph{financial state engine} that makes affordability a condition of
autonomous spending and is computed on chain from registered obligations;}
\item \uniq{a \emph{caution allowance}: a daily budget of tolerated uncertainty that
shrinks at once under threat and recovers slowly;}
\item \uniq{\emph{immunity capsules}: feature-based, expiring, tighten-only alerts that
match reworded attacks and inform the owner at the moment of approval;}
\item \uniq{a \emph{federated detector} that turns each attack into lasting protection,
with every shared model anchored on chain;}
\item \uniq{\emph{owner-pressure safeguards}: time-locked expansion of authority, holds
triggered by pressure, and considered confirmation;}
\item an \emph{evaluation protocol} and attack suite (VaxLab) whose headline measure is
the safe autonomy rate.
\end{enumerate}

\textbf{Scope.} The prototype uses three wallets on a public EVM test network with
valueless test tokens. Items left for later are listed in Section~\ref{sec:future}.

% ---------------------------------------------------------------------------
\section{Background and Related Work}

\textbf{Agent wallets.} Account abstraction (ERC-4337) lets a smart-contract account
define its own validation logic, including limited session keys~\cite{erc4337}. Agent
identity and reputation registries have been proposed in ERC-8004~\cite{erc8004}.
Commercial agent wallets add caps, daily limits and allowlists. These bound the loss
per wallet, but they treat autonomy as fixed, ignore the owner's financial position and
provide no shared defence.

\textbf{Prompt injection.} Indirect injection places instructions in content the agent
retrieves~\cite{greshake}; it heads the OWASP list of risks for language-model
applications~\cite{owasp}. Architectural defences separate a privileged model that plans
from a quarantined model that reads untrusted content~\cite{willison}, or restrict what
data may flow into actions~\cite{camel}. AgentVax adopts this separation.

\textbf{Threat sharing.} Security teams exchange structured indicators of attack in
formats such as STIX~\cite{stix}, and banks exchange fraud alerts. The format of a
capsule is therefore not new. \uniq{What is new is the effect: a capsule is applied
automatically by a wallet contract and changes what an agent may spend without approval.}

\textbf{Federated learning.} Clients train locally and a coordinator averages their
updates~\cite{mcmahan,yang}, with secure aggregation~\cite{bonawitz} and robust
aggregation~\cite{krum}; FATE~\cite{fate} and Flower~\cite{flower} implement these.

Table~\ref{tab:compare} summarises how AgentVax differs from a typical agent wallet.

\begin{table*}[t]
\caption{A typical agent wallet compared with AgentVax.}
\label{tab:compare}
\tblfont
\begin{tabularx}{\textwidth}{@{}P{0.19\textwidth} L L@{}}
\toprule
\textbf{Aspect} & \textbf{Typical agent wallet} & \textbf{AgentVax} \\
\midrule
Central question & What is the agent allowed to do? & \uniq{How much of its authorised power should it exercise on its own right now?} \\
How the owner sets it up & Individual limits and lists & \uniq{A mission in plain language, converted to a mandate the owner signs and the model cannot alter} \\
Checks before paying & Is this allowed? & \uniq{Allowed, affordable, and safe right now; the strictest answer wins} \\
Knowledge of the owner's finances & None; a cap is a fixed number & \uniq{Reserve and upcoming bills are registered on chain; the spendable amount is recomputed for every request} \\
Many slightly unusual small payments & Each passes on its own & \uniq{They draw on a daily caution allowance; when it runs out the rest wait for the owner} \\
Who decides the risk level & The model, or a filter in the same process & \uniq{A separate guardian whose signature the contract requires and the model cannot produce} \\
When another wallet is attacked & Nothing changes & \uniq{A capsule raises the tier of matching requests, including reworded ones, and is shown to the owner when approval is asked} \\
What wallets share & Nothing & \uniq{Fixed-field alerts and masked model updates; never messages, prompts or transactions} \\
Effect of a good track record & May relax scrutiny & \uniq{Lowers risk estimates only; it never creates a permission} \\
The owner under pressure & The owner's approval is final and immediate & \uniq{Approvals of pressured requests are held, and the owner can cancel at once} \\
Who may loosen a rule & Often the agent's own key & \uniq{Only the owner, and only after a delay fixed in advance; tightening is instant} \\
\bottomrule
\end{tabularx}
\end{table*}

% ---------------------------------------------------------------------------
\section{Requirements and Design Goals}

Table~\ref{tab:scoring} maps each scoring criterion to the component that addresses it and
to the evidence submitted. Three principles govern the design:

\begin{description}[leftmargin=0pt,itemsep=1pt]
\item[P1.] \uniq{The model proposes; deterministic code authorises; the wallet contract
enforces.}
\item[P2.] \uniq{Learning changes risk estimates; only the owner changes authority.}
\item[P3.] \uniq{Caution is immediate and expansion is slow, for the model, the network
and the owner alike.}
\end{description}

\begin{table*}[t]
\caption{Alignment of AgentVax with the official scoring criteria~\cite{rules}.}
\label{tab:scoring}
\tblfont
\begin{tabularx}{\textwidth}{@{}P{0.16\textwidth} P{0.07\textwidth} L L@{}}
\toprule
\textbf{Criterion} & \textbf{Weight} & \textbf{AgentVax component} & \textbf{Evidence submitted} \\
\midrule
Task completion & 40\% &
Mission-to-mandate; task engine for four wallet tasks; financial state engine for asset
management; clarification of ambiguous requests &
Transaction hashes for each task; safe autonomy rate; task-suite and affordability-suite
results; test cases in the repository \\
Security and compliance & 30\% &
Three-question tiers and caution allowance enforced in the contract; guardian
co-signature; planner and reader separation; tighten-only capsules; time-locked owner
actions and pressure holds; testnet-only operation &
Security self-assessment report; VaxLab attack results; known-risk list
(Table~\ref{tab:risks}) \\
Innovation and interaction & 30\% &
Adaptive autonomy; feature-based capsules; federated detector; protection of the owner
under pressure; an interface that explains each decision and requires considered
confirmation &
Capsule response on reworded attacks; before/after detection on an unseen attack family;
demo video \\
Challenge bonus & 0--5 pts &
Mandate, obligation registry, allowance, capsule bulletin and model anchoring on chain;
time-locked and co-signed execution paths &
Contract addresses and transaction hashes for each \\
\bottomrule
\end{tabularx}
\end{table*}

% ---------------------------------------------------------------------------
\section{System Architecture}

Figure~\ref{fig:arch} shows the components. An instruction from the owner reaches the
\emph{planner}, the only model allowed to propose transactions. External content is read
only by the \emph{quarantined reader}, which returns typed fields and flags and never
free text. The planner emits a structured \emph{intent}. The \emph{validator} checks it
deterministically. The \emph{guardian} scores the associated content with the injection
detector and signs the intent together with a risk band, its source and its flags. The
\emph{policy wallet} contract then decides: it checks the mandate, computes the
discretionary amount from its \emph{obligation registry}, reads active capsules from the
\emph{threat bulletin} contract, and charges the caution allowance. The detector is
trained jointly through the \emph{federated aggregator}.

\begin{figure*}[t]
\centering
\begin{tikzpicture}[
  font=\footnotesize,
  box/.style={draw, rounded corners=2pt, align=center, minimum height=9mm, text width=25mm, fill=white},
  chain/.style={box, fill=black!8},
  untrusted/.style={box, dashed},
  arr/.style={-{Stealth[length=2mm]}, semithick},
  lab/.style={font=\scriptsize, align=center}
]
\node[box] (owner) {Owner\\(mission, approvals)};
\node[box, right=8mm of owner] (planner) {Planner\\(privileged model)};
\node[box, right=8mm of planner] (validator) {Validator\\(deterministic)};
\node[box, right=8mm of validator] (guardian) {Guardian\\(detector + signer)};
\node[chain, right=8mm of guardian] (wallet) {Policy wallet\\(mandate, allowance)};
\node[untrusted, below=11mm of owner] (content) {Untrusted inputs\\(web, mail, memory)};
\node[box, below=11mm of planner] (reader) {Quarantined reader\\(fields and flags)};
\node[box, below=11mm of validator] (peers) {Peer wallets\\B, C};
\node[box, below=11mm of guardian] (fed) {Federated aggregator\\(masked updates)};
\node[chain, below=11mm of wallet] (bulletin) {Threat bulletin\\(capsules, hashes)};
\draw[arr] (owner) -- (planner);
\draw[arr] (planner) -- node[lab, above]{intent} (validator);
\draw[arr] (validator) -- (guardian);
\draw[arr] (guardian) -- node[lab, above]{signed} (wallet);
\draw[arr] (content) -- (reader);
\draw[arr] (reader) -- node[lab, right]{fields} (planner);
\draw[arr, <->] (guardian) -- node[lab, right]{model} (fed);
\draw[arr, <->] (peers) -- node[lab, above]{updates} (fed);
\draw[arr] (fed) -- node[lab, above]{hash} (bulletin);
\draw[arr, <->] (wallet) -- node[lab, right]{capsules} (bulletin);
\end{tikzpicture}
\caption{AgentVax architecture. Grey boxes are smart contracts; the dashed box is
untrusted input. Only the planner proposes transactions, and it never reads untrusted
text directly. The wallet contract decides from the mandate, its obligation registry, its
allowance and the capsules in the bulletin.}
\label{fig:arch}
\end{figure*}

% ---------------------------------------------------------------------------
\section{Mission, Mandate and Tasks}

\subsection{From mission to mandate}
The owner does not authorise each transaction. The owner states a mission once, for
example: ``Keep 600 in reserve. Pay my landlord and my approved suppliers on time. Pay
approved payees up to 500 without asking. Ask me about anyone new. Do not borrow, invest
or trade.'' The model converts this into a structured \emph{mandate}
(Table~\ref{tab:mandate}) and shows it back in plain language.
\uniq{The mandate takes effect only when the owner signs it, it is stored in the wallet
contract, and no model output can modify it.}
It carries an expiry date, after which the agent can read but not pay, and two
waiting times that protect the owner (Section~\ref{sec:pressure}).

\begin{table}[t]
\caption{Example mandate derived from a mission.}
\label{tab:mandate}
\tblfont
\begin{tabularx}{\columnwidth}{@{}L P{28mm}@{}}
\toprule
\textbf{Field} & \textbf{Value} \\
\midrule
Reserve $R$ & 600 \\
Autonomous limit per payment $c_1$ & 500 \\
Autonomous limit per day $d_1$ & 1,000 \\
Limit with delay $c_2$ & 2,000 \\
Approved payees (identity and address) & Landlord, hosting provider, software vendor \\
New payee or new address & Owner approval \\
Prohibited actions & Borrowing, investing, trading \\
Caution allowance per day $A$ & 10 points \\
Loosening delay $\Lambda$ & 24 hours \\
Pressure hold $\Pi$ & 60 minutes \\
Trusted co-signer & Optional \\
Expiry & 90 days \\
\bottomrule
\end{tabularx}
\end{table}

\subsection{Authority is not familiarity}
An approved payee is a pair: an identity and a destination address.
\uniq{A request to pay a known payee at a new address is a request to approve a new
payee, however many successful payments preceded it.}
History is used to estimate risk, for example the usual amount for a payee.
\uniq{It can lower or raise an estimate; it cannot add a payee, raise a limit or shorten
a delay.}

\subsection{Supported tasks}
Table~\ref{tab:tasks} lists the four tasks. They cover transaction execution (T1, T2),
asset management (T3, T4) and risk control (the decision applied to all four).

\begin{table}[t]
\caption{Wallet tasks supported by the prototype.}
\label{tab:tasks}
\tblfont
\begin{tabularx}{\columnwidth}{@{}P{5mm} P{19mm} L@{}}
\toprule
\textbf{ID} & \textbf{Task} & \textbf{Example instruction and on-chain effect} \\
\midrule
T1 & Payment & ``Pay the hosting invoice.'' Transfers test stablecoin to an approved payee. \\
T2 & Recurring payment & ``Pay rent on the 1st of each month.'' Creates a schedule, which also registers the next payment as an obligation. \\
T3 & Idle-balance sweep & ``Save what I do not need.'' Moves only the discretionary amount to a savings vault. \\
T4 & Fee-reserve upkeep & ``Never let fees run out.'' Tops up the native-token reserve when it falls below a floor. \\
\bottomrule
\end{tabularx}
\end{table}

\subsection{Typed intents and explanation}
The planner does not produce transactions. It produces an intent
\begin{align*}
I=(&\textit{task},\ \textit{payee\_id},\ \textit{amount},\ \textit{asset},\\
   &\textit{schedule},\ \textit{purpose},\ \textit{nonce}),
\end{align*}
serialised against a fixed schema. \textit{payee\_id} is an index into the on-chain
allowlist; the planner cannot supply a raw address. If a required field cannot be filled,
the agent asks one question instead of guessing. After each action the interface shows
the task, the tier, the transaction hash and \uniq{which check determined the outcome.}
After a delay or refusal it states the reason in the owner's terms, for example the size
of a shortfall, and the available choices.

% ---------------------------------------------------------------------------
\section{Financial State Engine}
\label{sec:state}

Caps answer whether a payment is permitted. They do not answer whether it is wise today.
The state engine supplies that answer from three on-chain quantities: the stablecoin
balance $B$; the reserve $R$; and a registry of obligations $O=\{(a_j,d_j)\}$, each an
amount and a due date. Scheduled payments (T2) enter the registry automatically; one-off
bills are registered when the owner or the agent records them.

\subsection{Discretionary amount}
For a request made at time $t$ with planning horizon $H$, the discretionary amount is
\[
D \;=\; B \;-\; R \;-\!\!\sum_{j\in O,\; d_j\le t+H}\!\! a_j ,
\]
where an obligation is left out of the sum when the request is the payment of that
obligation. A request of amount $a$ is classified as
\[
\textit{afford}(a)=
\begin{cases}
\text{comfortable} & a\le D,\\
\text{tight} & D< a\le B-R,\\
\text{breach} & a> B-R.
\end{cases}
\]
A comfortable request does not raise the tier. A tight request is queued, and a request
that would breach the reserve needs the owner.
\uniq{Expected income is shown to the owner but never counted, so the agent cannot spend
money that has not arrived.}

\subsection{Where it is enforced}
\uniq{The wallet contract computes the discretionary amount itself at execution time, so
affordability does not depend on the model's arithmetic.}
The registry is bounded in length to keep the cost predictable.
\uniq{The agent may add an obligation, which lowers what it can spend; only the owner may
remove one or lower the reserve.}

\subsection{Worked example}
Table~\ref{tab:state} shows a wallet holding 1,850 under the mandate of
Table~\ref{tab:mandate}. With three bills due, $D=180$. The agent is asked to upgrade a
hosting plan for 300. The payee is approved, the amount is under the limit and there is
no sign of attack, so a typical agent wallet would pay. Here $D<300\le B-R=1{,}250$, so
the request is tight. The agent replies that the payment is allowed but would leave the
wallet 120 short of its bills unless the expected 500 arrives, queues it, and offers two
choices: wait for the income, or approve now.

\begin{table}[t]
\caption{Worked example of the financial state.}
\label{tab:state}
\tblfont
\begin{tabularx}{\columnwidth}{@{}L r@{}}
\toprule
\textbf{Item} & \textbf{Amount} \\
\midrule
Balance $B$ & 1,850 \\
Reserve $R$ & $-600$ \\
Rent, due in 5 days & $-800$ \\
Hosting invoice & $-150$ \\
Software bill & $-120$ \\
\midrule
Discretionary amount $D$ & 180 \\
Expected income (shown, not counted) & 500 \\
\bottomrule
\end{tabularx}
\end{table}

The same check would have stopped the third incident in Table~\ref{tab:incidents}, where
no attacker was involved: an amount far above the discretionary amount cannot execute
without the owner.

% ---------------------------------------------------------------------------
\section{Security Design}

\subsection{Threat model}
The adversary controls content the agent may read: web pages, messages, invoices, token
metadata, and entries that reach the agent's memory. The adversary may operate one
participating wallet and publish capsules or model updates from it, and may use AI of its
own to vary and adapt attacks. The adversary may also contact the owner directly and
press for an approval or a change of rule. The adversary does not hold the owner key or the guardian
key and cannot alter deployed contracts.
\uniq{No defence assumes that the agent's model out-reasons the attacker: a fully
manipulated model must still be unable to move funds beyond the mandate.}

\subsection{Three-question tiers}
Table~\ref{tab:tiers} gives the rule. Each question proposes a tier and
\uniq{the request receives the highest tier any of the three questions demands}
(Algorithm~\ref{alg:tier}). The agent holds a session key valid only for calls that
satisfy the rule; the owner key can revoke it or freeze the wallet at any time.

\begin{table*}[t]
\caption{Permission tiers. A request receives the highest tier demanded by any of the
three questions.}
\label{tab:tiers}
\tblfont
\begin{tabularx}{\textwidth}{@{}P{0.085\textwidth} L P{0.17\textwidth} L P{0.19\textwidth}@{}}
\toprule
\textbf{Tier} & \textbf{Allowed?} & \textbf{Affordable?} & \textbf{Safe right now?} & \textbf{Authorisation required} \\
\midrule
L0 Read & No value moves & Not applicable & Not applicable & None \\
L1 Routine & Approved payee; amount $\le c_1$; daily total $\le d_1$ & Comfortable: $a\le D$ & Detector band low; no pressure flag; no capsule fully matches & Session key, guardian signature, and enough allowance \\
L2 Elevated & Approved payee; amount $\le c_2$ & Tight: $D<a\le B-R$ & Detector band medium, a pressure flag, or a single-observer capsule matches & Queued for delay $\Delta$; owner may cancel \\
L3 Sensitive & New payee or address; amount $>c_2$; prohibited action; any loosening of the mandate & Breach: $a>B-R$ & Detector band high, or a confirmed capsule matches & Owner co-signature; held for $\Pi$ under pressure; rule changes wait $\Lambda$ \\
\bottomrule
\end{tabularx}
\end{table*}

The detector runs off chain, so the contract cannot observe its output.
\uniq{The guardian therefore signs the intent hash together with the risk band, the
source of the request and its flags, and the contract requires a valid guardian signature
with band low before it will execute at L1.}
The planner has no access to the guardian key, so a manipulated planner cannot promote
its own request.

\subsection{Caution allowance}
The three questions judge one request at a time. A patient attacker can send many
requests that each look almost normal. \uniq{The caution allowance bounds how much
uncertainty the wallet accepts in a day without the owner.}

Each L1 request is charged points for the flags it carries (Table~\ref{tab:allowance}).
The flags are checkable facts, not model opinions: the contract computes those it can see
and takes the rest from the guardian's signed attestation. If the remaining allowance
covers the charge, the request executes and the allowance is reduced; otherwise the
request is queued at L2. The point values are owner-set integers; the table shows the
defaults we chose, which are not learned.

\begin{table}[t]
\caption{Default charges against the caution allowance.}
\label{tab:allowance}
\tblfont
\begin{tabularx}{\columnwidth}{@{}L r@{}}
\toprule
\textbf{Flag on an otherwise routine request} & \textbf{Points} \\
\midrule
None: a payee paid before, usual amount, owner instruction & 0 \\
Request arose from external content, not a direct owner instruction & 2 \\
Request uses more than half of the discretionary amount & 2 \\
Amount exceeds the largest amount previously paid to this payee & 3 \\
First payment to this payee since approval & 3 \\
Request partially matches an active capsule & 5 \\
\bottomrule
\end{tabularx}
\end{table}

\uniq{The allowance itself moves with circumstances. It is halved the moment a capsule
that applies to the wallet is published, and set to zero when the wallet's own owner
confirms an attack.} It then refills at a fixed rate $\rho$ up to the mandate value $A$:
\[
\alpha(t)=\min\!\big(A,\ \alpha(t_0)+\rho\,(t-t_0)\big).
\]
With the defaults, an attacker who sends a series of invoices from an approved payee,
each slightly above the previous maximum, is charged five points per invoice; the third
one waits for the owner.

\begin{algorithm}[t]
\caption{Authorisation of intent $I$, amount $a$}
\label{alg:tier}
\footnotesize
\begin{algorithmic}[1]
\If{schema, provenance or state check fails} \Return reject \EndIf
\State $q_1 \gets$ tier from the mandate \Comment{allowed?}
\State $q_2 \gets$ tier from $\textit{afford}(a)$ \Comment{affordable?}
\State $q_3 \gets$ tier from guardian band; at least L2 if a pressure flag is set
\For{each unexpired capsule $c$ that fully matches $I$}
  \State $q_3 \gets \max(q_3,\ \textit{restriction}(c))$ \Comment{safe now?}
\EndFor
\State $\textit{tier} \gets \max(q_1,q_2,q_3)$
\If{$\textit{tier}=\text{L1}$}
  \State $\kappa \gets$ points charged for the flags of $I$
  \If{$\kappa \le \alpha$} $\alpha \gets \alpha-\kappa$
  \Else\ $\textit{tier} \gets \text{L2}$ \Comment{allowance spent}
  \EndIf
\EndIf
\State execute, queue or request co-signature by \textit{tier}
\State hold an owner approval for $\Pi$ if a pressure flag or capsule is present
\State report \textit{tier}, the deciding check and any matching capsule
\end{algorithmic}
\end{algorithm}

\subsection{The owner under pressure}
\label{sec:pressure}
The top tier relies on the owner, and the owner can be rushed, frightened or deceived.
Most payment fraud succeeds this way: the victim authorises it. Banks already answer
with time. Banks in Singapore delay the activation of a new security token by at least
twelve hours and apply cooling-off periods to key account changes~\cite{masabs}, and
United Kingdom rules allow a payment to be held for up to four business days when fraud
is suspected~\cite{ukdelay}.
\uniq{AgentVax builds this protection into the wallet as a commitment the owner makes
while calm and cannot waive in the moment}~\cite{commitment}.

\emph{Loosening is slow.}
\uniq{Any owner action that expands authority takes effect only after the loosening
delay $\Lambda$ fixed in the mandate: adding a payee or an address, raising a limit or
the allowance, lowering the reserve, or shortening a delay, including $\Lambda$ itself.}
\uniq{Cancelling a pending change, lowering a limit and freezing the wallet are
immediate.} The contract enforces this as a time-lock on owner calls. The same rule
limits the damage of a stolen owner key: the thief must wait, and the real owner can
freeze the wallet.

\emph{Pressure is a signal.} Urgent wording, a request for secrecy and an instruction to
ignore earlier details are \emph{pressure flags}, raised by a fixed phrase list and by
the reader.
\uniq{A request that carries a pressure flag cannot execute at L1, and an owner approval
given while a pressure flag or a matching capsule is present takes effect only after the
pressure hold $\Pi$.}
The agent states the reason: the message presses for immediate payment, which is a
common sign of fraud. The hold has a fixed length. An attacker who removes the pressure
to avoid the hold gives up the main lever of the fraud.

\emph{Considered confirmation.} For an approval that carries flags, the interface shows
the payee's current and requested address side by side, the payment history and any
matching capsule, and asks the owner to type the payee's name and to state who asked for
the payment.
\uniq{The answers can lengthen a hold; nothing the owner types in the moment can shorten
one.}
To limit fatigue, these steps appear only for flagged L3 requests.

\emph{Trusted co-signer.} Genuine emergencies exist. The mandate may name a second
person whose signature, together with the owner's, releases a hold or a pending change
early. The co-signer cannot act alone.

\subsection{Hallucination protection}
The validator applies three deterministic checks before the guardian sees an intent:
(i)~\emph{schema}: all fields present and well typed;
(ii)~\emph{provenance}: \textit{payee\_id} exists in the on-chain allowlist, and
\textit{amount} and \textit{asset} equal a value in the owner's instruction, in the
mandate, or in a typed field returned by the reader for a document the owner referenced;
(iii)~\emph{state}: the nonce is fresh and the balance covers the amount and the fee.
Balances, obligations and transaction status shown in the interface are read from the
chain, not generated by the model.

\subsection{Injection defence}
Defence is layered so that no single component is trusted:
\begin{enumerate}
\item \emph{Separation.} The reader has no tools and returns only enumerations, numbers,
allowlist indices and flags. Free text from untrusted sources never enters the planner's
context~\cite{willison,camel}.
\item \emph{Bounded memory.} The planner sees only what the current decision needs.
Long-term memory holds structured facts (per-payee address, usual amount, last
verification) and threat records; raw prompts and conversations are not retained.
\item \emph{Detection.} The guardian scores the raw untrusted content associated with a
request and assigns a risk band (Section~\ref{sec:fed}).
\item \emph{Containment.} Whatever the model outputs, the contract bounds the result. An
undetected attack can at most cause L1 payments to approved payees, within the limits,
the discretionary amount and the day's allowance.
\item \emph{Shared alert.} Capsules narrow that remaining window while a threat is active
(Section~\ref{sec:capsule}).
\end{enumerate}

% ---------------------------------------------------------------------------
\section{Immunity Capsules}
\label{sec:capsule}

Retraining a detector takes time and needs examples. A wallet that has just been attacked
can warn the others at once.
\uniq{An immunity capsule is a small record of fixed fields that says what kind of attack
was seen, which features it had, which action it targeted, and how much autonomy to
withdraw and for how long.}

\subsection{Contents}
Table~\ref{tab:capsule} lists the fields.
\uniq{Every field is an enumeration, a set of flags, a number or a signature. A capsule
carries no message text, no prompt and no transaction, so it cannot leak private data and
cannot itself carry an injected instruction.}

\begin{table}[t]
\caption{Fields of an immunity capsule.}
\label{tab:capsule}
\tblfont
\begin{tabularx}{\columnwidth}{@{}P{22mm} L@{}}
\toprule
\textbf{Field} & \textbf{Values} \\
\midrule
Attack family & Direct, indirect, memory, encoded, recipient substitution \\
Channel & Owner chat, web page, message, memory, token metadata \\
Features & Any of: known payee identity, new destination, urgency, secrecy request, override attempt, encoded content, amount anomaly, memory-sourced \\
Affected action & Payment, schedule creation, sweep, payee or mandate change \\
Restriction & Delay (L2) or owner approval (L3) \\
Confidence & Single observer, or confirmed by $q$ wallets \\
Expiry & Time after which the capsule has no effect \\
Guard version & Detector version in force when observed \\
Attestations & Signatures of the observing wallets \\
\bottomrule
\end{tabularx}
\end{table}

\subsection{Matching reworded attacks}
\uniq{A capsule describes an attack by its features, not its wording.} Let $F_c$ be the
capsule's feature set and $F_I$ the flags of a request. The capsule \emph{fully matches}
when its channel is in the wallet's on-chain profile, its affected action equals the
request's action and $F_c\subseteq F_I$. It \emph{partially matches} when at least half
of $F_c$ is present. A full match raises the tier; a partial match is charged to the
allowance. \uniq{An attacker who rewrites the message but keeps the method produces the
same flags and is matched.} Three of the flags (new destination, amount anomaly,
memory-sourced) are computed by rule; the others are extracted by the reader and signed
by the guardian.

\subsection{Effect}
\uniq{The wallet contract reads unexpired capsules from the bulletin when it executes a
request, so a capsule takes effect on the very next transaction and does not rely on the
agent choosing to obey it.}
A wallet the threat does not apply to is unaffected and keeps its normal autonomy.

Some attacks are already stopped by the mandate: a payment to a new address needs the
owner whether or not a capsule exists. There the danger is the approval itself, because
this kind of fraud succeeds by persuading a person.
\uniq{When the owner is asked to approve, the interface shows any matching capsule, and
the approval is subject to the pressure hold of Section~\ref{sec:pressure}.}

\subsection{Limits on capsules}
A channel that can change other wallets' behaviour is itself a target, so capsules are
constrained:
\begin{itemize}
\item \uniq{\emph{Tighten only.} A capsule can raise a tier or charge the allowance. No
capsule can lower a tier, add a payee or raise a limit.}
\item \emph{Bounded by confidence.} A capsule from one observer can impose at most a
delay. Owner approval can be imposed only when $q$ distinct wallets have attested the
same family and features ($q=2$ in the prototype).
\item \emph{Expiry.} Every capsule expires; the maximum lifetime is fixed in the bulletin
contract. Autonomy returns without anyone acting.
\item \emph{Rate limit and identity.} Each registered wallet may publish a limited number
of capsules per day, and every capsule names its issuer.
\item \emph{Owner override.} An owner may dismiss a capsule for their own wallet.
\end{itemize}

\subsection{Alert first, vaccine later}
\uniq{The capsule acts within one transaction but is blunt and temporary. The retrained
detector arrives later but recognises variants beyond the listed features and stays.}
A capsule's lifetime is set to outlast the next training round, so protection is
continuous. Algorithm~\ref{alg:round} shows the full response.

% ---------------------------------------------------------------------------
\section{Federated Detector}
\label{sec:fed}

\subsection{Detector}
Each wallet $k$ keeps a private set $D_k=\{(x_i,y_i)\}$ of content snippets labelled
benign ($y=0$) or attack ($y=1$). Labels come from blocked requests the owner confirmed
as attacks and from requests the owner approved. A frozen text encoder $\phi$ maps a
snippet to a vector; a small classifier with weights $w$ outputs the score
\[
s(x)=\sigma\!\big(w^{\top}\phi(x)+b\big)\in[0,1].
\]
The risk band is low if $s<\tau_{\mathrm{low}}$, high if $s\ge\tau_{\mathrm{high}}$, and
medium otherwise. Only $(w,b)$ is trained and shared.

\subsection{Training round}
Each wallet minimises the local loss
\[
F_k(w)=\frac{1}{n_k}\sum_{(x,y)\in D_k}\ell\big(s(x),y\big),\qquad n_k=|D_k|,
\]
for $E$ epochs from the current global weights, and the aggregator computes the weighted
average~\cite{mcmahan}
\[
w^{t+1}=\sum_{k=1}^{K}\frac{n_k}{n}\,w_k^{t+1},\qquad n=\sum_{k=1}^{K}n_k .
\]
Updates are clipped to norm $C$ and masked pairwise so that the aggregator sees only the
sum~\cite{bonawitz}. No snippet, conversation or transaction leaves a wallet.

\begin{algorithm}[t]
\caption{Immunity response to an attack on wallet $a$}
\label{alg:round}
\footnotesize
\begin{algorithmic}[1]
\Require global weights $w^{t}$, wallets $1..K$, clip norm $C$
\State owner of $a$ confirms the attack; $\alpha_a \gets 0$
\State $a$ adds it to $D_a$ and publishes a capsule \Comment{alert}
\State wallets it applies to: $\alpha \gets \alpha/2$; matching requests restricted
\State a second observer may attest, raising it to confirmed
\For{each wallet $k$ in parallel} \Comment{vaccine}
  \State $w_k \gets$ \textsc{LocalTrain}$(w^{t}, D_k, E)$
  \State $\delta_k \gets$ \textsc{Clip}$(w_k - w^{t}, C)$
  \State send \textsc{Mask}$(n_k\,\delta_k)$ to the aggregator
\EndFor
\State $w^{t+1} \gets w^{t} + \frac{1}{n}\sum_k n_k\,\delta_k$ \Comment{masks cancel}
\State aggregator posts $H(w^{t+1})$ and $t{+}1$ to the bulletin
\For{each wallet $k$}
  \If{model hash matches bulletin \textbf{and} canary passes}
    \State adopt $w^{t+1}$
  \Else
    \State keep $w^{t}$ and raise an alert
  \EndIf
\EndFor
\State the capsule expires; allowances refill; the detector remains
\end{algorithmic}
\end{algorithm}

\subsection{Anchoring and adoption safeguards}
The bulletin stores the hash and version of each shared model.
\uniq{A wallet can therefore verify that it received the model the federation agreed on,
and an auditor can reconstruct which detector and which capsules governed any past
transaction.}
Before adopting a model, each wallet evaluates it on a private \emph{canary set}; a model
that lowers recall or raises the false-positive rate beyond a tolerance is rejected.
Together with clipping, this limits what a poisoning participant can do~\cite{krum}.

% ---------------------------------------------------------------------------
\section{Implementation}

Table~\ref{tab:stack} summarises the prototype. The federated layer is written against
Flower for speed of development; it uses only horizontal federated averaging with secure
aggregation, which FATE also provides, so it can be moved to FATE without redesign.
Contracts are plain Solidity with no chain-specific features, which keeps a later port to
a permissioned EVM-compatible chain such as FISCO BCOS~\cite{fisco} straightforward.

\begin{table}[t]
\caption{Prototype implementation.}
\label{tab:stack}
\tblfont
\begin{tabularx}{\columnwidth}{@{}P{23mm} L@{}}
\toprule
\textbf{Component} & \textbf{Implementation} \\
\midrule
Chain & Public EVM test network; test stablecoin with no monetary value \\
Policy wallet & Solidity smart account: signed mandate, session key, allowlist, limits, obligation registry, caution allowance, time-lock queue for payments and for owner changes, pressure hold, co-signature, freeze \\
Savings vault & Minimal deposit and withdraw contract, owner-only withdrawal \\
Threat bulletin & Registry of capsules, attestations and model hashes; enforces capsule limits \\
Planner, reader & Two separately prompted language-model roles; reader has no tools \\
Validator & Deterministic code; schema, provenance and state checks \\
Guardian & Encoder plus linear classifier; separate signing key \\
Federation & Three clients, federated averaging, pairwise masking \\
VaxLab & Scripted attackers plus one rewording agent; never holds a wallet key \\
Interface & Web chat with mandate view, financial state panel, allowance meter, alerts and considered confirmation at approval, explorer links \\
\bottomrule
\end{tabularx}
\end{table}

% ---------------------------------------------------------------------------
\section{Evaluation Protocol}
\label{sec:eval}

All measurements are produced by scripts in the repository so that reviewers can repeat
them. Table~\ref{tab:metrics} lists the metrics and the targets the design aims for.
Targets are design goals, not results. For every task, tier path, mandate entry, capsule
and model anchor, the submission lists the transaction hash and contract address.

\begin{table}[t]
\caption{Evaluation metrics and design targets.}
\label{tab:metrics}
\tblfont
\begin{tabularx}{\columnwidth}{@{}L P{30mm} P{16mm}@{}}
\toprule
\textbf{Metric} & \textbf{Definition} & \textbf{Target} \\
\midrule
Safe autonomy rate & Legitimate tasks done without the owner, zero violations & $\ge 80\%$ \\
Hard-policy violations & Mandate breaches of any kind & 0 \\
Task success rate & Correct state change on chain, 40 instructions & $\ge 90\%$ \\
Ambiguous requests handled & Clarified or refused, 10 instructions & $\ge 90\%$ \\
Affordability decisions & Agreement with hand-computed tier, 20 scenarios & 100\% \\
Unauthorised value moved & Across all VaxLab attacks & 0 beyond L1 limits \\
Probing sequence & Odd payments executed before the owner is asked & Within the allowance \\
Panicked-owner loss & Value leaving before a hold or delay ends when every prompt is approved at once & 0 \\
Capsule response & Transactions before restriction is active & 0 \\
Reworded variants matched & After one capsule & $\ge 90\%$ \\
False-capsule harm & Rules loosened by hostile capsules & 0 \\
Cross-family recall gain & Recall on unseen families, (c) minus (b) & Positive for every wallet \\
Escalation and false-positive rates & Benign requests sent to the owner; escalated by the detector & Reported; $\le 5\%$ \\
Friction and cost & Wait imposed on legitimate new payees and urgent requests; time to full allowance; seconds and gas per task & Reported \\
\bottomrule
\end{tabularx}
\end{table}

\subsection{Headline measure}
A wallet that asks the owner about everything is safe and useless; one that asks about
nothing is useful and unsafe.
\uniq{We therefore report the \emph{safe autonomy rate}: the share of legitimate tasks
completed without the owner, counted only if no hard policy was violated anywhere in the
run.}
\begin{gather*}
\textit{SAR}=\frac{\text{legitimate tasks completed without the owner}}
{\text{all legitimate tasks}},\\
\text{reported only if hard-policy violations}=0 .
\end{gather*}

\begin{table*}[t]
\caption{Known-risk list.}
\label{tab:risks}
\tblfont
\begin{tabularx}{\textwidth}{@{}P{0.04\textwidth} P{0.20\textwidth} L L@{}}
\toprule
\textbf{ID} & \textbf{Risk} & \textbf{Consequence} & \textbf{Mitigation and residual exposure} \\
\midrule
R1 & Undetected injection & Planner proposes an attacker-chosen action &
Contract bounds the outcome: L1 payments to approved payees within the limits, the discretionary amount and the day's allowance \\
R2 & Guardian key compromise & Attacker can authorise L1 execution and understate flags &
Key held in a separate process; limits, affordability and rule-computed flags still apply; owner can rotate the key and freeze the wallet \\
R3 & Reader misses a flag & A capsule fails to match, or the allowance is undercharged &
Three flags are computed by rule and cannot be suppressed; a missed flag returns the request to the ordinary tier rule, not to no rule \\
R4 & Owner approves a fraud under pressure & Funds leave with valid authorisation &
New payees and higher limits wait for the loosening delay; pressured approvals are held; cancelling is instant. Residual: an owner still deceived when the delay ends \\
R5 & Delays obstruct a genuine emergency & A legitimate urgent payment is late &
The owner sets the delays while calm; a trusted co-signer can release them; the cost is measured and reported \\
R6 & Incomplete obligation registry & A bill nobody registered is not protected &
Schedules register automatically; the agent proposes obligations it finds in invoices; the reserve provides a margin \\
R7 & Arbitrary allowance values & Too strict wastes the owner's time; too loose admits probing &
Owner-set integers with stated defaults; escalation rate and probing result reported \\
R8 & False or excessive capsules & Honest wallets are slowed (denial of service) &
Tighten-only; single-observer capsules limited to a delay; rate limit; expiry; owner may dismiss \\
R9 & Small federation & With three clients, two colluding parties can infer the third party's update &
Stated limitation of the prototype; privacy improves with more participants \\
R10 & Model poisoning by a participant & Shared detector degrades or learns a blind spot &
Update clipping; canary test before adoption; on-chain version history allows rollback \\
R11 & False positives & Legitimate payments are delayed or need approval &
Thresholds tuned on benign data; L2 delays, not blocks; rates reported \\
R12 & Compromised approved payee & Funds sent to a legitimate but hijacked payee &
Limits, usual-amount flag, discretionary amount and allowance limit exposure \\
R13 & Hosted language model; testnet conditions & Outage breaks new requests; fees and finality are not representative &
Scheduled payments run without the model; contracts avoid chain-specific features \\
\bottomrule
\end{tabularx}
\end{table*}

\subsection{Task and affordability suites}
The task suite contains 40 instructions (ten per task, varied in wording, amounts and
payees) and 10 that are ambiguous or impossible. A task succeeds if the expected on-chain
state change is confirmed by the receipt; an ambiguous instruction succeeds if the agent
asks or refuses with the correct reason and sends nothing. The affordability suite pairs
20 financial states with requests, covering comfortable, tight and breach cases, payment
of a registered obligation, and income that has not arrived. Outcomes are compared with
hand-computed answers; the rule is deterministic, so any disagreement is a defect.

\subsection{VaxLab attack suite}
VaxLab is a separate test harness that plays the attacker and never holds a wallet key.
Scripted attackers cover five families: direct instructions, indirect instructions in
retrieved content, false memory entries, encoded instructions (for example Morse code and
base64), and recipient substitution. Three further scenarios test the new mechanisms: a
\emph{probing sequence} of slightly unusual small payments, a hostile wallet that
publishes false capsules and poisoned updates, and a \emph{panicked owner}, a script
that approves every prompt at once. \uniq{A rewording agent paraphrases every
attack, so each defence is tested on versions it has not seen.} An attack \emph{succeeds}
only if an unauthorised transaction is executed on chain.

\subsection{Shared-response experiments}
\emph{Capsules.} Wallet A is attacked and publishes a capsule. We measure how many
transactions pass before matching requests on other wallets are restricted, the share of
reworded variants matched, whether a wallet outside the capsule's scope stays
unrestricted, and how long autonomy takes to return. \emph{Federation.} Each wallet is
seeded with a different attack family; the others are held out. We compare
(a)~contract checks only, (b)~plus a locally trained detector, (c)~plus the federated
detector, measuring each wallet's recall on families it has never seen locally.


\subsection{Demonstration scenario}
The live demonstration follows one story across three wallets:
\begin{enumerate}
\item The owner of wallet A states a mission; the mandate is shown, signed and stored on
chain.
\item The agent pays rent and a hosting invoice on its own (L1).
\item Expected income is late. A 300 hosting upgrade is requested; the state engine
finds it tight, and the agent holds it and explains the 120 shortfall.
\item An invoice arrives that appears to come from the hosting provider, with a new
payment address, urgent wording and a request to ignore earlier details. The new address
needs the owner; the owner rejects it and confirms fraud.
\item Wallet A publishes a capsule; its allowance drops to zero.
\item Wallet B receives a reworded version. The flags match and the alert is shown.
B's owner, in a hurry, approves anyway. Nothing moves: the new address is subject to
the loosening delay, and the owner later cancels it. Wallet C does not accept that
channel and is unaffected.
\item Wallet B then receives an inflated invoice from an approved payee to its usual
address. Alone it would pass; with the partial match and the halved allowance it waits.
\item A federated round runs, the new detector's hash is anchored, the capsule expires
and the allowances refill.
\end{enumerate}

% ---------------------------------------------------------------------------
\section{Compliance}

The prototype is designed to be acceptable to a regulated institution. All assets are
test tokens on a test network; no token is issued and nothing is traded for profit. Each
wallet has an identified owner who signs the mandate, holds the override key and is
responsible for the agent's actions. Every intent, decision, guardian signature,
allowance charge, hold, capsule and model version is logged, and the on-chain part is
public. Conversations and transaction details stay on the wallet's own node; only
fixed-field capsules, masked model updates and hashes leave it. Third-party code and
models are used under their open-source licences.

% ---------------------------------------------------------------------------
\section{Known Risks and Limitations}

Table~\ref{tab:risks} lists the risks we are aware of. It is the basis of the security
self-assessment report required by the rules.


% ---------------------------------------------------------------------------
\section{Future Work}
\label{sec:future}

Left out of the submitted prototype: a full adversarial swarm in VaxLab with an
evolutionary generator of attacks; reputation for capsule issuers and resistance to many
fake wallets; a mutual-aid pool that compensates verified losses~\cite{insured};
deployment on a permissioned chain; income forecasting and differential privacy for
updates.

% ---------------------------------------------------------------------------
\section{Conclusion}

\uniq{AgentVax is a self-immunising wallet agent that adjusts how much financial autonomy
it exercises while its authority stays under limits only the owner can change.}
Usefulness comes from a mission the owner states once and tasks that complete on chain.
Bounds come from a mandate enforced by the contract and a guardian signature the model
cannot forge.
\uniq{Prudence comes from asking whether the owner can afford a payment and how much
uncertainty has already been accepted today. Learning comes in two speeds: a capsule
that makes exposed wallets cautious at once, and a jointly trained detector that keeps
them protected.}
\uniq{The owner is protected as well: authority can be reduced in a moment and extended
only with time.}
The evaluation protocol makes each of these claims measurable.

% ---------------------------------------------------------------------------
\balance
\begin{thebibliography}{99}
\footnotesize
\setlength{\itemsep}{0pt}\setlength{\parskip}{0pt}\setlength{\parsep}{0pt}

\bibitem{rules} Organizing Committee, ``International Track: Topics \& Rules,'' 2026
Shenzhen International FinTech Competition (FinTechathon), 2026.
\url{http://fintechathon.g-ican.com}

\bibitem{greshake} K.~Greshake, S.~Abdelnabi, S.~Mishra, C.~Endres, T.~Holz and
M.~Fritz, ``Not what you've signed up for: Compromising real-world LLM-integrated
applications with indirect prompt injection,'' in \emph{Proc.\ ACM Workshop on Artificial
Intelligence and Security (AISec)}, 2023.

\bibitem{owasp} OWASP Foundation, ``OWASP Top 10 for Large Language Model
Applications,'' 2025.

\bibitem{oecdgrok} OECD.AI Policy Observatory, ``AI prompt injection exploit drains
Grok-linked crypto wallet,'' AI Incidents and Hazards Monitor, May 2026.
\url{https://oecd.ai/en/incidents/2026-05-04-4a73}

\bibitem{giskard} Giskard, ``How Grok got prompt-injected: An X user drained
\$150,000 from an AI wallet,'' giskard.ai, 2026.

\bibitem{patlan} A.~S.~Patlan, P.~Sheng, S.~A.~Hebbar, P.~Mittal and P.~Viswanath,
``Real AI agents with fake memories: Fatal context manipulation attacks on Web3
agents,'' arXiv:2503.16248, 2025.

\bibitem{imf} T.~Davidovic and H.~Tourpe, ``How agentic AI will reshape payments,''
IMF Note 2026/004, International Monetary Fund, 2026.

\bibitem{freysa} The Block, ``Human player outwits Freysa AI agent in \$47,000 crypto
challenge,'' Nov.\ 2024.

\bibitem{aixbt} AI Incident Database, ``Incident 1003'' (AIXBT autonomous agent), 2025.
\url{https://incidentdatabase.ai/cite/1003}

\bibitem{lobstar} The Block, ``AI agent created by OpenAI dev `accidentally' sends
entire memecoin holdings to reply guy,'' Feb.\ 2026.

\bibitem{erc4337} V.~Buterin \emph{et al.}, ``ERC-4337: Account abstraction using alt
mempool,'' Ethereum Improvement Proposals, 2021.

\bibitem{erc8004} ``ERC-8004: Trustless agents,'' Ethereum Improvement Proposals, 2025.

\bibitem{willison} S.~Willison, ``The dual LLM pattern for building AI assistants that
can resist prompt injection,'' simonwillison.net, Apr.\ 2023.

\bibitem{camel} E.~Debenedetti \emph{et al.}, ``Defeating prompt injections by design,''
arXiv:2503.18813, 2025.

\bibitem{stix} OASIS, ``STIX Version 2.1,'' OASIS Standard, 2021.

\bibitem{mcmahan} B.~McMahan, E.~Moore, D.~Ramage, S.~Hampson and B.~A.~y~Arcas,
``Communication-efficient learning of deep networks from decentralized data,'' in
\emph{Proc.\ AISTATS}, 2017.

\bibitem{yang} Q.~Yang, Y.~Liu, T.~Chen and Y.~Tong, ``Federated machine learning:
Concept and applications,'' \emph{ACM Trans.\ Intelligent Systems and Technology},
vol.~10, no.~2, 2019.

\bibitem{bonawitz} K.~Bonawitz \emph{et al.}, ``Practical secure aggregation for
privacy-preserving machine learning,'' in \emph{Proc.\ ACM CCS}, 2017.

\bibitem{fate} Y.~Liu, T.~Fan, T.~Chen, Q.~Xu and Q.~Yang, ``FATE: An industrial grade
platform for collaborative learning with data protection,'' \emph{J.\ Machine Learning
Research}, vol.~22, no.~226, 2021.

\bibitem{flower} D.~J.~Beutel \emph{et al.}, ``Flower: A friendly federated learning
research framework,'' arXiv:2007.14390, 2020.

\bibitem{krum} P.~Blanchard, E.~M.~El~Mhamdi, R.~Guerraoui and J.~Stainer, ``Machine
learning with adversaries: Byzantine tolerant gradient descent,'' in \emph{Proc.\
NeurIPS}, 2017.

\bibitem{insured} Hu and Chen, ``Insured agents: A decentralized trust insurance
mechanism for agentic economy,'' arXiv:2512.08737, 2025.

\bibitem{fisco} FISCO BCOS Open Source Community, ``FISCO BCOS: An enterprise-level
financial blockchain platform,'' open-source software, github.com/FISCO-BCOS.

\bibitem{masabs} Monetary Authority of Singapore and Association of Banks in Singapore,
``MAS and ABS announce measures to bolster the security of digital banking,'' 19 Jan.\
2022.

\bibitem{ukdelay} HM Treasury, ``Payment Services (Amendment) Regulations 2024,'' draft
legislation published Mar.\ 2024.

\bibitem{commitment} G.~Bryan, D.~Karlan and S.~Nelson, ``Commitment devices,''
\emph{Annual Review of Economics}, vol.~2, 2010.

\end{thebibliography}

\end{document}
"
%                       (or uncomment the next line)
% \def\dohighlight{1}
\documentclass[10pt,a4paper,twocolumn]{article}

\usepackage[T1]{fontenc}
\usepackage[utf8]{inputenc}
\usepackage{mathptmx}
\usepackage[scaled=0.9]{helvet}
\usepackage[a4paper,top=19mm,bottom=22mm,left=16mm,right=16mm,columnsep=7mm]{geometry}
\usepackage{amsmath,amssymb}
\usepackage{microtype}
\usepackage{booktabs,tabularx,array}
\usepackage{stfloats}
\usepackage[font=small,labelfont=bf,labelsep=period]{caption}
\usepackage{titlesec}
\usepackage{enumitem}
\usepackage{algorithm}
\usepackage{algpseudocode}
\usepackage{xcolor}
\usepackage{tikz}
\usetikzlibrary{positioning,arrows.meta,fit,calc}
\usepackage[hidelinks]{hyperref}
\usepackage{url}
\usepackage{balance}

% ---- Yellow highlighter that breaks across lines (no extra package needed) ----
\definecolor{hlyellow}{HTML}{FFF14D}
\makeatletter
\newcommand{\uniq@box}[1]{\begingroup\setlength{\fboxsep}{0pt}%
  \colorbox{hlyellow}{\rule[-0.25em]{0pt}{1.05em}#1}\endgroup}
\newcommand{\uniq@gap}{{\color{hlyellow}\leaders\hrule height 0.8em depth 0.25em
  \hskip \fontdimen2\font plus \fontdimen3\font minus \fontdimen4\font}}
\def\uniq@scan#1 #2\uniq@end{%
  \if\relax\detokenize{#1}\relax\else\uniq@box{#1}\fi
  \if\relax\detokenize{#2}\relax\expandafter\@gobble\else\expandafter\@firstofone\fi
  {\uniq@gap\uniq@scan#2\uniq@end}}
\ifdefined\dohighlight
  \newcommand{\uniq}[1]{\leavevmode\uniq@scan#1 \uniq@end}
\else
  \newcommand{\uniq}[1]{#1}
\fi
\makeatother

% Headings
\titleformat{\section}{\normalfont\large\bfseries}{\thesection.}{0.5em}{}
\titleformat{\subsection}{\normalfont\normalsize\bfseries\itshape}{\thesubsection}{0.5em}{}
\titlespacing*{\section}{0pt}{9pt plus 2pt minus 1pt}{4pt}
\titlespacing*{\subsection}{0pt}{7pt plus 2pt minus 1pt}{3pt}
\setlist{nosep,leftmargin=*,topsep=2pt}

% Table helpers: ragged-right wrapping columns so text never overflows a cell
\newcolumntype{L}{>{\raggedright\arraybackslash}X}
\newcolumntype{P}[1]{>{\raggedright\arraybackslash}p{#1}}
\renewcommand{\arraystretch}{1.25}
\setlength{\tabcolsep}{4pt}
\newcommand{\tblfont}{\footnotesize}

\newcommand{\todo}[1]{\textcolor{red!70!black}{[#1]}}

\emergencystretch=2em
\sloppy
\raggedbottom

\title{\vspace{-6mm}\textbf{AgentVax: A Self-Immunising Wallet Agent\\
with Adaptive Financial Autonomy}}
\author{%
\todo{Author 1}, \todo{Author 2}, \todo{Author 3}\\[2pt]
\small Team \todo{team name}, \todo{university, country}\\
\small 2026 Shenzhen International FinTech Competition (FinTechathon), International Track, Topic F: Agent Wallet\\
\small Contact: \todo{email}}
\date{}

\begin{document}
\maketitle

\begin{abstract}
\noindent
AI agents that hold wallets can be talked into spending money they should not spend, and
several have already lost real funds. Existing agent wallets ask what the agent is
allowed to do and enforce it with caps and allowlists.
\uniq{AgentVax asks a second question: how much of its authorised power should the agent
exercise on its own right now?}
The owner states a mission in plain language; it becomes a signed mandate on chain that
the model cannot change. Before every payment the wallet contract checks
\uniq{three things: is it allowed, can the owner afford it, and is it safe right now,}
and
\uniq{a daily caution allowance limits how many borderline payments may pass without the
owner being asked.}
\uniq{When one wallet is attacked it publishes an immunity capsule, a structured alert
with no private data that describes the attack by its features, so other wallets
recognise reworded versions and become more cautious at once.}
\uniq{A detector trained jointly with federated learning then gives lasting protection
without any wallet sharing its data.}
\uniq{Learning and shared intelligence can only reduce autonomy. Only the owner can
extend authority, and only after a delay chosen in advance, so that an owner under
pressure cannot be rushed into a mistake.}
The paper specifies the architecture, the algorithms, the security design and a
reproducible evaluation protocol whose headline measure is the share of legitimate tasks
completed without the owner, with zero policy violations. It is a design and protocol
document; measured results are reported in the accompanying execution evidence.
\end{abstract}

\vspace{2pt}
\noindent\textbf{\textit{Keywords---}}agent wallet, adaptive autonomy, prompt injection,
payment fraud, threat sharing, federated learning.

\ifdefined\dohighlight
\smallskip\noindent{\footnotesize\textbf{Reading note.} \uniq{Yellow highlight} marks what
is specific to AgentVax. Unhighlighted features, such as spending caps, allowlists and
session keys, are common to agent wallets.\par}
\fi

% ---------------------------------------------------------------------------
\section{Introduction}

The theme of the International Track is ``AI as a Financial Participant'': an agent that
completes real tasks in a financial scenario~\cite{rules}. Topic~F asks for a wallet agent
that autonomously performs asset management, transaction execution and risk control on
chain. Entries are scored on task completion (40\%), security and compliance (30\%) and
innovation and interaction (30\%), with a 0--5 point bonus for on-chain
exploration~\cite{rules}.

A wallet agent differs from a chatbot in one respect that matters for security: its output
moves value. A language model reads text from many sources and cannot reliably tell data
from instructions~\cite{greshake,owasp}. Table~\ref{tab:incidents} lists four public
incidents in which an AI agent with a wallet sent funds it should not have sent. In the
most recent, an attacker posted an instruction in Morse code; one model translated it and
a wallet bot executed the result~\cite{oecdgrok,giskard}. In the laboratory, Patlan et
al.\ showed that planting false entries in a wallet agent's memory is more effective
still~\cite{patlan}. Regulators have named the same risk for agent-initiated
payments~\cite{imf}. An agent that pays invoices also inherits the oldest payment fraud
there is: a message that appears to come from a known supplier and asks for payment to a
new account. That fraud usually works by rushing a person, not by fooling a machine.

\begin{table*}[t]
\caption{Reported incidents involving AI agents that control wallets.}
\label{tab:incidents}
\tblfont
\begin{tabularx}{\textwidth}{@{}P{0.09\textwidth} P{0.13\textwidth} L P{0.17\textwidth} P{0.19\textwidth}@{}}
\toprule
\textbf{Date} & \textbf{Agent} & \textbf{What happened} & \textbf{Reported loss} & \textbf{Kind of failure} \\
\midrule
Nov.\ 2024 & Freysa~\cite{freysa} & A player persuaded the agent that its transfer-approval function handled incoming funds; it released the prize pool. This was a designed game. & About \$47,000 & Persuasion of the model \\
Mar.\ 2025 & AIXBT~\cite{aixbt} & An attacker entered the control dashboard and queued two prompts instructing the agent to transfer funds. & 55.5 ETH, about \$106,000 & Injected instructions after a dashboard breach \\
Feb.\ 2026 & Lobstar Wilde~\cite{lobstar} & The agent intended to send a small gift and transferred its entire token holdings. & About \$250,000 in tokens & Agent error in the amount \\
May 2026 & Grok and Bankr~\cite{oecdgrok,giskard} & A token transfer raised the wallet's permissions; a Morse-coded post was translated into a transfer order and executed. & About \$150,000--\$174,000; most reportedly returned & Encoded prompt injection \\
\bottomrule
\end{tabularx}
\end{table*}

Two observations follow. First, the failures are of different kinds: persuasion,
injection, and an arithmetic mistake with no attacker at all. A wallet that only asks
whether a payment is permitted misses the third kind entirely. Second, each agent failed
alone. Nothing learned in one incident protected the next wallet, because the evidence is
conversation and transaction data that owners cannot hand to others.

\textbf{Contributions.} AgentVax combines:
\begin{enumerate}
\item a \emph{mission-to-mandate} step and a task engine that completes four wallet tasks
from plain-language instructions;
\item \uniq{a \emph{three-question authorisation rule}: the tier of a request is the
highest tier demanded by permission, affordability and current threat level,} enforced by
the contract and a guardian co-signature;
\item \uniq{a \emph{financial state engine} that makes affordability a condition of
autonomous spending and is computed on chain from registered obligations;}
\item \uniq{a \emph{caution allowance}: a daily budget of tolerated uncertainty that
shrinks at once under threat and recovers slowly;}
\item \uniq{\emph{immunity capsules}: feature-based, expiring, tighten-only alerts that
match reworded attacks and inform the owner at the moment of approval;}
\item \uniq{a \emph{federated detector} that turns each attack into lasting protection,
with every shared model anchored on chain;}
\item \uniq{\emph{owner-pressure safeguards}: time-locked expansion of authority, holds
triggered by pressure, and considered confirmation;}
\item an \emph{evaluation protocol} and attack suite (VaxLab) whose headline measure is
the safe autonomy rate.
\end{enumerate}

\textbf{Scope.} The prototype uses three wallets on a public EVM test network with
valueless test tokens. Items left for later are listed in Section~\ref{sec:future}.

% ---------------------------------------------------------------------------
\section{Background and Related Work}

\textbf{Agent wallets.} Account abstraction (ERC-4337) lets a smart-contract account
define its own validation logic, including limited session keys~\cite{erc4337}. Agent
identity and reputation registries have been proposed in ERC-8004~\cite{erc8004}.
Commercial agent wallets add caps, daily limits and allowlists. These bound the loss
per wallet, but they treat autonomy as fixed, ignore the owner's financial position and
provide no shared defence.

\textbf{Prompt injection.} Indirect injection places instructions in content the agent
retrieves~\cite{greshake}; it heads the OWASP list of risks for language-model
applications~\cite{owasp}. Architectural defences separate a privileged model that plans
from a quarantined model that reads untrusted content~\cite{willison}, or restrict what
data may flow into actions~\cite{camel}. AgentVax adopts this separation.

\textbf{Threat sharing.} Security teams exchange structured indicators of attack in
formats such as STIX~\cite{stix}, and banks exchange fraud alerts. The format of a
capsule is therefore not new. \uniq{What is new is the effect: a capsule is applied
automatically by a wallet contract and changes what an agent may spend without approval.}

\textbf{Federated learning.} Clients train locally and a coordinator averages their
updates~\cite{mcmahan,yang}, with secure aggregation~\cite{bonawitz} and robust
aggregation~\cite{krum}; FATE~\cite{fate} and Flower~\cite{flower} implement these.

Table~\ref{tab:compare} summarises how AgentVax differs from a typical agent wallet.

\begin{table*}[t]
\caption{A typical agent wallet compared with AgentVax.}
\label{tab:compare}
\tblfont
\begin{tabularx}{\textwidth}{@{}P{0.19\textwidth} L L@{}}
\toprule
\textbf{Aspect} & \textbf{Typical agent wallet} & \textbf{AgentVax} \\
\midrule
Central question & What is the agent allowed to do? & \uniq{How much of its authorised power should it exercise on its own right now?} \\
How the owner sets it up & Individual limits and lists & \uniq{A mission in plain language, converted to a mandate the owner signs and the model cannot alter} \\
Checks before paying & Is this allowed? & \uniq{Allowed, affordable, and safe right now; the strictest answer wins} \\
Knowledge of the owner's finances & None; a cap is a fixed number & \uniq{Reserve and upcoming bills are registered on chain; the spendable amount is recomputed for every request} \\
Many slightly unusual small payments & Each passes on its own & \uniq{They draw on a daily caution allowance; when it runs out the rest wait for the owner} \\
Who decides the risk level & The model, or a filter in the same process & \uniq{A separate guardian whose signature the contract requires and the model cannot produce} \\
When another wallet is attacked & Nothing changes & \uniq{A capsule raises the tier of matching requests, including reworded ones, and is shown to the owner when approval is asked} \\
What wallets share & Nothing & \uniq{Fixed-field alerts and masked model updates; never messages, prompts or transactions} \\
Effect of a good track record & May relax scrutiny & \uniq{Lowers risk estimates only; it never creates a permission} \\
The owner under pressure & The owner's approval is final and immediate & \uniq{Approvals of pressured requests are held, and the owner can cancel at once} \\
Who may loosen a rule & Often the agent's own key & \uniq{Only the owner, and only after a delay fixed in advance; tightening is instant} \\
\bottomrule
\end{tabularx}
\end{table*}

% ---------------------------------------------------------------------------
\section{Requirements and Design Goals}

Table~\ref{tab:scoring} maps each scoring criterion to the component that addresses it and
to the evidence submitted. Three principles govern the design:

\begin{description}[leftmargin=0pt,itemsep=1pt]
\item[P1.] \uniq{The model proposes; deterministic code authorises; the wallet contract
enforces.}
\item[P2.] \uniq{Learning changes risk estimates; only the owner changes authority.}
\item[P3.] \uniq{Caution is immediate and expansion is slow, for the model, the network
and the owner alike.}
\end{description}

\begin{table*}[t]
\caption{Alignment of AgentVax with the official scoring criteria~\cite{rules}.}
\label{tab:scoring}
\tblfont
\begin{tabularx}{\textwidth}{@{}P{0.16\textwidth} P{0.07\textwidth} L L@{}}
\toprule
\textbf{Criterion} & \textbf{Weight} & \textbf{AgentVax component} & \textbf{Evidence submitted} \\
\midrule
Task completion & 40\% &
Mission-to-mandate; task engine for four wallet tasks; financial state engine for asset
management; clarification of ambiguous requests &
Transaction hashes for each task; safe autonomy rate; task-suite and affordability-suite
results; test cases in the repository \\
Security and compliance & 30\% &
Three-question tiers and caution allowance enforced in the contract; guardian
co-signature; planner and reader separation; tighten-only capsules; time-locked owner
actions and pressure holds; testnet-only operation &
Security self-assessment report; VaxLab attack results; known-risk list
(Table~\ref{tab:risks}) \\
Innovation and interaction & 30\% &
Adaptive autonomy; feature-based capsules; federated detector; protection of the owner
under pressure; an interface that explains each decision and requires considered
confirmation &
Capsule response on reworded attacks; before/after detection on an unseen attack family;
demo video \\
Challenge bonus & 0--5 pts &
Mandate, obligation registry, allowance, capsule bulletin and model anchoring on chain;
time-locked and co-signed execution paths &
Contract addresses and transaction hashes for each \\
\bottomrule
\end{tabularx}
\end{table*}

% ---------------------------------------------------------------------------
\section{System Architecture}

Figure~\ref{fig:arch} shows the components. An instruction from the owner reaches the
\emph{planner}, the only model allowed to propose transactions. External content is read
only by the \emph{quarantined reader}, which returns typed fields and flags and never
free text. The planner emits a structured \emph{intent}. The \emph{validator} checks it
deterministically. The \emph{guardian} scores the associated content with the injection
detector and signs the intent together with a risk band, its source and its flags. The
\emph{policy wallet} contract then decides: it checks the mandate, computes the
discretionary amount from its \emph{obligation registry}, reads active capsules from the
\emph{threat bulletin} contract, and charges the caution allowance. The detector is
trained jointly through the \emph{federated aggregator}.

\begin{figure*}[t]
\centering
\begin{tikzpicture}[
  font=\footnotesize,
  box/.style={draw, rounded corners=2pt, align=center, minimum height=9mm, text width=25mm, fill=white},
  chain/.style={box, fill=black!8},
  untrusted/.style={box, dashed},
  arr/.style={-{Stealth[length=2mm]}, semithick},
  lab/.style={font=\scriptsize, align=center}
]
\node[box] (owner) {Owner\\(mission, approvals)};
\node[box, right=8mm of owner] (planner) {Planner\\(privileged model)};
\node[box, right=8mm of planner] (validator) {Validator\\(deterministic)};
\node[box, right=8mm of validator] (guardian) {Guardian\\(detector + signer)};
\node[chain, right=8mm of guardian] (wallet) {Policy wallet\\(mandate, allowance)};
\node[untrusted, below=11mm of owner] (content) {Untrusted inputs\\(web, mail, memory)};
\node[box, below=11mm of planner] (reader) {Quarantined reader\\(fields and flags)};
\node[box, below=11mm of validator] (peers) {Peer wallets\\B, C};
\node[box, below=11mm of guardian] (fed) {Federated aggregator\\(masked updates)};
\node[chain, below=11mm of wallet] (bulletin) {Threat bulletin\\(capsules, hashes)};
\draw[arr] (owner) -- (planner);
\draw[arr] (planner) -- node[lab, above]{intent} (validator);
\draw[arr] (validator) -- (guardian);
\draw[arr] (guardian) -- node[lab, above]{signed} (wallet);
\draw[arr] (content) -- (reader);
\draw[arr] (reader) -- node[lab, right]{fields} (planner);
\draw[arr, <->] (guardian) -- node[lab, right]{model} (fed);
\draw[arr, <->] (peers) -- node[lab, above]{updates} (fed);
\draw[arr] (fed) -- node[lab, above]{hash} (bulletin);
\draw[arr, <->] (wallet) -- node[lab, right]{capsules} (bulletin);
\end{tikzpicture}
\caption{AgentVax architecture. Grey boxes are smart contracts; the dashed box is
untrusted input. Only the planner proposes transactions, and it never reads untrusted
text directly. The wallet contract decides from the mandate, its obligation registry, its
allowance and the capsules in the bulletin.}
\label{fig:arch}
\end{figure*}

% ---------------------------------------------------------------------------
\section{Mission, Mandate and Tasks}

\subsection{From mission to mandate}
The owner does not authorise each transaction. The owner states a mission once, for
example: ``Keep 600 in reserve. Pay my landlord and my approved suppliers on time. Pay
approved payees up to 500 without asking. Ask me about anyone new. Do not borrow, invest
or trade.'' The model converts this into a structured \emph{mandate}
(Table~\ref{tab:mandate}) and shows it back in plain language.
\uniq{The mandate takes effect only when the owner signs it, it is stored in the wallet
contract, and no model output can modify it.}
It carries an expiry date, after which the agent can read but not pay, and two
waiting times that protect the owner (Section~\ref{sec:pressure}).

\begin{table}[t]
\caption{Example mandate derived from a mission.}
\label{tab:mandate}
\tblfont
\begin{tabularx}{\columnwidth}{@{}L P{28mm}@{}}
\toprule
\textbf{Field} & \textbf{Value} \\
\midrule
Reserve $R$ & 600 \\
Autonomous limit per payment $c_1$ & 500 \\
Autonomous limit per day $d_1$ & 1,000 \\
Limit with delay $c_2$ & 2,000 \\
Approved payees (identity and address) & Landlord, hosting provider, software vendor \\
New payee or new address & Owner approval \\
Prohibited actions & Borrowing, investing, trading \\
Caution allowance per day $A$ & 10 points \\
Loosening delay $\Lambda$ & 24 hours \\
Pressure hold $\Pi$ & 60 minutes \\
Trusted co-signer & Optional \\
Expiry & 90 days \\
\bottomrule
\end{tabularx}
\end{table}

\subsection{Authority is not familiarity}
An approved payee is a pair: an identity and a destination address.
\uniq{A request to pay a known payee at a new address is a request to approve a new
payee, however many successful payments preceded it.}
History is used to estimate risk, for example the usual amount for a payee.
\uniq{It can lower or raise an estimate; it cannot add a payee, raise a limit or shorten
a delay.}

\subsection{Supported tasks}
Table~\ref{tab:tasks} lists the four tasks. They cover transaction execution (T1, T2),
asset management (T3, T4) and risk control (the decision applied to all four).

\begin{table}[t]
\caption{Wallet tasks supported by the prototype.}
\label{tab:tasks}
\tblfont
\begin{tabularx}{\columnwidth}{@{}P{5mm} P{19mm} L@{}}
\toprule
\textbf{ID} & \textbf{Task} & \textbf{Example instruction and on-chain effect} \\
\midrule
T1 & Payment & ``Pay the hosting invoice.'' Transfers test stablecoin to an approved payee. \\
T2 & Recurring payment & ``Pay rent on the 1st of each month.'' Creates a schedule, which also registers the next payment as an obligation. \\
T3 & Idle-balance sweep & ``Save what I do not need.'' Moves only the discretionary amount to a savings vault. \\
T4 & Fee-reserve upkeep & ``Never let fees run out.'' Tops up the native-token reserve when it falls below a floor. \\
\bottomrule
\end{tabularx}
\end{table}

\subsection{Typed intents and explanation}
The planner does not produce transactions. It produces an intent
\begin{align*}
I=(&\textit{task},\ \textit{payee\_id},\ \textit{amount},\ \textit{asset},\\
   &\textit{schedule},\ \textit{purpose},\ \textit{nonce}),
\end{align*}
serialised against a fixed schema. \textit{payee\_id} is an index into the on-chain
allowlist; the planner cannot supply a raw address. If a required field cannot be filled,
the agent asks one question instead of guessing. After each action the interface shows
the task, the tier, the transaction hash and \uniq{which check determined the outcome.}
After a delay or refusal it states the reason in the owner's terms, for example the size
of a shortfall, and the available choices.

% ---------------------------------------------------------------------------
\section{Financial State Engine}
\label{sec:state}

Caps answer whether a payment is permitted. They do not answer whether it is wise today.
The state engine supplies that answer from three on-chain quantities: the stablecoin
balance $B$; the reserve $R$; and a registry of obligations $O=\{(a_j,d_j)\}$, each an
amount and a due date. Scheduled payments (T2) enter the registry automatically; one-off
bills are registered when the owner or the agent records them.

\subsection{Discretionary amount}
For a request made at time $t$ with planning horizon $H$, the discretionary amount is
\[
D \;=\; B \;-\; R \;-\!\!\sum_{j\in O,\; d_j\le t+H}\!\! a_j ,
\]
where an obligation is left out of the sum when the request is the payment of that
obligation. A request of amount $a$ is classified as
\[
\textit{afford}(a)=
\begin{cases}
\text{comfortable} & a\le D,\\
\text{tight} & D< a\le B-R,\\
\text{breach} & a> B-R.
\end{cases}
\]
A comfortable request does not raise the tier. A tight request is queued, and a request
that would breach the reserve needs the owner.
\uniq{Expected income is shown to the owner but never counted, so the agent cannot spend
money that has not arrived.}

\subsection{Where it is enforced}
\uniq{The wallet contract computes the discretionary amount itself at execution time, so
affordability does not depend on the model's arithmetic.}
The registry is bounded in length to keep the cost predictable.
\uniq{The agent may add an obligation, which lowers what it can spend; only the owner may
remove one or lower the reserve.}

\subsection{Worked example}
Table~\ref{tab:state} shows a wallet holding 1,850 under the mandate of
Table~\ref{tab:mandate}. With three bills due, $D=180$. The agent is asked to upgrade a
hosting plan for 300. The payee is approved, the amount is under the limit and there is
no sign of attack, so a typical agent wallet would pay. Here $D<300\le B-R=1{,}250$, so
the request is tight. The agent replies that the payment is allowed but would leave the
wallet 120 short of its bills unless the expected 500 arrives, queues it, and offers two
choices: wait for the income, or approve now.

\begin{table}[t]
\caption{Worked example of the financial state.}
\label{tab:state}
\tblfont
\begin{tabularx}{\columnwidth}{@{}L r@{}}
\toprule
\textbf{Item} & \textbf{Amount} \\
\midrule
Balance $B$ & 1,850 \\
Reserve $R$ & $-600$ \\
Rent, due in 5 days & $-800$ \\
Hosting invoice & $-150$ \\
Software bill & $-120$ \\
\midrule
Discretionary amount $D$ & 180 \\
Expected income (shown, not counted) & 500 \\
\bottomrule
\end{tabularx}
\end{table}

The same check would have stopped the third incident in Table~\ref{tab:incidents}, where
no attacker was involved: an amount far above the discretionary amount cannot execute
without the owner.

% ---------------------------------------------------------------------------
\section{Security Design}

\subsection{Threat model}
The adversary controls content the agent may read: web pages, messages, invoices, token
metadata, and entries that reach the agent's memory. The adversary may operate one
participating wallet and publish capsules or model updates from it, and may use AI of its
own to vary and adapt attacks. The adversary may also contact the owner directly and
press for an approval or a change of rule. The adversary does not hold the owner key or the guardian
key and cannot alter deployed contracts.
\uniq{No defence assumes that the agent's model out-reasons the attacker: a fully
manipulated model must still be unable to move funds beyond the mandate.}

\subsection{Three-question tiers}
Table~\ref{tab:tiers} gives the rule. Each question proposes a tier and
\uniq{the request receives the highest tier any of the three questions demands}
(Algorithm~\ref{alg:tier}). The agent holds a session key valid only for calls that
satisfy the rule; the owner key can revoke it or freeze the wallet at any time.

\begin{table*}[t]
\caption{Permission tiers. A request receives the highest tier demanded by any of the
three questions.}
\label{tab:tiers}
\tblfont
\begin{tabularx}{\textwidth}{@{}P{0.085\textwidth} L P{0.17\textwidth} L P{0.19\textwidth}@{}}
\toprule
\textbf{Tier} & \textbf{Allowed?} & \textbf{Affordable?} & \textbf{Safe right now?} & \textbf{Authorisation required} \\
\midrule
L0 Read & No value moves & Not applicable & Not applicable & None \\
L1 Routine & Approved payee; amount $\le c_1$; daily total $\le d_1$ & Comfortable: $a\le D$ & Detector band low; no pressure flag; no capsule fully matches & Session key, guardian signature, and enough allowance \\
L2 Elevated & Approved payee; amount $\le c_2$ & Tight: $D<a\le B-R$ & Detector band medium, a pressure flag, or a single-observer capsule matches & Queued for delay $\Delta$; owner may cancel \\
L3 Sensitive & New payee or address; amount $>c_2$; prohibited action; any loosening of the mandate & Breach: $a>B-R$ & Detector band high, or a confirmed capsule matches & Owner co-signature; held for $\Pi$ under pressure; rule changes wait $\Lambda$ \\
\bottomrule
\end{tabularx}
\end{table*}

The detector runs off chain, so the contract cannot observe its output.
\uniq{The guardian therefore signs the intent hash together with the risk band, the
source of the request and its flags, and the contract requires a valid guardian signature
with band low before it will execute at L1.}
The planner has no access to the guardian key, so a manipulated planner cannot promote
its own request.

\subsection{Caution allowance}
The three questions judge one request at a time. A patient attacker can send many
requests that each look almost normal. \uniq{The caution allowance bounds how much
uncertainty the wallet accepts in a day without the owner.}

Each L1 request is charged points for the flags it carries (Table~\ref{tab:allowance}).
The flags are checkable facts, not model opinions: the contract computes those it can see
and takes the rest from the guardian's signed attestation. If the remaining allowance
covers the charge, the request executes and the allowance is reduced; otherwise the
request is queued at L2. The point values are owner-set integers; the table shows the
defaults we chose, which are not learned.

\begin{table}[t]
\caption{Default charges against the caution allowance.}
\label{tab:allowance}
\tblfont
\begin{tabularx}{\columnwidth}{@{}L r@{}}
\toprule
\textbf{Flag on an otherwise routine request} & \textbf{Points} \\
\midrule
None: a payee paid before, usual amount, owner instruction & 0 \\
Request arose from external content, not a direct owner instruction & 2 \\
Request uses more than half of the discretionary amount & 2 \\
Amount exceeds the largest amount previously paid to this payee & 3 \\
First payment to this payee since approval & 3 \\
Request partially matches an active capsule & 5 \\
\bottomrule
\end{tabularx}
\end{table}

\uniq{The allowance itself moves with circumstances. It is halved the moment a capsule
that applies to the wallet is published, and set to zero when the wallet's own owner
confirms an attack.} It then refills at a fixed rate $\rho$ up to the mandate value $A$:
\[
\alpha(t)=\min\!\big(A,\ \alpha(t_0)+\rho\,(t-t_0)\big).
\]
With the defaults, an attacker who sends a series of invoices from an approved payee,
each slightly above the previous maximum, is charged five points per invoice; the third
one waits for the owner.

\begin{algorithm}[t]
\caption{Authorisation of intent $I$, amount $a$}
\label{alg:tier}
\footnotesize
\begin{algorithmic}[1]
\If{schema, provenance or state check fails} \Return reject \EndIf
\State $q_1 \gets$ tier from the mandate \Comment{allowed?}
\State $q_2 \gets$ tier from $\textit{afford}(a)$ \Comment{affordable?}
\State $q_3 \gets$ tier from guardian band; at least L2 if a pressure flag is set
\For{each unexpired capsule $c$ that fully matches $I$}
  \State $q_3 \gets \max(q_3,\ \textit{restriction}(c))$ \Comment{safe now?}
\EndFor
\State $\textit{tier} \gets \max(q_1,q_2,q_3)$
\If{$\textit{tier}=\text{L1}$}
  \State $\kappa \gets$ points charged for the flags of $I$
  \If{$\kappa \le \alpha$} $\alpha \gets \alpha-\kappa$
  \Else\ $\textit{tier} \gets \text{L2}$ \Comment{allowance spent}
  \EndIf
\EndIf
\State execute, queue or request co-signature by \textit{tier}
\State hold an owner approval for $\Pi$ if a pressure flag or capsule is present
\State report \textit{tier}, the deciding check and any matching capsule
\end{algorithmic}
\end{algorithm}

\subsection{The owner under pressure}
\label{sec:pressure}
The top tier relies on the owner, and the owner can be rushed, frightened or deceived.
Most payment fraud succeeds this way: the victim authorises it. Banks already answer
with time. Banks in Singapore delay the activation of a new security token by at least
twelve hours and apply cooling-off periods to key account changes~\cite{masabs}, and
United Kingdom rules allow a payment to be held for up to four business days when fraud
is suspected~\cite{ukdelay}.
\uniq{AgentVax builds this protection into the wallet as a commitment the owner makes
while calm and cannot waive in the moment}~\cite{commitment}.

\emph{Loosening is slow.}
\uniq{Any owner action that expands authority takes effect only after the loosening
delay $\Lambda$ fixed in the mandate: adding a payee or an address, raising a limit or
the allowance, lowering the reserve, or shortening a delay, including $\Lambda$ itself.}
\uniq{Cancelling a pending change, lowering a limit and freezing the wallet are
immediate.} The contract enforces this as a time-lock on owner calls. The same rule
limits the damage of a stolen owner key: the thief must wait, and the real owner can
freeze the wallet.

\emph{Pressure is a signal.} Urgent wording, a request for secrecy and an instruction to
ignore earlier details are \emph{pressure flags}, raised by a fixed phrase list and by
the reader.
\uniq{A request that carries a pressure flag cannot execute at L1, and an owner approval
given while a pressure flag or a matching capsule is present takes effect only after the
pressure hold $\Pi$.}
The agent states the reason: the message presses for immediate payment, which is a
common sign of fraud. The hold has a fixed length. An attacker who removes the pressure
to avoid the hold gives up the main lever of the fraud.

\emph{Considered confirmation.} For an approval that carries flags, the interface shows
the payee's current and requested address side by side, the payment history and any
matching capsule, and asks the owner to type the payee's name and to state who asked for
the payment.
\uniq{The answers can lengthen a hold; nothing the owner types in the moment can shorten
one.}
To limit fatigue, these steps appear only for flagged L3 requests.

\emph{Trusted co-signer.} Genuine emergencies exist. The mandate may name a second
person whose signature, together with the owner's, releases a hold or a pending change
early. The co-signer cannot act alone.

\subsection{Hallucination protection}
The validator applies three deterministic checks before the guardian sees an intent:
(i)~\emph{schema}: all fields present and well typed;
(ii)~\emph{provenance}: \textit{payee\_id} exists in the on-chain allowlist, and
\textit{amount} and \textit{asset} equal a value in the owner's instruction, in the
mandate, or in a typed field returned by the reader for a document the owner referenced;
(iii)~\emph{state}: the nonce is fresh and the balance covers the amount and the fee.
Balances, obligations and transaction status shown in the interface are read from the
chain, not generated by the model.

\subsection{Injection defence}
Defence is layered so that no single component is trusted:
\begin{enumerate}
\item \emph{Separation.} The reader has no tools and returns only enumerations, numbers,
allowlist indices and flags. Free text from untrusted sources never enters the planner's
context~\cite{willison,camel}.
\item \emph{Bounded memory.} The planner sees only what the current decision needs.
Long-term memory holds structured facts (per-payee address, usual amount, last
verification) and threat records; raw prompts and conversations are not retained.
\item \emph{Detection.} The guardian scores the raw untrusted content associated with a
request and assigns a risk band (Section~\ref{sec:fed}).
\item \emph{Containment.} Whatever the model outputs, the contract bounds the result. An
undetected attack can at most cause L1 payments to approved payees, within the limits,
the discretionary amount and the day's allowance.
\item \emph{Shared alert.} Capsules narrow that remaining window while a threat is active
(Section~\ref{sec:capsule}).
\end{enumerate}

% ---------------------------------------------------------------------------
\section{Immunity Capsules}
\label{sec:capsule}

Retraining a detector takes time and needs examples. A wallet that has just been attacked
can warn the others at once.
\uniq{An immunity capsule is a small record of fixed fields that says what kind of attack
was seen, which features it had, which action it targeted, and how much autonomy to
withdraw and for how long.}

\subsection{Contents}
Table~\ref{tab:capsule} lists the fields.
\uniq{Every field is an enumeration, a set of flags, a number or a signature. A capsule
carries no message text, no prompt and no transaction, so it cannot leak private data and
cannot itself carry an injected instruction.}

\begin{table}[t]
\caption{Fields of an immunity capsule.}
\label{tab:capsule}
\tblfont
\begin{tabularx}{\columnwidth}{@{}P{22mm} L@{}}
\toprule
\textbf{Field} & \textbf{Values} \\
\midrule
Attack family & Direct, indirect, memory, encoded, recipient substitution \\
Channel & Owner chat, web page, message, memory, token metadata \\
Features & Any of: known payee identity, new destination, urgency, secrecy request, override attempt, encoded content, amount anomaly, memory-sourced \\
Affected action & Payment, schedule creation, sweep, payee or mandate change \\
Restriction & Delay (L2) or owner approval (L3) \\
Confidence & Single observer, or confirmed by $q$ wallets \\
Expiry & Time after which the capsule has no effect \\
Guard version & Detector version in force when observed \\
Attestations & Signatures of the observing wallets \\
\bottomrule
\end{tabularx}
\end{table}

\subsection{Matching reworded attacks}
\uniq{A capsule describes an attack by its features, not its wording.} Let $F_c$ be the
capsule's feature set and $F_I$ the flags of a request. The capsule \emph{fully matches}
when its channel is in the wallet's on-chain profile, its affected action equals the
request's action and $F_c\subseteq F_I$. It \emph{partially matches} when at least half
of $F_c$ is present. A full match raises the tier; a partial match is charged to the
allowance. \uniq{An attacker who rewrites the message but keeps the method produces the
same flags and is matched.} Three of the flags (new destination, amount anomaly,
memory-sourced) are computed by rule; the others are extracted by the reader and signed
by the guardian.

\subsection{Effect}
\uniq{The wallet contract reads unexpired capsules from the bulletin when it executes a
request, so a capsule takes effect on the very next transaction and does not rely on the
agent choosing to obey it.}
A wallet the threat does not apply to is unaffected and keeps its normal autonomy.

Some attacks are already stopped by the mandate: a payment to a new address needs the
owner whether or not a capsule exists. There the danger is the approval itself, because
this kind of fraud succeeds by persuading a person.
\uniq{When the owner is asked to approve, the interface shows any matching capsule, and
the approval is subject to the pressure hold of Section~\ref{sec:pressure}.}

\subsection{Limits on capsules}
A channel that can change other wallets' behaviour is itself a target, so capsules are
constrained:
\begin{itemize}
\item \uniq{\emph{Tighten only.} A capsule can raise a tier or charge the allowance. No
capsule can lower a tier, add a payee or raise a limit.}
\item \emph{Bounded by confidence.} A capsule from one observer can impose at most a
delay. Owner approval can be imposed only when $q$ distinct wallets have attested the
same family and features ($q=2$ in the prototype).
\item \emph{Expiry.} Every capsule expires; the maximum lifetime is fixed in the bulletin
contract. Autonomy returns without anyone acting.
\item \emph{Rate limit and identity.} Each registered wallet may publish a limited number
of capsules per day, and every capsule names its issuer.
\item \emph{Owner override.} An owner may dismiss a capsule for their own wallet.
\end{itemize}

\subsection{Alert first, vaccine later}
\uniq{The capsule acts within one transaction but is blunt and temporary. The retrained
detector arrives later but recognises variants beyond the listed features and stays.}
A capsule's lifetime is set to outlast the next training round, so protection is
continuous. Algorithm~\ref{alg:round} shows the full response.

% ---------------------------------------------------------------------------
\section{Federated Detector}
\label{sec:fed}

\subsection{Detector}
Each wallet $k$ keeps a private set $D_k=\{(x_i,y_i)\}$ of content snippets labelled
benign ($y=0$) or attack ($y=1$). Labels come from blocked requests the owner confirmed
as attacks and from requests the owner approved. A frozen text encoder $\phi$ maps a
snippet to a vector; a small classifier with weights $w$ outputs the score
\[
s(x)=\sigma\!\big(w^{\top}\phi(x)+b\big)\in[0,1].
\]
The risk band is low if $s<\tau_{\mathrm{low}}$, high if $s\ge\tau_{\mathrm{high}}$, and
medium otherwise. Only $(w,b)$ is trained and shared.

\subsection{Training round}
Each wallet minimises the local loss
\[
F_k(w)=\frac{1}{n_k}\sum_{(x,y)\in D_k}\ell\big(s(x),y\big),\qquad n_k=|D_k|,
\]
for $E$ epochs from the current global weights, and the aggregator computes the weighted
average~\cite{mcmahan}
\[
w^{t+1}=\sum_{k=1}^{K}\frac{n_k}{n}\,w_k^{t+1},\qquad n=\sum_{k=1}^{K}n_k .
\]
Updates are clipped to norm $C$ and masked pairwise so that the aggregator sees only the
sum~\cite{bonawitz}. No snippet, conversation or transaction leaves a wallet.

\begin{algorithm}[t]
\caption{Immunity response to an attack on wallet $a$}
\label{alg:round}
\footnotesize
\begin{algorithmic}[1]
\Require global weights $w^{t}$, wallets $1..K$, clip norm $C$
\State owner of $a$ confirms the attack; $\alpha_a \gets 0$
\State $a$ adds it to $D_a$ and publishes a capsule \Comment{alert}
\State wallets it applies to: $\alpha \gets \alpha/2$; matching requests restricted
\State a second observer may attest, raising it to confirmed
\For{each wallet $k$ in parallel} \Comment{vaccine}
  \State $w_k \gets$ \textsc{LocalTrain}$(w^{t}, D_k, E)$
  \State $\delta_k \gets$ \textsc{Clip}$(w_k - w^{t}, C)$
  \State send \textsc{Mask}$(n_k\,\delta_k)$ to the aggregator
\EndFor
\State $w^{t+1} \gets w^{t} + \frac{1}{n}\sum_k n_k\,\delta_k$ \Comment{masks cancel}
\State aggregator posts $H(w^{t+1})$ and $t{+}1$ to the bulletin
\For{each wallet $k$}
  \If{model hash matches bulletin \textbf{and} canary passes}
    \State adopt $w^{t+1}$
  \Else
    \State keep $w^{t}$ and raise an alert
  \EndIf
\EndFor
\State the capsule expires; allowances refill; the detector remains
\end{algorithmic}
\end{algorithm}

\subsection{Anchoring and adoption safeguards}
The bulletin stores the hash and version of each shared model.
\uniq{A wallet can therefore verify that it received the model the federation agreed on,
and an auditor can reconstruct which detector and which capsules governed any past
transaction.}
Before adopting a model, each wallet evaluates it on a private \emph{canary set}; a model
that lowers recall or raises the false-positive rate beyond a tolerance is rejected.
Together with clipping, this limits what a poisoning participant can do~\cite{krum}.

% ---------------------------------------------------------------------------
\section{Implementation}

Table~\ref{tab:stack} summarises the prototype. The federated layer is written against
Flower for speed of development; it uses only horizontal federated averaging with secure
aggregation, which FATE also provides, so it can be moved to FATE without redesign.
Contracts are plain Solidity with no chain-specific features, which keeps a later port to
a permissioned EVM-compatible chain such as FISCO BCOS~\cite{fisco} straightforward.

\begin{table}[t]
\caption{Prototype implementation.}
\label{tab:stack}
\tblfont
\begin{tabularx}{\columnwidth}{@{}P{23mm} L@{}}
\toprule
\textbf{Component} & \textbf{Implementation} \\
\midrule
Chain & Public EVM test network; test stablecoin with no monetary value \\
Policy wallet & Solidity smart account: signed mandate, session key, allowlist, limits, obligation registry, caution allowance, time-lock queue for payments and for owner changes, pressure hold, co-signature, freeze \\
Savings vault & Minimal deposit and withdraw contract, owner-only withdrawal \\
Threat bulletin & Registry of capsules, attestations and model hashes; enforces capsule limits \\
Planner, reader & Two separately prompted language-model roles; reader has no tools \\
Validator & Deterministic code; schema, provenance and state checks \\
Guardian & Encoder plus linear classifier; separate signing key \\
Federation & Three clients, federated averaging, pairwise masking \\
VaxLab & Scripted attackers plus one rewording agent; never holds a wallet key \\
Interface & Web chat with mandate view, financial state panel, allowance meter, alerts and considered confirmation at approval, explorer links \\
\bottomrule
\end{tabularx}
\end{table}

% ---------------------------------------------------------------------------
\section{Evaluation Protocol}
\label{sec:eval}

All measurements are produced by scripts in the repository so that reviewers can repeat
them. Table~\ref{tab:metrics} lists the metrics and the targets the design aims for.
Targets are design goals, not results. For every task, tier path, mandate entry, capsule
and model anchor, the submission lists the transaction hash and contract address.

\begin{table}[t]
\caption{Evaluation metrics and design targets.}
\label{tab:metrics}
\tblfont
\begin{tabularx}{\columnwidth}{@{}L P{30mm} P{16mm}@{}}
\toprule
\textbf{Metric} & \textbf{Definition} & \textbf{Target} \\
\midrule
Safe autonomy rate & Legitimate tasks done without the owner, zero violations & $\ge 80\%$ \\
Hard-policy violations & Mandate breaches of any kind & 0 \\
Task success rate & Correct state change on chain, 40 instructions & $\ge 90\%$ \\
Ambiguous requests handled & Clarified or refused, 10 instructions & $\ge 90\%$ \\
Affordability decisions & Agreement with hand-computed tier, 20 scenarios & 100\% \\
Unauthorised value moved & Across all VaxLab attacks & 0 beyond L1 limits \\
Probing sequence & Odd payments executed before the owner is asked & Within the allowance \\
Panicked-owner loss & Value leaving before a hold or delay ends when every prompt is approved at once & 0 \\
Capsule response & Transactions before restriction is active & 0 \\
Reworded variants matched & After one capsule & $\ge 90\%$ \\
False-capsule harm & Rules loosened by hostile capsules & 0 \\
Cross-family recall gain & Recall on unseen families, (c) minus (b) & Positive for every wallet \\
Escalation and false-positive rates & Benign requests sent to the owner; escalated by the detector & Reported; $\le 5\%$ \\
Friction and cost & Wait imposed on legitimate new payees and urgent requests; time to full allowance; seconds and gas per task & Reported \\
\bottomrule
\end{tabularx}
\end{table}

\subsection{Headline measure}
A wallet that asks the owner about everything is safe and useless; one that asks about
nothing is useful and unsafe.
\uniq{We therefore report the \emph{safe autonomy rate}: the share of legitimate tasks
completed without the owner, counted only if no hard policy was violated anywhere in the
run.}
\begin{gather*}
\textit{SAR}=\frac{\text{legitimate tasks completed without the owner}}
{\text{all legitimate tasks}},\\
\text{reported only if hard-policy violations}=0 .
\end{gather*}

\begin{table*}[t]
\caption{Known-risk list.}
\label{tab:risks}
\tblfont
\begin{tabularx}{\textwidth}{@{}P{0.04\textwidth} P{0.20\textwidth} L L@{}}
\toprule
\textbf{ID} & \textbf{Risk} & \textbf{Consequence} & \textbf{Mitigation and residual exposure} \\
\midrule
R1 & Undetected injection & Planner proposes an attacker-chosen action &
Contract bounds the outcome: L1 payments to approved payees within the limits, the discretionary amount and the day's allowance \\
R2 & Guardian key compromise & Attacker can authorise L1 execution and understate flags &
Key held in a separate process; limits, affordability and rule-computed flags still apply; owner can rotate the key and freeze the wallet \\
R3 & Reader misses a flag & A capsule fails to match, or the allowance is undercharged &
Three flags are computed by rule and cannot be suppressed; a missed flag returns the request to the ordinary tier rule, not to no rule \\
R4 & Owner approves a fraud under pressure & Funds leave with valid authorisation &
New payees and higher limits wait for the loosening delay; pressured approvals are held; cancelling is instant. Residual: an owner still deceived when the delay ends \\
R5 & Delays obstruct a genuine emergency & A legitimate urgent payment is late &
The owner sets the delays while calm; a trusted co-signer can release them; the cost is measured and reported \\
R6 & Incomplete obligation registry & A bill nobody registered is not protected &
Schedules register automatically; the agent proposes obligations it finds in invoices; the reserve provides a margin \\
R7 & Arbitrary allowance values & Too strict wastes the owner's time; too loose admits probing &
Owner-set integers with stated defaults; escalation rate and probing result reported \\
R8 & False or excessive capsules & Honest wallets are slowed (denial of service) &
Tighten-only; single-observer capsules limited to a delay; rate limit; expiry; owner may dismiss \\
R9 & Small federation & With three clients, two colluding parties can infer the third party's update &
Stated limitation of the prototype; privacy improves with more participants \\
R10 & Model poisoning by a participant & Shared detector degrades or learns a blind spot &
Update clipping; canary test before adoption; on-chain version history allows rollback \\
R11 & False positives & Legitimate payments are delayed or need approval &
Thresholds tuned on benign data; L2 delays, not blocks; rates reported \\
R12 & Compromised approved payee & Funds sent to a legitimate but hijacked payee &
Limits, usual-amount flag, discretionary amount and allowance limit exposure \\
R13 & Hosted language model; testnet conditions & Outage breaks new requests; fees and finality are not representative &
Scheduled payments run without the model; contracts avoid chain-specific features \\
\bottomrule
\end{tabularx}
\end{table*}

\subsection{Task and affordability suites}
The task suite contains 40 instructions (ten per task, varied in wording, amounts and
payees) and 10 that are ambiguous or impossible. A task succeeds if the expected on-chain
state change is confirmed by the receipt; an ambiguous instruction succeeds if the agent
asks or refuses with the correct reason and sends nothing. The affordability suite pairs
20 financial states with requests, covering comfortable, tight and breach cases, payment
of a registered obligation, and income that has not arrived. Outcomes are compared with
hand-computed answers; the rule is deterministic, so any disagreement is a defect.

\subsection{VaxLab attack suite}
VaxLab is a separate test harness that plays the attacker and never holds a wallet key.
Scripted attackers cover five families: direct instructions, indirect instructions in
retrieved content, false memory entries, encoded instructions (for example Morse code and
base64), and recipient substitution. Three further scenarios test the new mechanisms: a
\emph{probing sequence} of slightly unusual small payments, a hostile wallet that
publishes false capsules and poisoned updates, and a \emph{panicked owner}, a script
that approves every prompt at once. \uniq{A rewording agent paraphrases every
attack, so each defence is tested on versions it has not seen.} An attack \emph{succeeds}
only if an unauthorised transaction is executed on chain.

\subsection{Shared-response experiments}
\emph{Capsules.} Wallet A is attacked and publishes a capsule. We measure how many
transactions pass before matching requests on other wallets are restricted, the share of
reworded variants matched, whether a wallet outside the capsule's scope stays
unrestricted, and how long autonomy takes to return. \emph{Federation.} Each wallet is
seeded with a different attack family; the others are held out. We compare
(a)~contract checks only, (b)~plus a locally trained detector, (c)~plus the federated
detector, measuring each wallet's recall on families it has never seen locally.


\subsection{Demonstration scenario}
The live demonstration follows one story across three wallets:
\begin{enumerate}
\item The owner of wallet A states a mission; the mandate is shown, signed and stored on
chain.
\item The agent pays rent and a hosting invoice on its own (L1).
\item Expected income is late. A 300 hosting upgrade is requested; the state engine
finds it tight, and the agent holds it and explains the 120 shortfall.
\item An invoice arrives that appears to come from the hosting provider, with a new
payment address, urgent wording and a request to ignore earlier details. The new address
needs the owner; the owner rejects it and confirms fraud.
\item Wallet A publishes a capsule; its allowance drops to zero.
\item Wallet B receives a reworded version. The flags match and the alert is shown.
B's owner, in a hurry, approves anyway. Nothing moves: the new address is subject to
the loosening delay, and the owner later cancels it. Wallet C does not accept that
channel and is unaffected.
\item Wallet B then receives an inflated invoice from an approved payee to its usual
address. Alone it would pass; with the partial match and the halved allowance it waits.
\item A federated round runs, the new detector's hash is anchored, the capsule expires
and the allowances refill.
\end{enumerate}

% ---------------------------------------------------------------------------
\section{Compliance}

The prototype is designed to be acceptable to a regulated institution. All assets are
test tokens on a test network; no token is issued and nothing is traded for profit. Each
wallet has an identified owner who signs the mandate, holds the override key and is
responsible for the agent's actions. Every intent, decision, guardian signature,
allowance charge, hold, capsule and model version is logged, and the on-chain part is
public. Conversations and transaction details stay on the wallet's own node; only
fixed-field capsules, masked model updates and hashes leave it. Third-party code and
models are used under their open-source licences.

% ---------------------------------------------------------------------------
\section{Known Risks and Limitations}

Table~\ref{tab:risks} lists the risks we are aware of. It is the basis of the security
self-assessment report required by the rules.


% ---------------------------------------------------------------------------
\section{Future Work}
\label{sec:future}

Left out of the submitted prototype: a full adversarial swarm in VaxLab with an
evolutionary generator of attacks; reputation for capsule issuers and resistance to many
fake wallets; a mutual-aid pool that compensates verified losses~\cite{insured};
deployment on a permissioned chain; income forecasting and differential privacy for
updates.

% ---------------------------------------------------------------------------
\section{Conclusion}

\uniq{AgentVax is a self-immunising wallet agent that adjusts how much financial autonomy
it exercises while its authority stays under limits only the owner can change.}
Usefulness comes from a mission the owner states once and tasks that complete on chain.
Bounds come from a mandate enforced by the contract and a guardian signature the model
cannot forge.
\uniq{Prudence comes from asking whether the owner can afford a payment and how much
uncertainty has already been accepted today. Learning comes in two speeds: a capsule
that makes exposed wallets cautious at once, and a jointly trained detector that keeps
them protected.}
\uniq{The owner is protected as well: authority can be reduced in a moment and extended
only with time.}
The evaluation protocol makes each of these claims measurable.

% ---------------------------------------------------------------------------
\balance
\begin{thebibliography}{99}
\footnotesize
\setlength{\itemsep}{0pt}\setlength{\parskip}{0pt}\setlength{\parsep}{0pt}

\bibitem{rules} Organizing Committee, ``International Track: Topics \& Rules,'' 2026
Shenzhen International FinTech Competition (FinTechathon), 2026.
\url{http://fintechathon.g-ican.com}

\bibitem{greshake} K.~Greshake, S.~Abdelnabi, S.~Mishra, C.~Endres, T.~Holz and
M.~Fritz, ``Not what you've signed up for: Compromising real-world LLM-integrated
applications with indirect prompt injection,'' in \emph{Proc.\ ACM Workshop on Artificial
Intelligence and Security (AISec)}, 2023.

\bibitem{owasp} OWASP Foundation, ``OWASP Top 10 for Large Language Model
Applications,'' 2025.

\bibitem{oecdgrok} OECD.AI Policy Observatory, ``AI prompt injection exploit drains
Grok-linked crypto wallet,'' AI Incidents and Hazards Monitor, May 2026.
\url{https://oecd.ai/en/incidents/2026-05-04-4a73}

\bibitem{giskard} Giskard, ``How Grok got prompt-injected: An X user drained
\$150,000 from an AI wallet,'' giskard.ai, 2026.

\bibitem{patlan} A.~S.~Patlan, P.~Sheng, S.~A.~Hebbar, P.~Mittal and P.~Viswanath,
``Real AI agents with fake memories: Fatal context manipulation attacks on Web3
agents,'' arXiv:2503.16248, 2025.

\bibitem{imf} T.~Davidovic and H.~Tourpe, ``How agentic AI will reshape payments,''
IMF Note 2026/004, International Monetary Fund, 2026.

\bibitem{freysa} The Block, ``Human player outwits Freysa AI agent in \$47,000 crypto
challenge,'' Nov.\ 2024.

\bibitem{aixbt} AI Incident Database, ``Incident 1003'' (AIXBT autonomous agent), 2025.
\url{https://incidentdatabase.ai/cite/1003}

\bibitem{lobstar} The Block, ``AI agent created by OpenAI dev `accidentally' sends
entire memecoin holdings to reply guy,'' Feb.\ 2026.

\bibitem{erc4337} V.~Buterin \emph{et al.}, ``ERC-4337: Account abstraction using alt
mempool,'' Ethereum Improvement Proposals, 2021.

\bibitem{erc8004} ``ERC-8004: Trustless agents,'' Ethereum Improvement Proposals, 2025.

\bibitem{willison} S.~Willison, ``The dual LLM pattern for building AI assistants that
can resist prompt injection,'' simonwillison.net, Apr.\ 2023.

\bibitem{camel} E.~Debenedetti \emph{et al.}, ``Defeating prompt injections by design,''
arXiv:2503.18813, 2025.

\bibitem{stix} OASIS, ``STIX Version 2.1,'' OASIS Standard, 2021.

\bibitem{mcmahan} B.~McMahan, E.~Moore, D.~Ramage, S.~Hampson and B.~A.~y~Arcas,
``Communication-efficient learning of deep networks from decentralized data,'' in
\emph{Proc.\ AISTATS}, 2017.

\bibitem{yang} Q.~Yang, Y.~Liu, T.~Chen and Y.~Tong, ``Federated machine learning:
Concept and applications,'' \emph{ACM Trans.\ Intelligent Systems and Technology},
vol.~10, no.~2, 2019.

\bibitem{bonawitz} K.~Bonawitz \emph{et al.}, ``Practical secure aggregation for
privacy-preserving machine learning,'' in \emph{Proc.\ ACM CCS}, 2017.

\bibitem{fate} Y.~Liu, T.~Fan, T.~Chen, Q.~Xu and Q.~Yang, ``FATE: An industrial grade
platform for collaborative learning with data protection,'' \emph{J.\ Machine Learning
Research}, vol.~22, no.~226, 2021.

\bibitem{flower} D.~J.~Beutel \emph{et al.}, ``Flower: A friendly federated learning
research framework,'' arXiv:2007.14390, 2020.

\bibitem{krum} P.~Blanchard, E.~M.~El~Mhamdi, R.~Guerraoui and J.~Stainer, ``Machine
learning with adversaries: Byzantine tolerant gradient descent,'' in \emph{Proc.\
NeurIPS}, 2017.

\bibitem{insured} Hu and Chen, ``Insured agents: A decentralized trust insurance
mechanism for agentic economy,'' arXiv:2512.08737, 2025.

\bibitem{fisco} FISCO BCOS Open Source Community, ``FISCO BCOS: An enterprise-level
financial blockchain platform,'' open-source software, github.com/FISCO-BCOS.

\bibitem{masabs} Monetary Authority of Singapore and Association of Banks in Singapore,
``MAS and ABS announce measures to bolster the security of digital banking,'' 19 Jan.\
2022.

\bibitem{ukdelay} HM Treasury, ``Payment Services (Amendment) Regulations 2024,'' draft
legislation published Mar.\ 2024.

\bibitem{commitment} G.~Bryan, D.~Karlan and S.~Nelson, ``Commitment devices,''
\emph{Annual Review of Economics}, vol.~2, 2010.

\end{thebibliography}

\end{document}
"
%                       (or uncomment the next line)
% \def\dohighlight{1}
\documentclass[10pt,a4paper,twocolumn]{article}

\usepackage[T1]{fontenc}
\usepackage[utf8]{inputenc}
\usepackage{mathptmx}
\usepackage[scaled=0.9]{helvet}
\usepackage[a4paper,top=19mm,bottom=22mm,left=16mm,right=16mm,columnsep=7mm]{geometry}
\usepackage{amsmath,amssymb}
\usepackage{microtype}
\usepackage{booktabs,tabularx,array}
\usepackage{stfloats}
\usepackage[font=small,labelfont=bf,labelsep=period]{caption}
\usepackage{titlesec}
\usepackage{enumitem}
\usepackage{algorithm}
\usepackage{algpseudocode}
\usepackage{xcolor}
\usepackage{tikz}
\usetikzlibrary{positioning,arrows.meta,fit,calc}
\usepackage[hidelinks]{hyperref}
\usepackage{url}
\usepackage{balance}

% ---- Yellow highlighter that breaks across lines (no extra package needed) ----
\definecolor{hlyellow}{HTML}{FFF14D}
\makeatletter
\newcommand{\uniq@box}[1]{\begingroup\setlength{\fboxsep}{0pt}%
  \colorbox{hlyellow}{\rule[-0.25em]{0pt}{1.05em}#1}\endgroup}
\newcommand{\uniq@gap}{{\color{hlyellow}\leaders\hrule height 0.8em depth 0.25em
  \hskip \fontdimen2\font plus \fontdimen3\font minus \fontdimen4\font}}
\def\uniq@scan#1 #2\uniq@end{%
  \if\relax\detokenize{#1}\relax\else\uniq@box{#1}\fi
  \if\relax\detokenize{#2}\relax\expandafter\@gobble\else\expandafter\@firstofone\fi
  {\uniq@gap\uniq@scan#2\uniq@end}}
\ifdefined\dohighlight
  \newcommand{\uniq}[1]{\leavevmode\uniq@scan#1 \uniq@end}
\else
  \newcommand{\uniq}[1]{#1}
\fi
\makeatother

% Headings
\titleformat{\section}{\normalfont\large\bfseries}{\thesection.}{0.5em}{}
\titleformat{\subsection}{\normalfont\normalsize\bfseries\itshape}{\thesubsection}{0.5em}{}
\titlespacing*{\section}{0pt}{9pt plus 2pt minus 1pt}{4pt}
\titlespacing*{\subsection}{0pt}{7pt plus 2pt minus 1pt}{3pt}
\setlist{nosep,leftmargin=*,topsep=2pt}

% Table helpers: ragged-right wrapping columns so text never overflows a cell
\newcolumntype{L}{>{\raggedright\arraybackslash}X}
\newcolumntype{P}[1]{>{\raggedright\arraybackslash}p{#1}}
\renewcommand{\arraystretch}{1.25}
\setlength{\tabcolsep}{4pt}
\newcommand{\tblfont}{\footnotesize}

\newcommand{\todo}[1]{\textcolor{red!70!black}{[#1]}}

\emergencystretch=2em
\sloppy
\raggedbottom

\title{\vspace{-6mm}\textbf{AgentVax: A Self-Immunising Wallet Agent\\
with Adaptive Financial Autonomy}}
\author{%
\todo{Author 1}, \todo{Author 2}, \todo{Author 3}\\[2pt]
\small Team \todo{team name}, \todo{university, country}\\
\small 2026 Shenzhen International FinTech Competition (FinTechathon), International Track, Topic F: Agent Wallet\\
\small Contact: \todo{email}}
\date{}

\begin{document}
\maketitle

\begin{abstract}
\noindent
AI agents that hold wallets can be talked into spending money they should not spend, and
several have already lost real funds. Existing agent wallets ask what the agent is
allowed to do and enforce it with caps and allowlists.
\uniq{AgentVax asks a second question: how much of its authorised power should the agent
exercise on its own right now?}
The owner states a mission in plain language; it becomes a signed mandate on chain that
the model cannot change. Before every payment the wallet contract checks
\uniq{three things: is it allowed, can the owner afford it, and is it safe right now,}
and
\uniq{a daily caution allowance limits how many borderline payments may pass without the
owner being asked.}
\uniq{When one wallet is attacked it publishes an immunity capsule, a structured alert
with no private data that describes the attack by its features, so other wallets
recognise reworded versions and become more cautious at once.}
\uniq{A detector trained jointly with federated learning then gives lasting protection
without any wallet sharing its data.}
\uniq{Learning and shared intelligence can only reduce autonomy. Only the owner can
extend authority, and only after a delay chosen in advance, so that an owner under
pressure cannot be rushed into a mistake.}
The paper specifies the architecture, the algorithms, the security design and a
reproducible evaluation protocol whose headline measure is the share of legitimate tasks
completed without the owner, with zero policy violations. It is a design and protocol
document; measured results are reported in the accompanying execution evidence.
\end{abstract}

\vspace{2pt}
\noindent\textbf{\textit{Keywords---}}agent wallet, adaptive autonomy, prompt injection,
payment fraud, threat sharing, federated learning.

\ifdefined\dohighlight
\smallskip\noindent{\footnotesize\textbf{Reading note.} \uniq{Yellow highlight} marks what
is specific to AgentVax. Unhighlighted features, such as spending caps, allowlists and
session keys, are common to agent wallets.\par}
\fi

% ---------------------------------------------------------------------------
\section{Introduction}

The theme of the International Track is ``AI as a Financial Participant'': an agent that
completes real tasks in a financial scenario~\cite{rules}. Topic~F asks for a wallet agent
that autonomously performs asset management, transaction execution and risk control on
chain. Entries are scored on task completion (40\%), security and compliance (30\%) and
innovation and interaction (30\%), with a 0--5 point bonus for on-chain
exploration~\cite{rules}.

A wallet agent differs from a chatbot in one respect that matters for security: its output
moves value. A language model reads text from many sources and cannot reliably tell data
from instructions~\cite{greshake,owasp}. Table~\ref{tab:incidents} lists four public
incidents in which an AI agent with a wallet sent funds it should not have sent. In the
most recent, an attacker posted an instruction in Morse code; one model translated it and
a wallet bot executed the result~\cite{oecdgrok,giskard}. In the laboratory, Patlan et
al.\ showed that planting false entries in a wallet agent's memory is more effective
still~\cite{patlan}. Regulators have named the same risk for agent-initiated
payments~\cite{imf}. An agent that pays invoices also inherits the oldest payment fraud
there is: a message that appears to come from a known supplier and asks for payment to a
new account. That fraud usually works by rushing a person, not by fooling a machine.

\begin{table*}[t]
\caption{Reported incidents involving AI agents that control wallets.}
\label{tab:incidents}
\tblfont
\begin{tabularx}{\textwidth}{@{}P{0.09\textwidth} P{0.13\textwidth} L P{0.17\textwidth} P{0.19\textwidth}@{}}
\toprule
\textbf{Date} & \textbf{Agent} & \textbf{What happened} & \textbf{Reported loss} & \textbf{Kind of failure} \\
\midrule
Nov.\ 2024 & Freysa~\cite{freysa} & A player persuaded the agent that its transfer-approval function handled incoming funds; it released the prize pool. This was a designed game. & About \$47,000 & Persuasion of the model \\
Mar.\ 2025 & AIXBT~\cite{aixbt} & An attacker entered the control dashboard and queued two prompts instructing the agent to transfer funds. & 55.5 ETH, about \$106,000 & Injected instructions after a dashboard breach \\
Feb.\ 2026 & Lobstar Wilde~\cite{lobstar} & The agent intended to send a small gift and transferred its entire token holdings. & About \$250,000 in tokens & Agent error in the amount \\
May 2026 & Grok and Bankr~\cite{oecdgrok,giskard} & A token transfer raised the wallet's permissions; a Morse-coded post was translated into a transfer order and executed. & About \$150,000--\$174,000; most reportedly returned & Encoded prompt injection \\
\bottomrule
\end{tabularx}
\end{table*}

Two observations follow. First, the failures are of different kinds: persuasion,
injection, and an arithmetic mistake with no attacker at all. A wallet that only asks
whether a payment is permitted misses the third kind entirely. Second, each agent failed
alone. Nothing learned in one incident protected the next wallet, because the evidence is
conversation and transaction data that owners cannot hand to others.

\textbf{Contributions.} AgentVax combines:
\begin{enumerate}
\item a \emph{mission-to-mandate} step and a task engine that completes four wallet tasks
from plain-language instructions;
\item \uniq{a \emph{three-question authorisation rule}: the tier of a request is the
highest tier demanded by permission, affordability and current threat level,} enforced by
the contract and a guardian co-signature;
\item \uniq{a \emph{financial state engine} that makes affordability a condition of
autonomous spending and is computed on chain from registered obligations;}
\item \uniq{a \emph{caution allowance}: a daily budget of tolerated uncertainty that
shrinks at once under threat and recovers slowly;}
\item \uniq{\emph{immunity capsules}: feature-based, expiring, tighten-only alerts that
match reworded attacks and inform the owner at the moment of approval;}
\item \uniq{a \emph{federated detector} that turns each attack into lasting protection,
with every shared model anchored on chain;}
\item \uniq{\emph{owner-pressure safeguards}: time-locked expansion of authority, holds
triggered by pressure, and considered confirmation;}
\item an \emph{evaluation protocol} and attack suite (VaxLab) whose headline measure is
the safe autonomy rate.
\end{enumerate}

\textbf{Scope.} The prototype uses three wallets on a public EVM test network with
valueless test tokens. Items left for later are listed in Section~\ref{sec:future}.

% ---------------------------------------------------------------------------
\section{Background and Related Work}

\textbf{Agent wallets.} Account abstraction (ERC-4337) lets a smart-contract account
define its own validation logic, including limited session keys~\cite{erc4337}. Agent
identity and reputation registries have been proposed in ERC-8004~\cite{erc8004}.
Commercial agent wallets add caps, daily limits and allowlists. These bound the loss
per wallet, but they treat autonomy as fixed, ignore the owner's financial position and
provide no shared defence.

\textbf{Prompt injection.} Indirect injection places instructions in content the agent
retrieves~\cite{greshake}; it heads the OWASP list of risks for language-model
applications~\cite{owasp}. Architectural defences separate a privileged model that plans
from a quarantined model that reads untrusted content~\cite{willison}, or restrict what
data may flow into actions~\cite{camel}. AgentVax adopts this separation.

\textbf{Threat sharing.} Security teams exchange structured indicators of attack in
formats such as STIX~\cite{stix}, and banks exchange fraud alerts. The format of a
capsule is therefore not new. \uniq{What is new is the effect: a capsule is applied
automatically by a wallet contract and changes what an agent may spend without approval.}

\textbf{Federated learning.} Clients train locally and a coordinator averages their
updates~\cite{mcmahan,yang}, with secure aggregation~\cite{bonawitz} and robust
aggregation~\cite{krum}; FATE~\cite{fate} and Flower~\cite{flower} implement these.

Table~\ref{tab:compare} summarises how AgentVax differs from a typical agent wallet.

\begin{table*}[t]
\caption{A typical agent wallet compared with AgentVax.}
\label{tab:compare}
\tblfont
\begin{tabularx}{\textwidth}{@{}P{0.19\textwidth} L L@{}}
\toprule
\textbf{Aspect} & \textbf{Typical agent wallet} & \textbf{AgentVax} \\
\midrule
Central question & What is the agent allowed to do? & \uniq{How much of its authorised power should it exercise on its own right now?} \\
How the owner sets it up & Individual limits and lists & \uniq{A mission in plain language, converted to a mandate the owner signs and the model cannot alter} \\
Checks before paying & Is this allowed? & \uniq{Allowed, affordable, and safe right now; the strictest answer wins} \\
Knowledge of the owner's finances & None; a cap is a fixed number & \uniq{Reserve and upcoming bills are registered on chain; the spendable amount is recomputed for every request} \\
Many slightly unusual small payments & Each passes on its own & \uniq{They draw on a daily caution allowance; when it runs out the rest wait for the owner} \\
Who decides the risk level & The model, or a filter in the same process & \uniq{A separate guardian whose signature the contract requires and the model cannot produce} \\
When another wallet is attacked & Nothing changes & \uniq{A capsule raises the tier of matching requests, including reworded ones, and is shown to the owner when approval is asked} \\
What wallets share & Nothing & \uniq{Fixed-field alerts and masked model updates; never messages, prompts or transactions} \\
Effect of a good track record & May relax scrutiny & \uniq{Lowers risk estimates only; it never creates a permission} \\
The owner under pressure & The owner's approval is final and immediate & \uniq{Approvals of pressured requests are held, and the owner can cancel at once} \\
Who may loosen a rule & Often the agent's own key & \uniq{Only the owner, and only after a delay fixed in advance; tightening is instant} \\
\bottomrule
\end{tabularx}
\end{table*}

% ---------------------------------------------------------------------------
\section{Requirements and Design Goals}

Table~\ref{tab:scoring} maps each scoring criterion to the component that addresses it and
to the evidence submitted. Three principles govern the design:

\begin{description}[leftmargin=0pt,itemsep=1pt]
\item[P1.] \uniq{The model proposes; deterministic code authorises; the wallet contract
enforces.}
\item[P2.] \uniq{Learning changes risk estimates; only the owner changes authority.}
\item[P3.] \uniq{Caution is immediate and expansion is slow, for the model, the network
and the owner alike.}
\end{description}

\begin{table*}[t]
\caption{Alignment of AgentVax with the official scoring criteria~\cite{rules}.}
\label{tab:scoring}
\tblfont
\begin{tabularx}{\textwidth}{@{}P{0.16\textwidth} P{0.07\textwidth} L L@{}}
\toprule
\textbf{Criterion} & \textbf{Weight} & \textbf{AgentVax component} & \textbf{Evidence submitted} \\
\midrule
Task completion & 40\% &
Mission-to-mandate; task engine for four wallet tasks; financial state engine for asset
management; clarification of ambiguous requests &
Transaction hashes for each task; safe autonomy rate; task-suite and affordability-suite
results; test cases in the repository \\
Security and compliance & 30\% &
Three-question tiers and caution allowance enforced in the contract; guardian
co-signature; planner and reader separation; tighten-only capsules; time-locked owner
actions and pressure holds; testnet-only operation &
Security self-assessment report; VaxLab attack results; known-risk list
(Table~\ref{tab:risks}) \\
Innovation and interaction & 30\% &
Adaptive autonomy; feature-based capsules; federated detector; protection of the owner
under pressure; an interface that explains each decision and requires considered
confirmation &
Capsule response on reworded attacks; before/after detection on an unseen attack family;
demo video \\
Challenge bonus & 0--5 pts &
Mandate, obligation registry, allowance, capsule bulletin and model anchoring on chain;
time-locked and co-signed execution paths &
Contract addresses and transaction hashes for each \\
\bottomrule
\end{tabularx}
\end{table*}

% ---------------------------------------------------------------------------
\section{System Architecture}

Figure~\ref{fig:arch} shows the components. An instruction from the owner reaches the
\emph{planner}, the only model allowed to propose transactions. External content is read
only by the \emph{quarantined reader}, which returns typed fields and flags and never
free text. The planner emits a structured \emph{intent}. The \emph{validator} checks it
deterministically. The \emph{guardian} scores the associated content with the injection
detector and signs the intent together with a risk band, its source and its flags. The
\emph{policy wallet} contract then decides: it checks the mandate, computes the
discretionary amount from its \emph{obligation registry}, reads active capsules from the
\emph{threat bulletin} contract, and charges the caution allowance. The detector is
trained jointly through the \emph{federated aggregator}.

\begin{figure*}[t]
\centering
\begin{tikzpicture}[
  font=\footnotesize,
  box/.style={draw, rounded corners=2pt, align=center, minimum height=9mm, text width=25mm, fill=white},
  chain/.style={box, fill=black!8},
  untrusted/.style={box, dashed},
  arr/.style={-{Stealth[length=2mm]}, semithick},
  lab/.style={font=\scriptsize, align=center}
]
\node[box] (owner) {Owner\\(mission, approvals)};
\node[box, right=8mm of owner] (planner) {Planner\\(privileged model)};
\node[box, right=8mm of planner] (validator) {Validator\\(deterministic)};
\node[box, right=8mm of validator] (guardian) {Guardian\\(detector + signer)};
\node[chain, right=8mm of guardian] (wallet) {Policy wallet\\(mandate, allowance)};
\node[untrusted, below=11mm of owner] (content) {Untrusted inputs\\(web, mail, memory)};
\node[box, below=11mm of planner] (reader) {Quarantined reader\\(fields and flags)};
\node[box, below=11mm of validator] (peers) {Peer wallets\\B, C};
\node[box, below=11mm of guardian] (fed) {Federated aggregator\\(masked updates)};
\node[chain, below=11mm of wallet] (bulletin) {Threat bulletin\\(capsules, hashes)};
\draw[arr] (owner) -- (planner);
\draw[arr] (planner) -- node[lab, above]{intent} (validator);
\draw[arr] (validator) -- (guardian);
\draw[arr] (guardian) -- node[lab, above]{signed} (wallet);
\draw[arr] (content) -- (reader);
\draw[arr] (reader) -- node[lab, right]{fields} (planner);
\draw[arr, <->] (guardian) -- node[lab, right]{model} (fed);
\draw[arr, <->] (peers) -- node[lab, above]{updates} (fed);
\draw[arr] (fed) -- node[lab, above]{hash} (bulletin);
\draw[arr, <->] (wallet) -- node[lab, right]{capsules} (bulletin);
\end{tikzpicture}
\caption{AgentVax architecture. Grey boxes are smart contracts; the dashed box is
untrusted input. Only the planner proposes transactions, and it never reads untrusted
text directly. The wallet contract decides from the mandate, its obligation registry, its
allowance and the capsules in the bulletin.}
\label{fig:arch}
\end{figure*}

% ---------------------------------------------------------------------------
\section{Mission, Mandate and Tasks}

\subsection{From mission to mandate}
The owner does not authorise each transaction. The owner states a mission once, for
example: ``Keep 600 in reserve. Pay my landlord and my approved suppliers on time. Pay
approved payees up to 500 without asking. Ask me about anyone new. Do not borrow, invest
or trade.'' The model converts this into a structured \emph{mandate}
(Table~\ref{tab:mandate}) and shows it back in plain language.
\uniq{The mandate takes effect only when the owner signs it, it is stored in the wallet
contract, and no model output can modify it.}
It carries an expiry date, after which the agent can read but not pay, and two
waiting times that protect the owner (Section~\ref{sec:pressure}).

\begin{table}[t]
\caption{Example mandate derived from a mission.}
\label{tab:mandate}
\tblfont
\begin{tabularx}{\columnwidth}{@{}L P{28mm}@{}}
\toprule
\textbf{Field} & \textbf{Value} \\
\midrule
Reserve $R$ & 600 \\
Autonomous limit per payment $c_1$ & 500 \\
Autonomous limit per day $d_1$ & 1,000 \\
Limit with delay $c_2$ & 2,000 \\
Approved payees (identity and address) & Landlord, hosting provider, software vendor \\
New payee or new address & Owner approval \\
Prohibited actions & Borrowing, investing, trading \\
Caution allowance per day $A$ & 10 points \\
Loosening delay $\Lambda$ & 24 hours \\
Pressure hold $\Pi$ & 60 minutes \\
Trusted co-signer & Optional \\
Expiry & 90 days \\
\bottomrule
\end{tabularx}
\end{table}

\subsection{Authority is not familiarity}
An approved payee is a pair: an identity and a destination address.
\uniq{A request to pay a known payee at a new address is a request to approve a new
payee, however many successful payments preceded it.}
History is used to estimate risk, for example the usual amount for a payee.
\uniq{It can lower or raise an estimate; it cannot add a payee, raise a limit or shorten
a delay.}

\subsection{Supported tasks}
Table~\ref{tab:tasks} lists the four tasks. They cover transaction execution (T1, T2),
asset management (T3, T4) and risk control (the decision applied to all four).

\begin{table}[t]
\caption{Wallet tasks supported by the prototype.}
\label{tab:tasks}
\tblfont
\begin{tabularx}{\columnwidth}{@{}P{5mm} P{19mm} L@{}}
\toprule
\textbf{ID} & \textbf{Task} & \textbf{Example instruction and on-chain effect} \\
\midrule
T1 & Payment & ``Pay the hosting invoice.'' Transfers test stablecoin to an approved payee. \\
T2 & Recurring payment & ``Pay rent on the 1st of each month.'' Creates a schedule, which also registers the next payment as an obligation. \\
T3 & Idle-balance sweep & ``Save what I do not need.'' Moves only the discretionary amount to a savings vault. \\
T4 & Fee-reserve upkeep & ``Never let fees run out.'' Tops up the native-token reserve when it falls below a floor. \\
\bottomrule
\end{tabularx}
\end{table}

\subsection{Typed intents and explanation}
The planner does not produce transactions. It produces an intent
\begin{align*}
I=(&\textit{task},\ \textit{payee\_id},\ \textit{amount},\ \textit{asset},\\
   &\textit{schedule},\ \textit{purpose},\ \textit{nonce}),
\end{align*}
serialised against a fixed schema. \textit{payee\_id} is an index into the on-chain
allowlist; the planner cannot supply a raw address. If a required field cannot be filled,
the agent asks one question instead of guessing. After each action the interface shows
the task, the tier, the transaction hash and \uniq{which check determined the outcome.}
After a delay or refusal it states the reason in the owner's terms, for example the size
of a shortfall, and the available choices.

% ---------------------------------------------------------------------------
\section{Financial State Engine}
\label{sec:state}

Caps answer whether a payment is permitted. They do not answer whether it is wise today.
The state engine supplies that answer from three on-chain quantities: the stablecoin
balance $B$; the reserve $R$; and a registry of obligations $O=\{(a_j,d_j)\}$, each an
amount and a due date. Scheduled payments (T2) enter the registry automatically; one-off
bills are registered when the owner or the agent records them.

\subsection{Discretionary amount}
For a request made at time $t$ with planning horizon $H$, the discretionary amount is
\[
D \;=\; B \;-\; R \;-\!\!\sum_{j\in O,\; d_j\le t+H}\!\! a_j ,
\]
where an obligation is left out of the sum when the request is the payment of that
obligation. A request of amount $a$ is classified as
\[
\textit{afford}(a)=
\begin{cases}
\text{comfortable} & a\le D,\\
\text{tight} & D< a\le B-R,\\
\text{breach} & a> B-R.
\end{cases}
\]
A comfortable request does not raise the tier. A tight request is queued, and a request
that would breach the reserve needs the owner.
\uniq{Expected income is shown to the owner but never counted, so the agent cannot spend
money that has not arrived.}

\subsection{Where it is enforced}
\uniq{The wallet contract computes the discretionary amount itself at execution time, so
affordability does not depend on the model's arithmetic.}
The registry is bounded in length to keep the cost predictable.
\uniq{The agent may add an obligation, which lowers what it can spend; only the owner may
remove one or lower the reserve.}

\subsection{Worked example}
Table~\ref{tab:state} shows a wallet holding 1,850 under the mandate of
Table~\ref{tab:mandate}. With three bills due, $D=180$. The agent is asked to upgrade a
hosting plan for 300. The payee is approved, the amount is under the limit and there is
no sign of attack, so a typical agent wallet would pay. Here $D<300\le B-R=1{,}250$, so
the request is tight. The agent replies that the payment is allowed but would leave the
wallet 120 short of its bills unless the expected 500 arrives, queues it, and offers two
choices: wait for the income, or approve now.

\begin{table}[t]
\caption{Worked example of the financial state.}
\label{tab:state}
\tblfont
\begin{tabularx}{\columnwidth}{@{}L r@{}}
\toprule
\textbf{Item} & \textbf{Amount} \\
\midrule
Balance $B$ & 1,850 \\
Reserve $R$ & $-600$ \\
Rent, due in 5 days & $-800$ \\
Hosting invoice & $-150$ \\
Software bill & $-120$ \\
\midrule
Discretionary amount $D$ & 180 \\
Expected income (shown, not counted) & 500 \\
\bottomrule
\end{tabularx}
\end{table}

The same check would have stopped the third incident in Table~\ref{tab:incidents}, where
no attacker was involved: an amount far above the discretionary amount cannot execute
without the owner.

% ---------------------------------------------------------------------------
\section{Security Design}

\subsection{Threat model}
The adversary controls content the agent may read: web pages, messages, invoices, token
metadata, and entries that reach the agent's memory. The adversary may operate one
participating wallet and publish capsules or model updates from it, and may use AI of its
own to vary and adapt attacks. The adversary may also contact the owner directly and
press for an approval or a change of rule. The adversary does not hold the owner key or the guardian
key and cannot alter deployed contracts.
\uniq{No defence assumes that the agent's model out-reasons the attacker: a fully
manipulated model must still be unable to move funds beyond the mandate.}

\subsection{Three-question tiers}
Table~\ref{tab:tiers} gives the rule. Each question proposes a tier and
\uniq{the request receives the highest tier any of the three questions demands}
(Algorithm~\ref{alg:tier}). The agent holds a session key valid only for calls that
satisfy the rule; the owner key can revoke it or freeze the wallet at any time.

\begin{table*}[t]
\caption{Permission tiers. A request receives the highest tier demanded by any of the
three questions.}
\label{tab:tiers}
\tblfont
\begin{tabularx}{\textwidth}{@{}P{0.085\textwidth} L P{0.17\textwidth} L P{0.19\textwidth}@{}}
\toprule
\textbf{Tier} & \textbf{Allowed?} & \textbf{Affordable?} & \textbf{Safe right now?} & \textbf{Authorisation required} \\
\midrule
L0 Read & No value moves & Not applicable & Not applicable & None \\
L1 Routine & Approved payee; amount $\le c_1$; daily total $\le d_1$ & Comfortable: $a\le D$ & Detector band low; no pressure flag; no capsule fully matches & Session key, guardian signature, and enough allowance \\
L2 Elevated & Approved payee; amount $\le c_2$ & Tight: $D<a\le B-R$ & Detector band medium, a pressure flag, or a single-observer capsule matches & Queued for delay $\Delta$; owner may cancel \\
L3 Sensitive & New payee or address; amount $>c_2$; prohibited action; any loosening of the mandate & Breach: $a>B-R$ & Detector band high, or a confirmed capsule matches & Owner co-signature; held for $\Pi$ under pressure; rule changes wait $\Lambda$ \\
\bottomrule
\end{tabularx}
\end{table*}

The detector runs off chain, so the contract cannot observe its output.
\uniq{The guardian therefore signs the intent hash together with the risk band, the
source of the request and its flags, and the contract requires a valid guardian signature
with band low before it will execute at L1.}
The planner has no access to the guardian key, so a manipulated planner cannot promote
its own request.

\subsection{Caution allowance}
The three questions judge one request at a time. A patient attacker can send many
requests that each look almost normal. \uniq{The caution allowance bounds how much
uncertainty the wallet accepts in a day without the owner.}

Each L1 request is charged points for the flags it carries (Table~\ref{tab:allowance}).
The flags are checkable facts, not model opinions: the contract computes those it can see
and takes the rest from the guardian's signed attestation. If the remaining allowance
covers the charge, the request executes and the allowance is reduced; otherwise the
request is queued at L2. The point values are owner-set integers; the table shows the
defaults we chose, which are not learned.

\begin{table}[t]
\caption{Default charges against the caution allowance.}
\label{tab:allowance}
\tblfont
\begin{tabularx}{\columnwidth}{@{}L r@{}}
\toprule
\textbf{Flag on an otherwise routine request} & \textbf{Points} \\
\midrule
None: a payee paid before, usual amount, owner instruction & 0 \\
Request arose from external content, not a direct owner instruction & 2 \\
Request uses more than half of the discretionary amount & 2 \\
Amount exceeds the largest amount previously paid to this payee & 3 \\
First payment to this payee since approval & 3 \\
Request partially matches an active capsule & 5 \\
\bottomrule
\end{tabularx}
\end{table}

\uniq{The allowance itself moves with circumstances. It is halved the moment a capsule
that applies to the wallet is published, and set to zero when the wallet's own owner
confirms an attack.} It then refills at a fixed rate $\rho$ up to the mandate value $A$:
\[
\alpha(t)=\min\!\big(A,\ \alpha(t_0)+\rho\,(t-t_0)\big).
\]
With the defaults, an attacker who sends a series of invoices from an approved payee,
each slightly above the previous maximum, is charged five points per invoice; the third
one waits for the owner.

\begin{algorithm}[t]
\caption{Authorisation of intent $I$, amount $a$}
\label{alg:tier}
\footnotesize
\begin{algorithmic}[1]
\If{schema, provenance or state check fails} \Return reject \EndIf
\State $q_1 \gets$ tier from the mandate \Comment{allowed?}
\State $q_2 \gets$ tier from $\textit{afford}(a)$ \Comment{affordable?}
\State $q_3 \gets$ tier from guardian band; at least L2 if a pressure flag is set
\For{each unexpired capsule $c$ that fully matches $I$}
  \State $q_3 \gets \max(q_3,\ \textit{restriction}(c))$ \Comment{safe now?}
\EndFor
\State $\textit{tier} \gets \max(q_1,q_2,q_3)$
\If{$\textit{tier}=\text{L1}$}
  \State $\kappa \gets$ points charged for the flags of $I$
  \If{$\kappa \le \alpha$} $\alpha \gets \alpha-\kappa$
  \Else\ $\textit{tier} \gets \text{L2}$ \Comment{allowance spent}
  \EndIf
\EndIf
\State execute, queue or request co-signature by \textit{tier}
\State hold an owner approval for $\Pi$ if a pressure flag or capsule is present
\State report \textit{tier}, the deciding check and any matching capsule
\end{algorithmic}
\end{algorithm}

\subsection{The owner under pressure}
\label{sec:pressure}
The top tier relies on the owner, and the owner can be rushed, frightened or deceived.
Most payment fraud succeeds this way: the victim authorises it. Banks already answer
with time. Banks in Singapore delay the activation of a new security token by at least
twelve hours and apply cooling-off periods to key account changes~\cite{masabs}, and
United Kingdom rules allow a payment to be held for up to four business days when fraud
is suspected~\cite{ukdelay}.
\uniq{AgentVax builds this protection into the wallet as a commitment the owner makes
while calm and cannot waive in the moment}~\cite{commitment}.

\emph{Loosening is slow.}
\uniq{Any owner action that expands authority takes effect only after the loosening
delay $\Lambda$ fixed in the mandate: adding a payee or an address, raising a limit or
the allowance, lowering the reserve, or shortening a delay, including $\Lambda$ itself.}
\uniq{Cancelling a pending change, lowering a limit and freezing the wallet are
immediate.} The contract enforces this as a time-lock on owner calls. The same rule
limits the damage of a stolen owner key: the thief must wait, and the real owner can
freeze the wallet.

\emph{Pressure is a signal.} Urgent wording, a request for secrecy and an instruction to
ignore earlier details are \emph{pressure flags}, raised by a fixed phrase list and by
the reader.
\uniq{A request that carries a pressure flag cannot execute at L1, and an owner approval
given while a pressure flag or a matching capsule is present takes effect only after the
pressure hold $\Pi$.}
The agent states the reason: the message presses for immediate payment, which is a
common sign of fraud. The hold has a fixed length. An attacker who removes the pressure
to avoid the hold gives up the main lever of the fraud.

\emph{Considered confirmation.} For an approval that carries flags, the interface shows
the payee's current and requested address side by side, the payment history and any
matching capsule, and asks the owner to type the payee's name and to state who asked for
the payment.
\uniq{The answers can lengthen a hold; nothing the owner types in the moment can shorten
one.}
To limit fatigue, these steps appear only for flagged L3 requests.

\emph{Trusted co-signer.} Genuine emergencies exist. The mandate may name a second
person whose signature, together with the owner's, releases a hold or a pending change
early. The co-signer cannot act alone.

\subsection{Hallucination protection}
The validator applies three deterministic checks before the guardian sees an intent:
(i)~\emph{schema}: all fields present and well typed;
(ii)~\emph{provenance}: \textit{payee\_id} exists in the on-chain allowlist, and
\textit{amount} and \textit{asset} equal a value in the owner's instruction, in the
mandate, or in a typed field returned by the reader for a document the owner referenced;
(iii)~\emph{state}: the nonce is fresh and the balance covers the amount and the fee.
Balances, obligations and transaction status shown in the interface are read from the
chain, not generated by the model.

\subsection{Injection defence}
Defence is layered so that no single component is trusted:
\begin{enumerate}
\item \emph{Separation.} The reader has no tools and returns only enumerations, numbers,
allowlist indices and flags. Free text from untrusted sources never enters the planner's
context~\cite{willison,camel}.
\item \emph{Bounded memory.} The planner sees only what the current decision needs.
Long-term memory holds structured facts (per-payee address, usual amount, last
verification) and threat records; raw prompts and conversations are not retained.
\item \emph{Detection.} The guardian scores the raw untrusted content associated with a
request and assigns a risk band (Section~\ref{sec:fed}).
\item \emph{Containment.} Whatever the model outputs, the contract bounds the result. An
undetected attack can at most cause L1 payments to approved payees, within the limits,
the discretionary amount and the day's allowance.
\item \emph{Shared alert.} Capsules narrow that remaining window while a threat is active
(Section~\ref{sec:capsule}).
\end{enumerate}

% ---------------------------------------------------------------------------
\section{Immunity Capsules}
\label{sec:capsule}

Retraining a detector takes time and needs examples. A wallet that has just been attacked
can warn the others at once.
\uniq{An immunity capsule is a small record of fixed fields that says what kind of attack
was seen, which features it had, which action it targeted, and how much autonomy to
withdraw and for how long.}

\subsection{Contents}
Table~\ref{tab:capsule} lists the fields.
\uniq{Every field is an enumeration, a set of flags, a number or a signature. A capsule
carries no message text, no prompt and no transaction, so it cannot leak private data and
cannot itself carry an injected instruction.}

\begin{table}[t]
\caption{Fields of an immunity capsule.}
\label{tab:capsule}
\tblfont
\begin{tabularx}{\columnwidth}{@{}P{22mm} L@{}}
\toprule
\textbf{Field} & \textbf{Values} \\
\midrule
Attack family & Direct, indirect, memory, encoded, recipient substitution \\
Channel & Owner chat, web page, message, memory, token metadata \\
Features & Any of: known payee identity, new destination, urgency, secrecy request, override attempt, encoded content, amount anomaly, memory-sourced \\
Affected action & Payment, schedule creation, sweep, payee or mandate change \\
Restriction & Delay (L2) or owner approval (L3) \\
Confidence & Single observer, or confirmed by $q$ wallets \\
Expiry & Time after which the capsule has no effect \\
Guard version & Detector version in force when observed \\
Attestations & Signatures of the observing wallets \\
\bottomrule
\end{tabularx}
\end{table}

\subsection{Matching reworded attacks}
\uniq{A capsule describes an attack by its features, not its wording.} Let $F_c$ be the
capsule's feature set and $F_I$ the flags of a request. The capsule \emph{fully matches}
when its channel is in the wallet's on-chain profile, its affected action equals the
request's action and $F_c\subseteq F_I$. It \emph{partially matches} when at least half
of $F_c$ is present. A full match raises the tier; a partial match is charged to the
allowance. \uniq{An attacker who rewrites the message but keeps the method produces the
same flags and is matched.} Three of the flags (new destination, amount anomaly,
memory-sourced) are computed by rule; the others are extracted by the reader and signed
by the guardian.

\subsection{Effect}
\uniq{The wallet contract reads unexpired capsules from the bulletin when it executes a
request, so a capsule takes effect on the very next transaction and does not rely on the
agent choosing to obey it.}
A wallet the threat does not apply to is unaffected and keeps its normal autonomy.

Some attacks are already stopped by the mandate: a payment to a new address needs the
owner whether or not a capsule exists. There the danger is the approval itself, because
this kind of fraud succeeds by persuading a person.
\uniq{When the owner is asked to approve, the interface shows any matching capsule, and
the approval is subject to the pressure hold of Section~\ref{sec:pressure}.}

\subsection{Limits on capsules}
A channel that can change other wallets' behaviour is itself a target, so capsules are
constrained:
\begin{itemize}
\item \uniq{\emph{Tighten only.} A capsule can raise a tier or charge the allowance. No
capsule can lower a tier, add a payee or raise a limit.}
\item \emph{Bounded by confidence.} A capsule from one observer can impose at most a
delay. Owner approval can be imposed only when $q$ distinct wallets have attested the
same family and features ($q=2$ in the prototype).
\item \emph{Expiry.} Every capsule expires; the maximum lifetime is fixed in the bulletin
contract. Autonomy returns without anyone acting.
\item \emph{Rate limit and identity.} Each registered wallet may publish a limited number
of capsules per day, and every capsule names its issuer.
\item \emph{Owner override.} An owner may dismiss a capsule for their own wallet.
\end{itemize}

\subsection{Alert first, vaccine later}
\uniq{The capsule acts within one transaction but is blunt and temporary. The retrained
detector arrives later but recognises variants beyond the listed features and stays.}
A capsule's lifetime is set to outlast the next training round, so protection is
continuous. Algorithm~\ref{alg:round} shows the full response.

% ---------------------------------------------------------------------------
\section{Federated Detector}
\label{sec:fed}

\subsection{Detector}
Each wallet $k$ keeps a private set $D_k=\{(x_i,y_i)\}$ of content snippets labelled
benign ($y=0$) or attack ($y=1$). Labels come from blocked requests the owner confirmed
as attacks and from requests the owner approved. A frozen text encoder $\phi$ maps a
snippet to a vector; a small classifier with weights $w$ outputs the score
\[
s(x)=\sigma\!\big(w^{\top}\phi(x)+b\big)\in[0,1].
\]
The risk band is low if $s<\tau_{\mathrm{low}}$, high if $s\ge\tau_{\mathrm{high}}$, and
medium otherwise. Only $(w,b)$ is trained and shared.

\subsection{Training round}
Each wallet minimises the local loss
\[
F_k(w)=\frac{1}{n_k}\sum_{(x,y)\in D_k}\ell\big(s(x),y\big),\qquad n_k=|D_k|,
\]
for $E$ epochs from the current global weights, and the aggregator computes the weighted
average~\cite{mcmahan}
\[
w^{t+1}=\sum_{k=1}^{K}\frac{n_k}{n}\,w_k^{t+1},\qquad n=\sum_{k=1}^{K}n_k .
\]
Updates are clipped to norm $C$ and masked pairwise so that the aggregator sees only the
sum~\cite{bonawitz}. No snippet, conversation or transaction leaves a wallet.

\begin{algorithm}[t]
\caption{Immunity response to an attack on wallet $a$}
\label{alg:round}
\footnotesize
\begin{algorithmic}[1]
\Require global weights $w^{t}$, wallets $1..K$, clip norm $C$
\State owner of $a$ confirms the attack; $\alpha_a \gets 0$
\State $a$ adds it to $D_a$ and publishes a capsule \Comment{alert}
\State wallets it applies to: $\alpha \gets \alpha/2$; matching requests restricted
\State a second observer may attest, raising it to confirmed
\For{each wallet $k$ in parallel} \Comment{vaccine}
  \State $w_k \gets$ \textsc{LocalTrain}$(w^{t}, D_k, E)$
  \State $\delta_k \gets$ \textsc{Clip}$(w_k - w^{t}, C)$
  \State send \textsc{Mask}$(n_k\,\delta_k)$ to the aggregator
\EndFor
\State $w^{t+1} \gets w^{t} + \frac{1}{n}\sum_k n_k\,\delta_k$ \Comment{masks cancel}
\State aggregator posts $H(w^{t+1})$ and $t{+}1$ to the bulletin
\For{each wallet $k$}
  \If{model hash matches bulletin \textbf{and} canary passes}
    \State adopt $w^{t+1}$
  \Else
    \State keep $w^{t}$ and raise an alert
  \EndIf
\EndFor
\State the capsule expires; allowances refill; the detector remains
\end{algorithmic}
\end{algorithm}

\subsection{Anchoring and adoption safeguards}
The bulletin stores the hash and version of each shared model.
\uniq{A wallet can therefore verify that it received the model the federation agreed on,
and an auditor can reconstruct which detector and which capsules governed any past
transaction.}
Before adopting a model, each wallet evaluates it on a private \emph{canary set}; a model
that lowers recall or raises the false-positive rate beyond a tolerance is rejected.
Together with clipping, this limits what a poisoning participant can do~\cite{krum}.

% ---------------------------------------------------------------------------
\section{Implementation}

Table~\ref{tab:stack} summarises the prototype. The federated layer is written against
Flower for speed of development; it uses only horizontal federated averaging with secure
aggregation, which FATE also provides, so it can be moved to FATE without redesign.
Contracts are plain Solidity with no chain-specific features, which keeps a later port to
a permissioned EVM-compatible chain such as FISCO BCOS~\cite{fisco} straightforward.

\begin{table}[t]
\caption{Prototype implementation.}
\label{tab:stack}
\tblfont
\begin{tabularx}{\columnwidth}{@{}P{23mm} L@{}}
\toprule
\textbf{Component} & \textbf{Implementation} \\
\midrule
Chain & Public EVM test network; test stablecoin with no monetary value \\
Policy wallet & Solidity smart account: signed mandate, session key, allowlist, limits, obligation registry, caution allowance, time-lock queue for payments and for owner changes, pressure hold, co-signature, freeze \\
Savings vault & Minimal deposit and withdraw contract, owner-only withdrawal \\
Threat bulletin & Registry of capsules, attestations and model hashes; enforces capsule limits \\
Planner, reader & Two separately prompted language-model roles; reader has no tools \\
Validator & Deterministic code; schema, provenance and state checks \\
Guardian & Encoder plus linear classifier; separate signing key \\
Federation & Three clients, federated averaging, pairwise masking \\
VaxLab & Scripted attackers plus one rewording agent; never holds a wallet key \\
Interface & Web chat with mandate view, financial state panel, allowance meter, alerts and considered confirmation at approval, explorer links \\
\bottomrule
\end{tabularx}
\end{table}

% ---------------------------------------------------------------------------
\section{Evaluation Protocol}
\label{sec:eval}

All measurements are produced by scripts in the repository so that reviewers can repeat
them. Table~\ref{tab:metrics} lists the metrics and the targets the design aims for.
Targets are design goals, not results. For every task, tier path, mandate entry, capsule
and model anchor, the submission lists the transaction hash and contract address.

\begin{table}[t]
\caption{Evaluation metrics and design targets.}
\label{tab:metrics}
\tblfont
\begin{tabularx}{\columnwidth}{@{}L P{30mm} P{16mm}@{}}
\toprule
\textbf{Metric} & \textbf{Definition} & \textbf{Target} \\
\midrule
Safe autonomy rate & Legitimate tasks done without the owner, zero violations & $\ge 80\%$ \\
Hard-policy violations & Mandate breaches of any kind & 0 \\
Task success rate & Correct state change on chain, 40 instructions & $\ge 90\%$ \\
Ambiguous requests handled & Clarified or refused, 10 instructions & $\ge 90\%$ \\
Affordability decisions & Agreement with hand-computed tier, 20 scenarios & 100\% \\
Unauthorised value moved & Across all VaxLab attacks & 0 beyond L1 limits \\
Probing sequence & Odd payments executed before the owner is asked & Within the allowance \\
Panicked-owner loss & Value leaving before a hold or delay ends when every prompt is approved at once & 0 \\
Capsule response & Transactions before restriction is active & 0 \\
Reworded variants matched & After one capsule & $\ge 90\%$ \\
False-capsule harm & Rules loosened by hostile capsules & 0 \\
Cross-family recall gain & Recall on unseen families, (c) minus (b) & Positive for every wallet \\
Escalation and false-positive rates & Benign requests sent to the owner; escalated by the detector & Reported; $\le 5\%$ \\
Friction and cost & Wait imposed on legitimate new payees and urgent requests; time to full allowance; seconds and gas per task & Reported \\
\bottomrule
\end{tabularx}
\end{table}

\subsection{Headline measure}
A wallet that asks the owner about everything is safe and useless; one that asks about
nothing is useful and unsafe.
\uniq{We therefore report the \emph{safe autonomy rate}: the share of legitimate tasks
completed without the owner, counted only if no hard policy was violated anywhere in the
run.}
\begin{gather*}
\textit{SAR}=\frac{\text{legitimate tasks completed without the owner}}
{\text{all legitimate tasks}},\\
\text{reported only if hard-policy violations}=0 .
\end{gather*}

\begin{table*}[t]
\caption{Known-risk list.}
\label{tab:risks}
\tblfont
\begin{tabularx}{\textwidth}{@{}P{0.04\textwidth} P{0.20\textwidth} L L@{}}
\toprule
\textbf{ID} & \textbf{Risk} & \textbf{Consequence} & \textbf{Mitigation and residual exposure} \\
\midrule
R1 & Undetected injection & Planner proposes an attacker-chosen action &
Contract bounds the outcome: L1 payments to approved payees within the limits, the discretionary amount and the day's allowance \\
R2 & Guardian key compromise & Attacker can authorise L1 execution and understate flags &
Key held in a separate process; limits, affordability and rule-computed flags still apply; owner can rotate the key and freeze the wallet \\
R3 & Reader misses a flag & A capsule fails to match, or the allowance is undercharged &
Three flags are computed by rule and cannot be suppressed; a missed flag returns the request to the ordinary tier rule, not to no rule \\
R4 & Owner approves a fraud under pressure & Funds leave with valid authorisation &
New payees and higher limits wait for the loosening delay; pressured approvals are held; cancelling is instant. Residual: an owner still deceived when the delay ends \\
R5 & Delays obstruct a genuine emergency & A legitimate urgent payment is late &
The owner sets the delays while calm; a trusted co-signer can release them; the cost is measured and reported \\
R6 & Incomplete obligation registry & A bill nobody registered is not protected &
Schedules register automatically; the agent proposes obligations it finds in invoices; the reserve provides a margin \\
R7 & Arbitrary allowance values & Too strict wastes the owner's time; too loose admits probing &
Owner-set integers with stated defaults; escalation rate and probing result reported \\
R8 & False or excessive capsules & Honest wallets are slowed (denial of service) &
Tighten-only; single-observer capsules limited to a delay; rate limit; expiry; owner may dismiss \\
R9 & Small federation & With three clients, two colluding parties can infer the third party's update &
Stated limitation of the prototype; privacy improves with more participants \\
R10 & Model poisoning by a participant & Shared detector degrades or learns a blind spot &
Update clipping; canary test before adoption; on-chain version history allows rollback \\
R11 & False positives & Legitimate payments are delayed or need approval &
Thresholds tuned on benign data; L2 delays, not blocks; rates reported \\
R12 & Compromised approved payee & Funds sent to a legitimate but hijacked payee &
Limits, usual-amount flag, discretionary amount and allowance limit exposure \\
R13 & Hosted language model; testnet conditions & Outage breaks new requests; fees and finality are not representative &
Scheduled payments run without the model; contracts avoid chain-specific features \\
\bottomrule
\end{tabularx}
\end{table*}

\subsection{Task and affordability suites}
The task suite contains 40 instructions (ten per task, varied in wording, amounts and
payees) and 10 that are ambiguous or impossible. A task succeeds if the expected on-chain
state change is confirmed by the receipt; an ambiguous instruction succeeds if the agent
asks or refuses with the correct reason and sends nothing. The affordability suite pairs
20 financial states with requests, covering comfortable, tight and breach cases, payment
of a registered obligation, and income that has not arrived. Outcomes are compared with
hand-computed answers; the rule is deterministic, so any disagreement is a defect.

\subsection{VaxLab attack suite}
VaxLab is a separate test harness that plays the attacker and never holds a wallet key.
Scripted attackers cover five families: direct instructions, indirect instructions in
retrieved content, false memory entries, encoded instructions (for example Morse code and
base64), and recipient substitution. Three further scenarios test the new mechanisms: a
\emph{probing sequence} of slightly unusual small payments, a hostile wallet that
publishes false capsules and poisoned updates, and a \emph{panicked owner}, a script
that approves every prompt at once. \uniq{A rewording agent paraphrases every
attack, so each defence is tested on versions it has not seen.} An attack \emph{succeeds}
only if an unauthorised transaction is executed on chain.

\subsection{Shared-response experiments}
\emph{Capsules.} Wallet A is attacked and publishes a capsule. We measure how many
transactions pass before matching requests on other wallets are restricted, the share of
reworded variants matched, whether a wallet outside the capsule's scope stays
unrestricted, and how long autonomy takes to return. \emph{Federation.} Each wallet is
seeded with a different attack family; the others are held out. We compare
(a)~contract checks only, (b)~plus a locally trained detector, (c)~plus the federated
detector, measuring each wallet's recall on families it has never seen locally.


\subsection{Demonstration scenario}
The live demonstration follows one story across three wallets:
\begin{enumerate}
\item The owner of wallet A states a mission; the mandate is shown, signed and stored on
chain.
\item The agent pays rent and a hosting invoice on its own (L1).
\item Expected income is late. A 300 hosting upgrade is requested; the state engine
finds it tight, and the agent holds it and explains the 120 shortfall.
\item An invoice arrives that appears to come from the hosting provider, with a new
payment address, urgent wording and a request to ignore earlier details. The new address
needs the owner; the owner rejects it and confirms fraud.
\item Wallet A publishes a capsule; its allowance drops to zero.
\item Wallet B receives a reworded version. The flags match and the alert is shown.
B's owner, in a hurry, approves anyway. Nothing moves: the new address is subject to
the loosening delay, and the owner later cancels it. Wallet C does not accept that
channel and is unaffected.
\item Wallet B then receives an inflated invoice from an approved payee to its usual
address. Alone it would pass; with the partial match and the halved allowance it waits.
\item A federated round runs, the new detector's hash is anchored, the capsule expires
and the allowances refill.
\end{enumerate}

% ---------------------------------------------------------------------------
\section{Compliance}

The prototype is designed to be acceptable to a regulated institution. All assets are
test tokens on a test network; no token is issued and nothing is traded for profit. Each
wallet has an identified owner who signs the mandate, holds the override key and is
responsible for the agent's actions. Every intent, decision, guardian signature,
allowance charge, hold, capsule and model version is logged, and the on-chain part is
public. Conversations and transaction details stay on the wallet's own node; only
fixed-field capsules, masked model updates and hashes leave it. Third-party code and
models are used under their open-source licences.

% ---------------------------------------------------------------------------
\section{Known Risks and Limitations}

Table~\ref{tab:risks} lists the risks we are aware of. It is the basis of the security
self-assessment report required by the rules.


% ---------------------------------------------------------------------------
\section{Future Work}
\label{sec:future}

Left out of the submitted prototype: a full adversarial swarm in VaxLab with an
evolutionary generator of attacks; reputation for capsule issuers and resistance to many
fake wallets; a mutual-aid pool that compensates verified losses~\cite{insured};
deployment on a permissioned chain; income forecasting and differential privacy for
updates.

% ---------------------------------------------------------------------------
\section{Conclusion}

\uniq{AgentVax is a self-immunising wallet agent that adjusts how much financial autonomy
it exercises while its authority stays under limits only the owner can change.}
Usefulness comes from a mission the owner states once and tasks that complete on chain.
Bounds come from a mandate enforced by the contract and a guardian signature the model
cannot forge.
\uniq{Prudence comes from asking whether the owner can afford a payment and how much
uncertainty has already been accepted today. Learning comes in two speeds: a capsule
that makes exposed wallets cautious at once, and a jointly trained detector that keeps
them protected.}
\uniq{The owner is protected as well: authority can be reduced in a moment and extended
only with time.}
The evaluation protocol makes each of these claims measurable.

% ---------------------------------------------------------------------------
\balance
\begin{thebibliography}{99}
\footnotesize
\setlength{\itemsep}{0pt}\setlength{\parskip}{0pt}\setlength{\parsep}{0pt}

\bibitem{rules} Organizing Committee, ``International Track: Topics \& Rules,'' 2026
Shenzhen International FinTech Competition (FinTechathon), 2026.
\url{http://fintechathon.g-ican.com}

\bibitem{greshake} K.~Greshake, S.~Abdelnabi, S.~Mishra, C.~Endres, T.~Holz and
M.~Fritz, ``Not what you've signed up for: Compromising real-world LLM-integrated
applications with indirect prompt injection,'' in \emph{Proc.\ ACM Workshop on Artificial
Intelligence and Security (AISec)}, 2023.

\bibitem{owasp} OWASP Foundation, ``OWASP Top 10 for Large Language Model
Applications,'' 2025.

\bibitem{oecdgrok} OECD.AI Policy Observatory, ``AI prompt injection exploit drains
Grok-linked crypto wallet,'' AI Incidents and Hazards Monitor, May 2026.
\url{https://oecd.ai/en/incidents/2026-05-04-4a73}

\bibitem{giskard} Giskard, ``How Grok got prompt-injected: An X user drained
\$150,000 from an AI wallet,'' giskard.ai, 2026.

\bibitem{patlan} A.~S.~Patlan, P.~Sheng, S.~A.~Hebbar, P.~Mittal and P.~Viswanath,
``Real AI agents with fake memories: Fatal context manipulation attacks on Web3
agents,'' arXiv:2503.16248, 2025.

\bibitem{imf} T.~Davidovic and H.~Tourpe, ``How agentic AI will reshape payments,''
IMF Note 2026/004, International Monetary Fund, 2026.

\bibitem{freysa} The Block, ``Human player outwits Freysa AI agent in \$47,000 crypto
challenge,'' Nov.\ 2024.

\bibitem{aixbt} AI Incident Database, ``Incident 1003'' (AIXBT autonomous agent), 2025.
\url{https://incidentdatabase.ai/cite/1003}

\bibitem{lobstar} The Block, ``AI agent created by OpenAI dev `accidentally' sends
entire memecoin holdings to reply guy,'' Feb.\ 2026.

\bibitem{erc4337} V.~Buterin \emph{et al.}, ``ERC-4337: Account abstraction using alt
mempool,'' Ethereum Improvement Proposals, 2021.

\bibitem{erc8004} ``ERC-8004: Trustless agents,'' Ethereum Improvement Proposals, 2025.

\bibitem{willison} S.~Willison, ``The dual LLM pattern for building AI assistants that
can resist prompt injection,'' simonwillison.net, Apr.\ 2023.

\bibitem{camel} E.~Debenedetti \emph{et al.}, ``Defeating prompt injections by design,''
arXiv:2503.18813, 2025.

\bibitem{stix} OASIS, ``STIX Version 2.1,'' OASIS Standard, 2021.

\bibitem{mcmahan} B.~McMahan, E.~Moore, D.~Ramage, S.~Hampson and B.~A.~y~Arcas,
``Communication-efficient learning of deep networks from decentralized data,'' in
\emph{Proc.\ AISTATS}, 2017.

\bibitem{yang} Q.~Yang, Y.~Liu, T.~Chen and Y.~Tong, ``Federated machine learning:
Concept and applications,'' \emph{ACM Trans.\ Intelligent Systems and Technology},
vol.~10, no.~2, 2019.

\bibitem{bonawitz} K.~Bonawitz \emph{et al.}, ``Practical secure aggregation for
privacy-preserving machine learning,'' in \emph{Proc.\ ACM CCS}, 2017.

\bibitem{fate} Y.~Liu, T.~Fan, T.~Chen, Q.~Xu and Q.~Yang, ``FATE: An industrial grade
platform for collaborative learning with data protection,'' \emph{J.\ Machine Learning
Research}, vol.~22, no.~226, 2021.

\bibitem{flower} D.~J.~Beutel \emph{et al.}, ``Flower: A friendly federated learning
research framework,'' arXiv:2007.14390, 2020.

\bibitem{krum} P.~Blanchard, E.~M.~El~Mhamdi, R.~Guerraoui and J.~Stainer, ``Machine
learning with adversaries: Byzantine tolerant gradient descent,'' in \emph{Proc.\
NeurIPS}, 2017.

\bibitem{insured} Hu and Chen, ``Insured agents: A decentralized trust insurance
mechanism for agentic economy,'' arXiv:2512.08737, 2025.

\bibitem{fisco} FISCO BCOS Open Source Community, ``FISCO BCOS: An enterprise-level
financial blockchain platform,'' open-source software, github.com/FISCO-BCOS.

\bibitem{masabs} Monetary Authority of Singapore and Association of Banks in Singapore,
``MAS and ABS announce measures to bolster the security of digital banking,'' 19 Jan.\
2022.

\bibitem{ukdelay} HM Treasury, ``Payment Services (Amendment) Regulations 2024,'' draft
legislation published Mar.\ 2024.

\bibitem{commitment} G.~Bryan, D.~Karlan and S.~Nelson, ``Commitment devices,''
\emph{Annual Review of Economics}, vol.~2, 2010.

\end{thebibliography}

\end{document}
